\chapter{Beberapa Variabel dan Turunan Parsial} \label{pd:chapter}

%%%%%%%%%%%%%%%%%%%%%%%%%%%%%%%%%%%%%%%%%%%%%%%%%%%%%%%%%%%%%%%%%%%%%%%%%%%%%%

\section{Ruang vektor, pemetaan linear, dan kekonveksan}
\label{sec:vectorspaces}

%mbxINTROSUBSECTION

\sectionnotes{3 kuliah}

\subsection{Ruang vektor}

Ruang Euklides $\R^n$ telah muncul dalam bab ruang metrik.
Dalam bab ini, kita memperluas kalkulus diferensial yang kita bangun
untuk satu variabel ke beberapa variabel.  Gagasan utama dalam
kalkulus diferensial adalah menghampiri fungsi diferensiabel dengan fungsi
linear
(menghampiri grafik dengan objek datar seperti garis atau bidang).
Untuk beberapa variabel, kita perlu memperkenalkan sedikit aljabar
linear sebelum dapat melanjutkan pembahasan.
Kita mulai dengan ruang vektor dan pemetaan linear antarruang vektor.

Notasi $\vec{v}$ atau cetak tebal
$\mathbf{v}$ lazim digunakan untuk elemen $\R^n$,
terutama dalam ilmu terapan.
Kita cukup menggunakan $v$ biasa, sebagaimana lazim dalam matematika.
Artinya, $v \in \R^n$ adalah sebuah \emph{\myindex{vektor}}, yaitu
$v = (v_1,v_2,\ldots,v_n)$\glsadd{not:vector} merupakan tupel-$n$ dari
bilangan real.\footnote{Subskrip digunakan untuk berbagai keperluan,
sehingga terkadang kita memiliki beberapa vektor yang juga
dibedakan dengan subskrip, seperti barisan berhingga atau tak berhingga dari
vektor-vektor $y_1,y_2,\ldots$.}
Adakalanya vektor berguna untuk ditulis dan diperlakukan sebagai
\emph{\myindex{vektor kolom}}, yaitu matriks berukuran $n$-kali-$1$:
\glsadd{not:vectorcol}
\begin{equation*}
v =
(v_1,v_2,\ldots,v_n) =
{ \left[
\begin{smallmatrix}
v_1 \\ v_2 \\ \vdots \\ v_n
\end{smallmatrix}
\right]
}
\end{equation*}
Kita menggunakan bentuk ini ketika diperlukan.
Bilangan real kita sebut
\emph{\myindex{skalar}} untuk membedakannya dari vektor.

Kita sering memandang vektor sebagai arah dan besar, lalu menggambar
vektor sebagai anak panah.  Vektor $(v_1,v_2,\ldots,v_n)$ direpresentasikan
oleh anak panah dari titik asal ke titik $(v_1,v_2,\ldots,v_n)$.
Ketika kita memandang vektor sebagai anak panah, pangkalnya
tidak harus berada di titik asal; sebuah vektor hanyalah arah dan
besar, tanpa memuat informasi tentang letak titik pangkalnya.
Terdapat struktur aljabar alami ketika vektor dipandang sebagai anak panah.
Kita dapat menjumlahkan vektor sebagai anak panah dengan menyusuri satu vektor lalu vektor lainnya.
Kita juga dapat mengalikan vektor dengan skalar dengan mengubah
skalanya.  Lihat \figureref{fig:vecasarrow}.
\begin{myfigureht}
\myincludepdft{vecasarrow}{%
Di sebelah kiri terdapat diagram vektor (v subskrip 1, v subskrip 2)
pada bidang x subskrip 1 x subskrip 2 sebagai anak panah dari titik asal
ke titik (v subskrip 1, v subskrip 2).
Di sebelah kanan terdapat diagram penjumlahan vektor dan perkalian
skalar.
Pertama, terdapat gambar anak panah v yang disusul oleh
vektor w yang diletakkan pada ujungnya, kemudian vektor
v tambah w yang bermula di pangkal v dan berakhir di ujung w.
Selanjutnya terdapat gambar anak panah v dan anak panah 2 v, dengan
anak panah 2 v searah dengan v tetapi panjangnya dua
kali panjang v.}
\caption{Vektor sebagai anak panah dalam $\R^2$, serta arti penjumlahan dan
perkalian skalar.\label{fig:vecasarrow}}
\end{myfigureht}

Sebuah vektor juga dapat merepresentasikan sebuah titik dalam $\R^n$.
Biasanya, kita memandang $v \in \R^n$
sebagai titik jika $\R^n$ dipandang sebagai ruang metrik, dan memandangnya
sebagai anak panah jika kita memperhatikan \emph{struktur ruang vektor} pada
$\R^n$ (penjumlahan dan perkalian skalar).
Mari kita definisikan pengertian abstrak ruang vektor, sebab terdapat
banyak ruang vektor lain selain $\R^n$.

\begin{defn}
Misalkan $X$ adalah suatu himpunan yang dilengkapi dengan
operasi penjumlahan, $+ \colon X \times X \to X$,
dan operasi perkalian skalar, $\cdot \colon \R \times X \to X$,
(biasanya kita menulis $ax$ alih-alih $a \cdot x$).
Himpunan $X$ disebut \emph{\myindex{ruang vektor}}
(atau \emph{\myindex{ruang vektor real}})
jika syarat-syarat berikut dipenuhi:
\begin{enumerate}[(i)]
\item
%mbxSTARTIGNORE
\makebox[2.15in][l]{(Penjumlahan bersifat asosiatif)}
%mbxENDIGNORE
%mbxlatex (Penjumlahan bersifat asosiatif) \quad
Jika $u, v, w \in X$, maka $u+(v+w) = (u+v)+w$.
\item
%mbxSTARTIGNORE
\makebox[2.15in][l]{(Penjumlahan bersifat komutatif)}
%mbxENDIGNORE
%mbxlatex (Penjumlahan bersifat komutatif) \quad
Jika $u, v \in X$, maka $u+v = v+u$.
\item \label{vecspacedefn:addidentity}
%mbxSTARTIGNORE
\makebox[2.15in][l]{(Identitas aditif)}
%mbxENDIGNORE
%mbxlatex (Identitas aditif) \quad
Terdapat $0 \in X$ sehingga $v+0=v$ untuk setiap $v \in X$.
\item \label{vecspacedefn:addinverse}
%mbxSTARTIGNORE
\makebox[2.15in][l]{(Invers aditif)}
%mbxENDIGNORE
%mbxlatex (Invers aditif) \quad
Untuk setiap $v \in X$, terdapat $-v \in X$
sehingga $v+(-v)=0$.
\item
%mbxSTARTIGNORE
\makebox[2.15in][l]{(Hukum distributif)}
%mbxENDIGNORE
%mbxlatex (Hukum distributif) \quad
Jika $a \in \R$ dan $u,v \in X$, maka
$a(u+v) = au+av$.
\item
%mbxSTARTIGNORE
\makebox[2.15in][l]{(Hukum distributif)}
%mbxENDIGNORE
%mbxlatex (Hukum distributif) \quad
Jika $a,b \in \R$ dan $v \in X$, maka
$(a+b)v = av+bv$.
\item
%mbxSTARTIGNORE
\makebox[2.15in][l]{(Perkalian bersifat asosiatif)}
%mbxENDIGNORE
%mbxlatex (Perkalian bersifat asosiatif) \quad
Jika $a,b \in \R$ dan $v \in X$, maka
$(ab)v = a(bv)$.
\item
%mbxSTARTIGNORE
\makebox[2.15in][l]{(Identitas perkalian)}
%mbxENDIGNORE
%mbxlatex (Identitas perkalian) \quad
$1v = v$ untuk setiap $v \in X$.
\end{enumerate}
Elemen-elemen ruang vektor biasanya disebut \emph{vektor},
meskipun elemen-elemen tersebut
bukan elemen $\R^n$ (vektor dalam pengertian \myquote{tradisional}).
%
Jika $Y \subset X$ adalah himpunan bagian yang juga merupakan ruang vektor dengan
operasi yang sama, maka $Y$ disebut \emph{\myindex{subruang}} atau
\emph{\myindex{subruang vektor}} dari $X$.
\end{defn}

Perkalian dengan skalar berlaku sebagaimana yang kita harapkan.
Sebagai contoh, $2v = (1+1)v = 1v + 1v = v+v$, demikian pula $3v = v+v+v$, dan
seterusnya.
Satu fakta khusus yang sering kita gunakan ialah $0 v = 0$, dengan
nol di sebelah kiri adalah $0 \in \R$ dan nol di sebelah kanan adalah $0 \in X$.
Untuk melihatnya, mulailah dengan
$0v = (0+0)v = 0v + 0v$, lalu tambahkan $-(0v)$ pada kedua ruas
untuk memperoleh $0 = 0v$.  Demikian pula, $-v = (-1)v$, yang diperoleh dari
$(-1)v+v = (-1)v + 1v = (-1+1)v = 0v = 0$.
Mulai sekarang, fakta-fakta aljabar semacam ini, yang langsung mengikuti
definisi, akan digunakan tanpa pembuktian.

\begin{example}
Himpunan $\R^n$ merupakan ruang vektor;
vektor nolnya adalah $(0,\ldots,0)$, sedangkan penjumlahan
dan perkalian dengan skalar dilakukan komponen demi komponen:
Jika $a \in \R$, $v = (v_1,v_2,\ldots,v_n) \in \R^n$, dan $w =
(w_1,w_2,\ldots,w_n) \in \R^n$, maka
\begin{align*}
& v+w \coloneqq
(v_1,v_2,\ldots,v_n) +
(w_1,w_2,\ldots,w_n) 
=
(v_1+w_1,v_2+w_2,\ldots,v_n+w_n) , \\
& a v \coloneqq
a (v_1,v_2,\ldots,v_n) =
(a v_1, a v_2,\ldots, a v_n) .
\end{align*}
\end{example}

Kita terutama akan membahas ruang vektor \myquote{berdimensi hingga} yang dapat
dipandang sebagai himpunan bagian dari $\R^n$, tetapi ruang vektor lain juga berguna dalam
analisis.
Bahkan ruang vektor yang lebih sederhana seperti ini sebaiknya dipandang
secara abstrak, bukan sebagai $\R^n$.

\begin{example}
Sebuah contoh trivial ruang vektor adalah
$X \coloneqq \{ 0 \}$.  Operasinya didefinisikan dengan cara yang jelas:
$0 + 0 \coloneqq 0$ dan $a0 \coloneqq 0$.  Vektor nol harus selalu ada,
sehingga semua ruang vektor merupakan himpunan tak kosong, dan $X$ ini adalah ruang
vektor terkecil yang mungkin.
\end{example}

\begin{example}
Ruang $C([0,1],\R)$ yang terdiri atas fungsi-fungsi kontinu pada interval $[0,1]$
merupakan ruang vektor.  Untuk dua fungsi $f$ dan $g$ dalam $C([0,1],\R)$ serta
$a \in \R$, kita mendefinisikan $f+g$ dan $af$ dengan cara yang jelas:
\begin{equation*}
(f+g)(x) \coloneqq f(x) + g(x), \qquad (af) (x) \coloneqq a\bigl(f(x)\bigr) .
\end{equation*}
Elemen 0 adalah fungsi yang identik nol.  Sebagai latihan, periksalah
bahwa semua syarat ruang vektor dipenuhi.
Ruang $C^1([0,1],\R)$ yang terdiri atas fungsi-fungsi yang diferensiabel dan memiliki turunan kontinu merupakan
subruang dari $C([0,1],\R)$.
\end{example}

\begin{example}
Ruang polinom $c_0 + c_1 t + c_2 t^2 + \cdots + c_m t^m$
(dengan derajat~$m$ sebarang) merupakan ruang vektor,
yang dinotasikan dengan $\R[t]$\glsadd{not:realpoly} (koefisiennya real dan
variabelnya adalah $t$).  Operasinya didefinisikan dengan cara yang sama seperti untuk
fungsi di atas.
Misalkan terdapat
dua polinom, yang satu berderajat $m$ dan yang lain berderajat $n$.  Agar sederhana, andaikan $n
\geq m$.  Maka
\begin{multline*}
(c_0 + c_1 t + c_2 t^2 + \cdots + c_m t^m)
+
(d_0 + d_1 t + d_2 t^2 + \cdots + d_n t^n)
= \\
(c_0+d_0) + (c_1+d_1) t + (c_2 + d_2) t^2 + \cdots + (c_m+d_m) t^m
+ d_{m+1} t^{m+1} + \cdots + d_n t^n
\end{multline*}
dan
\begin{equation*}
a(c_0 + c_1 t + c_2 t^2 + \cdots + c_m t^m)
=
(ac_0) + (ac_1) t + (ac_2) t^2 + \cdots + (ac_m) t^m  .
\end{equation*}
Walaupun sekilas tampak demikian, $\R[t]$ tidak ekuivalen dengan $\R^n$ untuk $n$ mana pun.
Secara khusus, ruang ini tidak \myquote{berdimensi hingga.}  Sebentar lagi kita akan merumuskan
pengertian ini secara tepat.  Kita dapat membentuk suatu
subruang vektor berdimensi hingga dengan membatasi derajatnya.  Sebagai contoh,
jika $\sP_n$ adalah himpunan polinom berderajat $n$ atau kurang,
maka $\sP_n$ merupakan ruang vektor berdimensi hingga, dan dapat kita
identifikasikan dengan $\R^{n+1}$.

Variabel $t$ di atas sebenarnya hanyalah penanda tempat formal.
Dengan menetapkan $t$ sama dengan suatu bilangan real, kita memperoleh sebuah fungsi.
Jadi, ruang $\R[t]$ dapat dipandang sebagai subruang dari $C(\R,\R)$.
Jika kita membatasi rentang nilai $t$ pada $[0,1]$, $\R[t]$ dapat diidentifikasikan dengan
subruang dari $C([0,1],\R)$.
\end{example}

\begin{prop}
Agar $S \subset X$ menjadi subruang vektor
dari ruang vektor $X$, kita hanya perlu memeriksa:
%If $X$ is a vector space, for $S \subset X$
%to be a vector subspace, we only need to check
\begin{enumerate}[1)]
\item
$0 \in S$.
\item
$S$ tertutup terhadap penjumlahan: Jika $x,y \in S$, maka $x+y \in S$.
\item
$S$ tertutup terhadap perkalian skalar:
Jika $x \in S$ dan $a \in \R$, maka $ax \in S$.
\end{enumerate}
\end{prop}

Butir 2) dan 3)
memastikan bahwa penjumlahan dan perkalian skalar memang terdefinisi pada~$S$.
Butir 1) diperlukan
untuk memenuhi butir \ref{vecspacedefn:addidentity} dalam definisi ruang
vektor.  Keberadaan
invers aditif $-v$, yaitu butir \ref{vecspacedefn:addinverse},
mengikuti karena $-v = (-1)v$ dan butir 3) menyatakan bahwa
$-v \in S$ jika $v \in S$.  Semua sifat lainnya berupa persamaan
yang sudah berlaku dalam $X$ sehingga juga harus berlaku dalam suatu himpunan bagian.

\medskip

Kita dapat menggunakan lapangan selain $\R$ dalam definisi
(sebagai contoh, bilangan kompleks $\C$ juga lazim digunakan),
tetapi kita akan tetap menggunakan
bilangan real\footnote{Jika Anda menginginkan ruang vektor yang sangat tidak biasa atas lapangan lain,
$\R$ sendiri merupakan ruang vektor atas lapangan $\Q$.}.

\subsection{Kombinasi linear dan dimensi}

\begin{defn}
Misalkan $X$ adalah suatu ruang vektor,
$x_1, x_2, \ldots, x_m \in X$ adalah vektor, dan
$a_1, a_2, \ldots, a_m \in \R$ adalah skalar.  Maka
\begin{equation*}
a_1 x_1 + 
a_2 x_2 +  \cdots
+ a_m x_m
\end{equation*}
disebut \emph{\myindex{kombinasi linear}} dari vektor-vektor $x_1, x_2,
\ldots, x_m$.

Untuk suatu subhimpunan $Y \subset X$, nyatakan $\spn(Y)$,\glsadd{not:span}
atau \emph{\myindex{rentang}} dari $Y$,
sebagai himpunan semua kombinasi linear dari semua subhimpunan hingga $Y$.
Kita katakan $Y$ \emph{merentang} $\spn(Y)$.
Secara konvensi, tetapkan $\spn(\emptyset) \coloneqq \{ 0 \}$.
\end{defn}

\begin{example}
Misalkan $Y \coloneqq \bigl\{ (1,1) \bigr\} \subset \R^2$.  Maka
\begin{equation*}
\spn(Y)
=
\bigl\{ (x,x) \in \R^2 : x \in \R \bigr\} .
\end{equation*}
Dengan kata lain, $\spn(Y)$ adalah garis yang melalui titik asal dan titik $(1,1)$.
\end{example}

\begin{example} \label{example:vecspr2span}
Misalkan $Y \coloneqq \bigl\{ (1,1), (0,1) \bigr\} \subset \R^2$.  Maka
\begin{equation*}
\spn(Y)
=
\R^2 ,
\end{equation*}
karena setiap titik $(x,y) \in \R^2$ dapat ditulis sebagai
suatu kombinasi linear
\begin{equation*}
(x,y) = x (1,1) + (y-x) (0,1) .
\end{equation*}
\end{example}

\begin{example}
Misalkan $Y \coloneqq \{ 1, t, t^2, t^3, \ldots \} \subset \R[t]$ dan
$E \coloneqq \{ 1, t^2, t^4, t^6, \ldots \} \subset \R[t]$.
Rentang $Y$ adalah seluruh himpunan polinom,
\begin{equation*}
\spn(Y) = \R[t] .
\end{equation*}
Rentang $E$ adalah himpunan polinom yang hanya memuat pangkat genap dari $t$.
\end{example}

Misalkan kita memiliki dua kombinasi linear vektor dari $Y$.  Kombinasi linear
pertama menggunakan vektor $\{ x_1,x_2,\ldots,x_m \}$, sedangkan yang lain menggunakan
$\{ \widetilde{x}_1,\widetilde{x}_2,\ldots,\widetilde{x}_n \}$.  Jumlah keduanya dapat ditulis
menggunakan vektor-vektor dari gabungan
$\{ x_1,x_2,\ldots,x_m \} \cup
\{ \widetilde{x}_1,\widetilde{x}_2,\ldots,\widetilde{x}_n \}$:
\begin{multline*}
(a_1 x_1 + 
a_2 x_2 +  \cdots
+ a_m x_m)
+
(b_1 \widetilde{x}_1 + 
b_2 \widetilde{x}_2 +  \cdots
+ b_n \widetilde{x}_n)
\\
=
a_1 x_1 + 
a_2 x_2 +  \cdots
+ a_m x_m
+
b_1 \widetilde{x}_1 + 
b_2 \widetilde{x}_2 +  \cdots
+ b_n \widetilde{x}_n .
\end{multline*}
Jadi, jumlah tersebut juga merupakan kombinasi linear vektor-vektor dari~$Y$.
Demikian pula,
kelipatan skalar dari suatu kombinasi linear vektor-vektor dari~$Y$
merupakan kombinasi linear vektor-vektor dari~$Y$:
\begin{equation*}
b (a_1 x_1 + 
a_2 x_2 +  \cdots
+ a_m x_m)
=
b a_1  x_1 + 
b a_2 x_2 +  \cdots
+ b a_m x_m .
\end{equation*}
Terakhir, $0 \in \spn(Y)$; jika $Y$ tak kosong, $0 = 0 v$ untuk suatu $v \in Y$.
Kita telah membuktikan proposisi berikut.

\begin{prop}
Misalkan $X$ adalah ruang vektor dan $Y \subset X$ suatu subhimpunan.
Maka himpunan $\spn(Y)$ merupakan subruang vektor dari $X$.
\end{prop}

Setiap kombinasi linear unsur-unsur suatu subruang merupakan unsur dari
subruang itu.  Jadi, $\spn(Y)$ adalah subruang terkecil yang memuat $Y$.
Secara khusus,
jika $Y$ sudah merupakan subruang vektor, maka $\spn(Y) = Y$.

\begin{defn}
Suatu himpunan vektor $\{ x_1, x_2, \ldots, x_m \} \subset X$ disebut 
\emph{\myindex{bebas linear}}\footnote{%
Untuk himpunan tak hingga $Y \subset X$, kita katakan $Y$ bebas
linear jika setiap subhimpunan hingga dari $Y$ bebas linear
dalam pengertian yang diberikan di sini.
Namun, keadaan ini hanya muncul pada dimensi tak berhingga.}
jika persamaan
\begin{equation} \label{eq:lincomb}
a_1 x_1 + a_2 x_2 + \cdots + a_m x_m = 0
\end{equation}
hanya memiliki solusi trivial $a_1 = a_2 = \cdots = a_m = 0$.
Secara konvensi, $\emptyset$ bebas linear.
Suatu himpunan yang tidak bebas linear disebut
\emph{\myindex{bergantung linear}}.
%
Suatu himpunan vektor bebas linear $B \subset X$ sedemikian sehingga
$\spn(B) = X$ 
disebut \emph{\myindex{basis}} dari $X$.
Pada umumnya, kita memandang basis bukan hanya sebagai suatu himpunan, melainkan sebagai
suatu tupel-$m$ terurut:
$x_1,x_2,\ldots,x_m$.

Misalkan $d$ adalah bilangan bulat terbesar sedemikian sehingga $X$ memuat suatu himpunan
yang terdiri atas $d$ vektor bebas linear.  Kita sebut $d$ sebagai
\emph{\myindex{dimensi}} dari $X$,
dan kita tulis $\dim \, X \coloneqq d$.
Jika $X$ memuat suatu himpunan yang terdiri atas $d$ vektor bebas linear 
untuk nilai $d$ sebesar apa pun, kita katakan
$X$ berdimensi tak berhingga dan menulis $\dim \, X \coloneqq \infty$.
%
Untuk ruang vektor trivial $\{ 0 \}$, kita definisikan $\dim \, \{ 0 \} \coloneqq 0$.
\end{defn}

Setiap subhimpunan dari himpunan bebas linear jelas bebas linear.
Jadi, jika suatu himpunan memuat $d$ vektor bebas linear, himpunan itu juga memuat
himpunan yang terdiri atas $m$ vektor bebas linear untuk setiap $m \leq d$.
Selain itu, jika suatu himpunan tidak memiliki $d+1$
vektor bebas linear, tidak ada himpunan yang terdiri atas lebih dari $d+1$ vektor yang bebas
linear.
Jadi, $X$ berdimensi $d$ jika terdapat
himpunan yang terdiri atas $d$ vektor bebas linear, tetapi tidak ada
himpunan yang terdiri atas $d+1$ vektor yang bebas linear.

Tidak ada unsur dalam suatu himpunan bebas linear yang dapat bernilai nol,
dan suatu himpunan dengan satu unsur tak nol selalu bebas linear.
Secara khusus, $\{ 0 \}$ adalah satu-satunya ruang vektor berdimensi $0$.
Setiap ruang vektor lainnya memiliki dimensi positif
atau berdimensi tak berhingga.
Karena himpunan kosong bebas linear, himpunan itu merupakan basis
dari $\{ 0 \}$.

Sebagai contoh,
himpunan $Y$ yang terdiri atas dua vektor dalam
\exampleref{example:vecspr2span} merupakan basis dari $\R^2$, sehingga
$\dim \, \R^2 \geq 2$.
Sebentar lagi kita akan melihat bahwa setiap subruang vektor dari $\R^n$
berdimensi hingga, dan dimensinya kurang dari atau sama dengan $n$.
Jadi, setiap himpunan yang terdiri atas 3 vektor dalam $\R^2$ bergantung linear, dan
$\dim \, \R^2 = 2$.

Jika suatu himpunan bergantung linear, maka salah satu
vektornya merupakan kombinasi linear dari vektor-vektor lainnya.
Dalam \eqref{eq:lincomb},
jika $a_k \neq 0$, kita selesaikan untuk $x_k$:
\begin{equation*}
x_k = \frac{-a_1}{a_k} x_1 + \cdots + 
\frac{-a_{k-1}}{a_k} x_{k-1} +
\frac{-a_{k+1}}{a_k} x_{k+1} +
\cdots + 
\frac{-a_m}{a_k} x_m .
\end{equation*}
Vektor $x_k$ memiliki sedikitnya dua representasi yang berbeda
sebagai kombinasi linear vektor-vektor $\{ x_1,x_2,\ldots,x_m \}$, yaitu representasi di atas dan
$x_k$ itu sendiri. 
Sebagai contoh, himpunan $\bigl\{ (0,1), (2,3), (1,0) \bigr\}$ dalam $\R^2$ bergantung linear:
\begin{equation*}
3(0,1) - (2,3) + 2(1,0) = 0 ,
\qquad \text{sehingga} \qquad
(2,3) = 3(0,1) + 2(1,0) .
\end{equation*}

\begin{prop}
Misalkan suatu ruang vektor $X$ memiliki basis 
$B = \{ x_1, x_2, \ldots, x_n \}$.  Maka
setiap $y \in X$ memiliki representasi tunggal berbentuk
\begin{equation*}
y = \sum_{k=1}^n a_k \, x_k
\end{equation*}
untuk skalar-skalar $a_1, a_2, \ldots, a_n$ tertentu.
\end{prop}

\begin{proof}
Karena $X$ merupakan rentang dari $B$,
setiap $y \in X$ adalah kombinasi linear unsur-unsur $B$.
Misalkan
\begin{equation*}
y = \sum_{k=1}^n a_k \, x_k = \sum_{k=1}^n b_k \, x_k .
\end{equation*}
Maka
\begin{equation*}
\sum_{k=1}^n (a_k-b_k) x_k = 0 .
\end{equation*}
Karena basis tersebut bebas linear, $a_k = b_k$ untuk setiap $k$, sehingga
representasinya tunggal.
\end{proof}

Untuk $\R^n$,
kita definisikan
\emph{\myindex{basis standar}} dari $\R^n$:\glsadd{not:standardbasis}
\begin{equation*}
e_1 \coloneqq (1,0,0,\ldots,0) , \quad
e_2 \coloneqq (0,1,0,\ldots,0) , \quad \ldots, \quad
e_n \coloneqq (0,0,0,\ldots,1) .
\end{equation*}
Kita menggunakan huruf $e_k$ yang sama untuk setiap $\R^n$, dan
ruang $\R^n$ yang sedang kita gunakan dipahami dari konteks.
Perhitungan langsung menunjukkan bahwa $\{ e_1, e_2, \ldots, e_n \}$ benar-benar
merupakan basis dari $\R^n$; himpunan itu merentang $\R^n$ dan
bebas linear.  Bahkan,
\begin{equation*}
a = (a_1,a_2,\ldots,a_n) = \sum_{k=1}^n a_k \, e_k .
\end{equation*}

\begin{prop} \label{mv:dimprop}
\pagebreak[2]
Misalkan $X$ adalah ruang vektor dan $d$ bilangan bulat tak negatif.
\begin{enumerate}[(i)]
\item \label{mv:dimprop:i}
Jika $X$ direntang oleh $d$ vektor, maka $\dim \, X \leq d$.
\item \label{mv:dimprop:ii}
Jika $T$ merupakan himpunan bebas linear dan
$v \in X \setminus \spn (T)$, maka
$T \cup \{ v \}$ bebas linear.
\item \label{mv:dimprop:iii}
$\dim \, X = d$ jika dan hanya jika $X$ memiliki basis yang terdiri atas $d$
vektor.
Secara khusus, $\dim \, \R^n = n$.
\item \label{mv:dimprop:iv}
Jika $Y \subset X$ merupakan subruang vektor dan $\dim \, X = d$,
maka $\dim \, Y \leq d$.
\item \label{mv:dimprop:v}
Jika $\dim \, X = d$ dan suatu himpunan $T$ yang terdiri atas $d$ vektor merentang $X$,
maka $T$ bebas linear.
\item \label{mv:dimprop:vi}
Jika $\dim \, X = d$ dan suatu himpunan $T$ yang terdiri atas $m$ vektor
bebas linear, maka terdapat himpunan $S$ yang terdiri atas $d-m$
vektor sedemikian sehingga $T \cup S$ merupakan basis dari $X$.
\end{enumerate}
\end{prop}

Secara khusus, butir terakhir menyatakan bahwa jika $\dim X = d$ dan $T$ merupakan
himpunan yang terdiri atas $d$ vektor bebas linear, maka $T$ merentang $X$.
Perlu pula diperhatikan bahwa butir \ref{mv:dimprop:iii} menyiratkan bahwa setiap
basis dari ruang vektor berdimensi hingga memiliki jumlah unsur yang sama.

\begin{proof}
Semua pernyataan berlaku secara trivial ketika $d=0$, jadi anggap $d \geq 1$.

Kita mulai dengan \ref{mv:dimprop:i}.
Misalkan $S \coloneqq \{ x_1 , x_2, \ldots, x_d \}$ merentang $X$, dan
$T \coloneqq \{ y_1, y_2, \ldots, y_m \}$ merupakan subhimpunan bebas
linear dari $X$.  Kita hendak menunjukkan bahwa $m \leq d$.
Karena $S$ merentang $X$,
tuliskan
\begin{equation*}
y_1 = \sum_{k=1}^d a_{k,1} x_k ,
\end{equation*}
untuk bilangan-bilangan $a_{1,1},a_{2,1},\ldots,a_{d,1}$ tertentu.
Salah satu
$a_{k,1}$ tidak nol, sebab jika tidak, $y_1$ akan bernilai nol.
Tanpa mengurangi keumuman, misalkan $a_{1,1} \neq 0$.  Selesaikan
\begin{equation*}
x_1 = \frac{1}{a_{1,1}} y_1 - \sum_{k=2}^d \frac{a_{k,1}}{a_{1,1}} x_k .
\end{equation*}
Secara khusus, $\{ y_1 , x_2, \ldots, x_d \}$ merentang $X$, karena $x_1$ dapat
diperoleh dari $\{ y_1 , x_2, \ldots, x_d \}$.  Oleh karena itu, terdapat bilangan
$a_{1,2},a_{2,2},\ldots,a_{d,2}$ sedemikian sehingga
\begin{equation*}
y_2 = a_{1,2} y_1 + \sum_{k=2}^d a_{k,2} x_k .
\end{equation*}
Karena $T$ bebas linear---sehingga $\{ y_1, y_2 \}$ juga bebas
linear---salah satu $a_{k,2}$
untuk $k \geq 2$ harus tidak nol.  Tanpa mengurangi keumuman, misalkan 
$a_{2,2} \neq 0$.  Selesaikan
\begin{equation*}
x_2 = \frac{1}{a_{2,2}} y_2 - \frac{a_{1,2}}{a_{2,2}} y_1 - \sum_{k=3}^d
\frac{a_{k,2}}{a_{2,2}} x_k .
\avoidbreak
\end{equation*}
Secara khusus,
$\{ y_1 , y_2, x_3, \ldots, x_d \}$ merentang $X$.

Kita lanjutkan prosedur ini.  Jika $m < d$, pembuktian selesai.  Misalkan
$m \geq d$.
Setelah $d$ langkah, kita memperoleh bahwa 
$\{ y_1 , y_2, \ldots, y_d \}$ merentang $X$.  Setiap
vektor lain $v$ dalam $X$ merupakan kombinasi linear dari
$\{ y_1 , y_2, \ldots, y_d \}$, sehingga tidak dapat menjadi anggota $T$ karena $T$
bebas linear.  Jadi, $m = d$.

Kita lanjutkan dengan \ref{mv:dimprop:ii}.
Misalkan 
$T = \{x_1,x_2,\ldots,x_m\}$ bebas linear,
tidak merentang $X$, dan $v \in X \setminus \spn (T)$.
Misalkan 
$a_1 x_1 + a_2 x_2 + \cdots + a_m x_m + a_{m+1} v = 0$
untuk skalar-skalar $a_1,a_2,\ldots,a_{m+1}$ tertentu.
Jika $a_{m+1} \neq 0$, maka $v$ akan menjadi kombinasi linear dari $T$,
jadi $a_{m+1} = 0$.  Kemudian, karena $T$ bebas linear, $a_1=a_2=\cdots=a_m = 0$.
Jadi, $T \cup \{ v \}$ bebas linear.

Kita beralih ke \ref{mv:dimprop:iii}.
Jika $\dim \, X = d$,
maka harus ada suatu himpunan bebas linear $T$ yang terdiri atas $d$ vektor,
dan $T$ harus merentang $X$; jika tidak, kita dapat memilih himpunan vektor bebas
linear yang lebih besar melalui \ref{mv:dimprop:ii}.
Jadi, kita memiliki basis yang terdiri atas $d$ vektor.
Sebaliknya, jika kita memiliki basis yang terdiri atas $d$ vektor,
dimensinya sedikitnya $d$ karena suatu basis bebas linear.
Suatu basis juga merentang $X$, sehingga
berdasarkan \ref{mv:dimprop:i}, kita tahu bahwa dimensinya paling besar $d$.
Oleh sebab itu, dimensi $X$ harus sama dengan $d$.
Pernyataan \myquote{secara khusus} diperoleh dengan memperhatikan bahwa
$\{ e_1, e_2, \ldots, e_n \}$ merupakan basis dari $\R^n$.

Untuk membuktikan \ref{mv:dimprop:iv},
misalkan $Y \subset X$ merupakan subruang vektor,
dengan $\dim \, X = d$.  Karena $X$ tidak dapat memuat $d+1$ vektor bebas
linear, demikian pula $Y$.

Untuk \ref{mv:dimprop:v}, misalkan $T$ merupakan himpunan yang terdiri atas $m$ vektor
yang bergantung linear
dan merentang $X$.
Kita akan menunjukkan bahwa $m > d$.
Salah satu
vektor merupakan kombinasi linear dari vektor-vektor lainnya.  Jika kita menghapusnya
dari $T$, kita memperoleh himpunan yang terdiri atas $m-1$ vektor yang masih merentang $X$.
Oleh karena itu,
$d = \dim \, X \leq m-1$ menurut~\ref{mv:dimprop:i}.

Untuk \ref{mv:dimprop:vi}, misalkan $T = \{ x_1, x_2, \ldots, x_m \}$ merupakan
himpunan bebas linear.  Pertama, $m \leq d$ menurut definisi dimensi.
Jika $m=d$, himpunan $T$ harus merentang $X$ seperti dalam bukti \ref{mv:dimprop:iii};
jika tidak, kita dapat menambahkan vektor lain pada $T$.
Jika $m < d$,  $T$ tidak dapat merentang $X$ menurut \ref{mv:dimprop:iii}.
Jadi, carilah $v$ yang tidak berada dalam rentang~$T$.
Berdasarkan \ref{mv:dimprop:ii}, himpunan $T \cup \{ v \}$ merupakan
himpunan bebas linear yang terdiri atas $m+1$ unsur.
Oleh karena itu, kita ulangi prosedur ini sebanyak $d-m$ kali untuk
memperoleh himpunan yang terdiri atas $d$ vektor bebas linear.
Sekali lagi, himpunan tersebut harus merentang $X$; jika tidak, kita dapat menambahkan satu vektor lagi.
\end{proof}

\subsection{Pemetaan linear}

Ketika $Y \neq \R$,
suatu fungsi $f \colon X \to Y$ sering disebut
\emph{\myindex{pemetaan}} atau \emph{\myindex{peta}}
alih-alih \emph{fungsi}.

\begin{defn}
Suatu pemetaan $A \colon X \to Y$ dari ruang vektor $X$ ke $Y$
disebut \emph{\myindex{linear}} (kita juga mengatakan bahwa $A$ adalah suatu
\emph{\myindex{transformasi linear}}\index{transformasi, linear}
atau \emph{\myindex{operator linear}}\index{operator, linear})
jika untuk setiap $a \in \R$ dan setiap $x,y \in X$,
\begin{equation*}
A(a x) = a A(x) \qquad \text{dan} \qquad A(x+y) = A(x)+A(y) .
\end{equation*}
Kita biasanya menulis $Ax$ alih-alih $A(x)$ jika $A$ linear.
Jika $A$ injektif dan surjektif, kita mengatakan bahwa $A$
\emph{invertibel}\index{transformasi linear invertibel},
dan menotasikan inversnya dengan $A^{-1}$.
Jika $A \colon X \to X$ linear, kita mengatakan bahwa $A$ adalah suatu
\emph{operator linear pada $X$}.

Kita menulis $L(X,Y)$\glsadd{not:LXY} untuk himpunan pemetaan linear dari $X$ ke
$Y$, dan $L(X)$\glsadd{not:LX} untuk himpunan operator linear pada $X$.
Jika $a \in \R$ dan $A,B \in L(X,Y)$, definisikan
pemetaan $aA$ dan $A+B$ dengan
\begin{equation*}
(aA)(x) \coloneqq aAx ,
\qquad
(A+B)(x) \coloneqq Ax + Bx .
\end{equation*}
Jika $A \in L(Y,Z)$ dan $B \in L(X,Y)$, definisikan
pemetaan $AB \colon X \to Z$ sebagai komposisi $A \circ B$,
\begin{equation*}
ABx \coloneqq A(Bx) .
\end{equation*}
Terakhir, notasikan dengan $I \in L(X)$ suatu \emph{\myindex{identitas}}: 
operator linear yang memenuhi $Ix = x$ untuk setiap $x$.
\end{defn}

\begin{prop} \label{prop:LXYvs}
Misalkan $X$, $Y$, dan $Z$ adalah ruang vektor.
\begin{enumerate}[(i)]
\item Jika $A \in L(X,Y)$, maka $A0 = 0$.
\item Jika $A,B \in L(X,Y)$, maka $A+B \in L(X,Y)$.
\item Jika $A \in L(X,Y)$ dan $a \in \R$, maka $aA \in L(X,Y)$.
\item Jika $A \in L(Y,Z)$ dan $B \in L(X,Y)$, maka $AB \in L(X,Z)$. 
\item Jika $A \in L(X,Y)$ invertibel, maka $A^{-1} \in L(Y,X)$.
\end{enumerate}
\end{prop}

Secara khusus, $L(X,Y)$ merupakan ruang vektor, dengan $0 \in L(X,Y)$ sebagai
pemetaan linear yang memetakan setiap elemen ke $0$.
Karena $L(X)$ bukan hanya ruang vektor, melainkan juga
memiliki operasi perkalian (komposisi),
ruang ini disebut suatu \emph{\myindex{aljabar}}.

\begin{proof}
Empat butir pertama kita serahkan sebagai latihan singkat,
\exerciseref{exercise:LXYvs}.
Mari kita buktikan butir terakhir.
Misalkan $a \in \R$ dan $y \in Y$.  Karena $A$ surjektif, terdapat 
$x \in X$ sedemikian sehingga $y = Ax$.
Karena pemetaan tersebut juga injektif, $A^{-1}(Az) = z$ untuk setiap $z \in X$.
Dengan demikian,
\begin{equation*}
A^{-1}(ay)
=
A^{-1}(aAx)
=
A^{-1}\bigl(A(ax)\bigr)
= ax
= aA^{-1}(y).
\end{equation*}
Serupa dengan itu, misalkan $y_1,y_2 \in Y$ dan $x_1, x_2 \in X$ sedemikian sehingga
$Ax_1 = y_1$ dan 
$Ax_2 = y_2$; maka
\begin{equation*}
A^{-1}(y_1+y_2)
=
A^{-1}(Ax_1+Ax_2)
=
A^{-1}\bigl(A(x_1+x_2)\bigr)
= x_1+x_2
= A^{-1}(y_1) + A^{-1}(y_2). \qedhere
\end{equation*}
\end{proof}

\begin{prop} \label{mv:lindefonbasis}
Jika $A \in L(X,Y)$ linear, maka pemetaan tersebut sepenuhnya ditentukan
oleh nilai-nilainya pada suatu basis dari $X$.
Lebih lanjut, jika $B$ adalah suatu basis dari $X$,
maka setiap fungsi $\widetilde{A} \colon B \to Y$ dapat diperluas menjadi fungsi
linear $A$ pada $X$.
\end{prop}

Kita hanya membuktikan proposisi ini untuk ruang berdimensi hingga, sebab kita
tidak memerlukan ruang berdimensi tak berhingga.%
\footnote{Untuk ruang berdimensi tak berhingga, pembuktiannya pada dasarnya sama, tetapi
sedikit lebih rumit untuk dituliskan.
Terlebih lagi, kita bahkan belum mendefinisikan apa yang dimaksud dengan basis untuk
ruang berdimensi tak berhingga.}

\begin{proof}
Misalkan $\{ x_1, x_2, \ldots, x_n \}$ adalah suatu basis dari $X$, dan tetapkan 
$y_k \coloneqq A x_k$.  Setiap $x \in X$ memiliki penyajian unik
\begin{equation*}
x = \sum_{k=1}^n b_k \, x_k
\end{equation*}
untuk beberapa bilangan $b_1,b_2,\ldots,b_n$.  Berdasarkan linearitas,
\begin{equation*}
Ax = 
A\sum_{k=1}^n b_k x_k
=
\sum_{k=1}^n b_k \, Ax_k
=
\sum_{k=1}^n b_k \, y_k .
\end{equation*}
Pernyataan \myquote{lebih lanjut} diperoleh dengan menetapkan $y_k \coloneqq \widetilde{A}(x_k)$,
kemudian, untuk 
$x = \sum_{k=1}^n b_k \, x_k$,
mendefinisikan perluasannya sebagai
$A(x) \coloneqq \sum_{k=1}^n b_k \, y_k$.  Fungsi ini terdefinisi dengan baik berkat
keunikan penyajian $x$.
Pembaca dipersilakan memeriksa bahwa $A$ linear.
\end{proof}

Untuk suatu pemetaan linear, cukup memeriksa injektivitas di titik asal.
Artinya, jika satu-satunya $x$ yang memenuhi $Ax =0$ adalah $x=0$, maka $A$ injektif,
sebab jika $Ay=Az$, maka $A(y-z) = 0$.  Atas alasan ini, 
kita sering mempelajari \emph{\myindex{ruang nol}} dari $A$, yaitu
$\{ x \in X : Ax = 0 \}$.
Untuk ruang vektor berdimensi hingga (dan hanya dalam dimensi hingga),
kita memperoleh kasus khusus berikut dari
teorema rank--nulitas dalam aljabar linear.

\begin{prop} \label{mv:prop:lin11onto}
Jika $X$ adalah ruang vektor berdimensi hingga dan $A \in L(X)$, maka $A$ injektif jika dan hanya jika pemetaan tersebut surjektif.
\end{prop}

\begin{proof}
Misalkan $\{ x_1,x_2,\ldots,x_n \}$ adalah suatu basis dari $X$.
Pertama, andaikan $A$ injektif.  Misalkan $c_1,c_2,\ldots,c_n$ adalah skalar sedemikian sehingga
\begin{equation*}
0 =
\sum_{k=1}^n c_k \, Ax_k =
A\sum_{k=1}^n c_k \, x_k 
.
\end{equation*}
Karena $A$ injektif,
satu-satunya vektor yang dipetakan ke 0 adalah vektor 0 itu sendiri.  
Oleh karena itu,
\begin{equation*}
0 =
\sum_{k=1}^n c_k \, x_k,
\end{equation*}
maka $c_k = 0$ untuk setiap $k$ karena $\{ x_1,x_2,\ldots,x_n \}$ adalah suatu basis.
Jadi, $\{ Ax_1, Ax_2, \ldots, Ax_n \}$ bebas linear.
Berdasarkan \propref{mv:dimprop}
dan fakta bahwa dimensinya adalah $n$, kita menyimpulkan bahwa
$\{ Ax_1, Ax_2, \ldots, Ax_n \}$ merentang $X$.  Akibatnya, $A$ surjektif,
sebab setiap $y \in X$ dapat ditulis sebagai
\begin{equation*}
y = \sum_{k=1}^n a_k \, Ax_k =
A\sum_{k=1}^n a_k \, x_k .
\avoidbreak
\end{equation*}

Untuk arah sebaliknya, andaikan $A$ surjektif.
Andaikan bahwa untuk suatu pilihan skalar
$c_1,c_2,\ldots,c_n$,
\begin{equation*}
0 = A\sum_{k=1}^n c_k \, x_k =
\sum_{k=1}^n c_k \, Ax_k .
\end{equation*}
Karena $A$ surjektif dan linear, citra basis tersebut merentang ruang sasaran;
$\{ Ax_1, Ax_2, \ldots, Ax_n \}$ merentang $X$.
Maka, berdasarkan \propref{mv:dimprop}, himpunan tersebut bebas linear,
dan $c_k = 0$ untuk setiap $k$.  Dengan kata lain, jika $Ax = 0$, maka $x=0$.
Jadi, $A$ injektif.
\end{proof}

Pembuktian proposisi berikut kita serahkan sebagai latihan.

\begin{prop} \label{prop:LXYfinitedim}
Jika $X$ dan $Y$ adalah ruang vektor berdimensi hingga, maka $L(X,Y)$
juga berdimensi hingga.
\end{prop}

Kita dapat mengidentifikasi ruang vektor berdimensi hingga
$X$ yang berdimensi $n$ dengan $\R^n$, asalkan kita menetapkan suatu basis
$\{ x_1, x_2, \ldots, x_n \}$ pada $X$.  Artinya, kita mendefinisikan pemetaan linear
bijektif $A \in L(X,\R^n)$ dengan
$Ax_k \coloneqq e_k$, dengan $\{ e_1, e_2, \ldots, e_n \}$ sebagai basis standar
di $\R^n$.  Kita memperoleh
korespondensi\glsadd{not:mapsto}
\begin{equation*}
\sum_{k=1}^n c_k \, x_k \, \in X
\quad
\overset{A}{\mapsto}
\quad
(c_1,c_2,\ldots,c_n) \, \in \R^n .
\end{equation*}

\subsection{Kekonveksan}

Suatu himpunan bagian $U$ dari ruang vektor disebut \emph{\myindex{konveks}}
jika untuk setiap $x,y \in U$, ruas garis dari
$x$ ke $y$ terletak dalam $U$.  Dengan kata lain, \emph{\myindex{kombinasi konveks}}
$(1-t)x+ty$ berada dalam $U$ untuk semua $t \in [0,1]$.
Kita menulis $[x,y]$ untuk ruas garis ini.\glsadd{not:segmentinRn}
Lihat \figureref{mv:convexcomb}.

\begin{myfigureht}
\myincludepdft{convexset}{%
Sebuah daerah oval berarsir berlabel U.  Dua titik,
x dan y, ditandai di dalamnya, beserta ruas garis di antara keduanya yang
seluruhnya terletak di dalam U.  Sebuah titik di bagian tengah ruas tersebut
diberi label satu dikurangi t dikali x ditambah t dikali y.}
\caption{Kekonveksan.\label{mv:convexcomb}}
\end{myfigureht}

Dalam $\R$, himpunan konveks tepat merupakan interval, yang juga
tepat merupakan himpunan terhubung.
Dalam dua dimensi atau lebih,
terdapat banyak himpunan terhubung yang tidak konveks.  Sebagai contoh,
himpunan $\R^2 \setminus \{0\}$ terhubung,
tetapi tidak konveks---untuk sebarang $x \in \R^2 \setminus \{0\}$ dengan $y \coloneqq -x$,
kita memperoleh
$(\nicefrac{1}{2})x + (\nicefrac{1}{2})y = 0$, yang tidak berada dalam himpunan tersebut.
Bola (dalam metrik standar) di $\R^n$ bersifat konveks.  Hasil ini
cukup berguna untuk dinyatakan sebagai proposisi,
tetapi pembuktiannya kita tinggalkan sebagai latihan.

\begin{prop}
Misalkan $x \in \R^n$ dan $r > 0$.  Bola $B(x,r) \subset \R^n$
%(menggunakan metrik standar pada $\R^n$)
bersifat konveks.
\end{prop}

\begin{example}
Kombinasi konveks, khususnya, merupakan kombinasi linear.
Jadi setiap subruang vektor $V$ dari ruang vektor $X$ bersifat konveks.
\end{example}

\begin{example}
%Contoh yang agak lebih rumit diberikan sebagai berikut.
Misalkan $C([0,1],\R)$ adalah ruang vektor fungsi kontinu bernilai real pada $[0,1]$.
Misalkan $V \subset C([0,1],\R)$ adalah himpunan fungsi $f$ sedemikian sehingga
\begin{equation*}
\int_0^1 f(x)\,dx \leq 1 \qquad \text{dan} \qquad
f(x) \geq 0 \enspace \text{untuk semua } x \in [0,1] .
\end{equation*}
Maka $V$ konveks.  Ambil $t \in [0,1]$, dan perhatikan bahwa jika $f,g \in V$,
maka $(1-t) f(x) + t g(x) \geq 0$ untuk semua $x$.  Lebih lanjut,
\begin{equation*}
\int_0^1 \bigl((1-t)f(x) + t g(x)\bigr) \,dx
=
(1-t) \int_0^1 f(x) \,dx
+ t \int_0^1 g(x) \,dx \leq 1 .
\end{equation*}
Perhatikan bahwa $V$ bukan subruang vektor dari $C([0,1],\R)$.  Fungsi
$f(x) \coloneqq 1$ berada dalam $V$, tetapi $2f$ dan $-f$ tidak.
\end{example}

\begin{prop} \label{prop:intersectionconvex}
Irisan dua himpunan konveks bersifat konveks.  Bahkan,
jika $\{ C_\lambda \}_{\lambda \in I}$ adalah
sebarang koleksi himpunan konveks dalam suatu ruang vektor, maka
\begin{equation*}
C \coloneqq \bigcap_{\lambda \in I} C_\lambda
\qquad \text{bersifat konveks.}
\end{equation*}
\end{prop}

\begin{proof}
Jika $x, y \in C$, maka $x,y \in C_\lambda$ untuk semua
$\lambda \in I$, sehingga jika $t \in [0,1]$, maka $(1-t)x + ty \in
C_\lambda$ untuk semua $\lambda \in I$.  Oleh karena itu, $(1-t)x + ty \in C$ dan $C$
bersifat konveks.
\end{proof}

Salah satu konstruksi yang berguna dengan menggunakan irisan himpunan konveks ialah
\emph{\myindex{selubung konveks}}.  Untuk himpunan bagian $S$ dari ruang
vektor $X$, definisikan selubung konveks $S$ sebagai irisan semua himpunan konveks
yang memuat $S$:
\begin{equation*}
\operatorname{co}(S) \coloneqq
\bigcap \{ C \subset X : S \subset C, \text{ dan } C \text{ konveks} \} .
\end{equation*}
Dengan kata lain, selubung konveks adalah himpunan konveks terkecil yang memuat $S$.
Menurut \propref{prop:intersectionconvex}, irisan himpunan-himpunan konveks
bersifat konveks.
Jadi selubung konveks bersifat konveks.

\begin{example}
Selubung konveks dari $\{ 0, 1 \}$ dalam $\R$ adalah $[0,1]$.  Bukti:
Himpunan konveks yang memuat 0 dan 1 harus memuat $[0,1]$, sehingga
$[0,1] \subset \operatorname{co}(\{0,1\})$.
Himpunan $[0,1]$ konveks dan memuat $\{0,1\}$, sehingga
$\operatorname{co}(\{0,1\}) \subset [0,1]$.
\end{example}

Pemetaan linear mempertahankan himpunan konveks.  Jadi, dalam arti tertentu, himpunan
konveks adalah jenis himpunan yang tepat ketika membahas pemetaan linear atau perubahan
koordinat.

\begin{prop}
Misalkan $X,Y$ ruang vektor, $A \in L(X,Y)$, dan $C \subset X$ konveks.
Maka $A(C)$ konveks.
\end{prop}

\begin{proof}
Ambil dua titik $p,q \in A(C)$.  Pilih $u,v \in C$ sedemikian sehingga
$Au = p$ dan $Av=q$.  Karena $C$ konveks,
$(1-t)u+t v \in C$
untuk semua $t \in [0,1]$.  Jadi
\begin{equation*}
(1-t)p+t q 
=
(1-t)Au+tAv
=
A\bigl((1-t)u+tv\bigr)
\in A(C) .  \qedhere
\end{equation*}
\end{proof}


\subsection{Latihan}

\begin{exercise}
Tunjukkan bahwa dalam $\R^n$
(dengan metrik Euklides standar),
untuk setiap $x \in \R^n$ dan setiap $r > 0$,
bola $B(x,r)$ bersifat konveks.
\end{exercise}


\begin{exercise}
Periksalah bahwa $\R^n$ merupakan ruang vektor.
\end{exercise}

\begin{exercise}
Misalkan $X$ ruang vektor.
Buktikan bahwa suatu himpunan berhingga vektor $\{ x_1,x_2,\ldots,x_n \} \subset X$
bebas linear jika dan hanya jika untuk setiap $k=1,2,\ldots,n$
\begin{equation*}
\spn\bigl( \{ x_1,\ldots,x_{k-1},x_{k+1},\ldots,x_n \}\bigr) \subsetneq
\spn\bigl( \{ x_1,x_2,\ldots,x_n \}\bigr) .
\end{equation*}
Artinya, rentang himpunan setelah satu vektor dihapus benar-benar lebih kecil.
\end{exercise}

\begin{exercise} \label{exercise:intzerosubspace}
Tunjukkan bahwa himpunan $X \subset C([0,1],\R)$ yang terdiri atas fungsi-fungsi
dengan $\int_0^1 f = 0$ merupakan subruang vektor.
Bandingkan dengan \exerciseref{exercise:intoneconvex}.
\end{exercise}

\begin{exercise}[Menantang]
Buktikan bahwa $C([0,1],\R)$ merupakan ruang vektor berdimensi tak hingga
dengan operasi yang didefinisikan secara alami:
$s=f+g$ dan $m=af$ didefinisikan oleh
$s(x) \coloneqq f(x)+g(x)$ dan
$m(x) \coloneqq a f(x)$.
Petunjuk: Untuk dimensinya, pikirkan fungsi-fungsi yang tidak nol hanya
pada interval $(\nicefrac{1}{n+1},\nicefrac{1}{n})$.
\end{exercise}

\begin{exercise} \label{exercise:continuouskernel}
Misalkan $k \colon [0,1]^2 \to \R$ kontinu.  Tunjukkan bahwa
$L \colon C([0,1],\R) \to C([0,1],\R)$ yang didefinisikan oleh
\begin{equation*}
Lf(y) \coloneqq \int_0^1 k(x,y)f(x)\,dx
\end{equation*}
merupakan operator linear.  Pertama, tunjukkan bahwa $L$ terdefinisi dengan baik dengan
menunjukkan bahwa $Lf$ kontinu setiap kali $f$ kontinu, lalu
tunjukkan bahwa $L$ linear.
\end{exercise}

\begin{exercise}
Misalkan $\sP_n$ ruang vektor polinom dalam satu variabel yang berderajat $n$
atau kurang.  Tunjukkan bahwa $\sP_n$ merupakan ruang vektor berdimensi $n+1$.
\end{exercise}

\begin{exercise}
Misalkan $\R[t]$ ruang vektor polinom dalam satu variabel $t$.  Misalkan
$D \colon \R[t] \to \R[t]$ operator turunan (turunan terhadap $t$).
Tunjukkan bahwa $D$ merupakan operator linear.
\end{exercise}

\begin{exercise}
Mari kita tunjukkan bahwa \propref{mv:prop:lin11onto} hanya berlaku dalam dimensi
berhingga.  Ambil ruang polinom $\R[t]$ dan definisikan operator $A \colon \R[t] \to \R[t]$
dengan $A\bigl(P(t)\bigr) \coloneqq tP(t)$.  Tunjukkan bahwa $A$ linear dan satu-satu, tetapi
tidak surjektif.
\end{exercise}

\begin{exercise}
Selesaikan pembuktian \propref{mv:lindefonbasis} untuk kasus berdimensi hingga.
Artinya, misalkan
$\{ x_1, x_2,\ldots x_n \}$ adalah basis bagi $X$,
$\{ y_1, y_2,\ldots y_n \} \subset Y$, dan definisikan suatu
fungsi
\begin{equation*}
A(x) \coloneqq \sum_{k=1}^n b_k y_k, \qquad \text{jika} \quad x=\sum_{k=1}^n b_k x_k .
\end{equation*}
Buktikan bahwa $A \colon X \to Y$ linear.
\end{exercise}


\begin{exercise}
Buktikan \propref{prop:LXYfinitedim}.  Petunjuk: Suatu transformasi linear ditentukan oleh
nilainya pada sebuah basis.  Jadi, untuk dua basis
$\{ x_1,\ldots,x_n \}$ dan
$\{ y_1,\ldots,y_m \}$ masing-masing bagi $X$ dan $Y$, tinjau operator linear
$A_{jk}$ yang memenuhi $A_{jk} x_j = y_k$, serta
$A_{jk} x_\ell = 0$ jika $\ell \neq j$.
\end{exercise}

\begin{samepage}
\begin{exercise}[Mudah]
Misalkan $X$ dan $Y$ ruang vektor dan $A \in L(X,Y)$ suatu pemetaan
linear.
\begin{enumerate}[a)]
\item
Tunjukkan bahwa ruang nol $N \coloneqq \{ x \in X : Ax = 0 \}$ merupakan
ruang vektor.
\item
Tunjukkan bahwa daerah hasil $R \coloneqq \{ y \in Y : Ax = y \text{ untuk suatu } x \in X \}$ merupakan
ruang vektor.
\end{enumerate}
\end{exercise}
\end{samepage}

\begin{exercise}[Mudah]
Berikan contoh yang menunjukkan bahwa gabungan himpunan-himpunan konveks belum tentu konveks.
\end{exercise}

\begin{exercise}
Hitung selubung konveks dari himpunan tiga titik
$\bigl\{ (0,0), (0,1), (1,1) \bigr\}$
dalam $\R^2$.
\end{exercise}

\begin{exercise}
Tunjukkan bahwa himpunan $\bigl\{ (x,y) \in \R^2 : y > x^2 \bigr\}$ bersifat konveks.
\end{exercise}

\begin{exercise} \label{exercise:intoneconvex}
Tunjukkan bahwa himpunan $X \subset C([0,1],\R)$ yang terdiri atas fungsi-fungsi
dengan $\int_0^1 f = 1$ merupakan himpunan konveks, tetapi bukan subruang vektor.
Bandingkan dengan \exerciseref{exercise:intzerosubspace}.
\end{exercise}


\begin{exercise}
Tunjukkan bahwa setiap himpunan konveks dalam $\R^n$ terhubung dalam topologi
standar pada $\R^n$.
\end{exercise}

\begin{exercise}
Misalkan $K \subset \R^2$ himpunan konveks sedemikian sehingga satu-satunya titik
berbentuk $(x,0)$ dalam $K$ adalah titik $(0,0)$.  Misalkan pula bahwa
$(0,1) \in K$ dan $(1,1) \in K$.  Tunjukkan bahwa jika $(x,y) \in K$ dan $x
\neq 0$, maka $y > 0$.
\end{exercise}

\begin{exercise}
Buktikan bahwa sebarang irisan subruang-subruang vektor
merupakan subruang vektor.
Artinya, jika $X$ ruang vektor dan
$\{ V_\lambda \}_{\lambda \in I}$ adalah
sebarang koleksi subruang vektor dari $X$,
maka
$\bigcap_{\lambda \in I} V_\lambda$ merupakan subruang vektor dari $X$.
\end{exercise}

\begin{exercise}[Mudah] \label{exercise:LXYvs}
Selesaikan pembuktian \propref{prop:LXYvs}, yaitu buktikan empat
butir pertama dalam proposisi tersebut.
\end{exercise}

%%%%%%%%%%%%%%%%%%%%%%%%%%%%%%%%%%%%%%%%%%%%%%%%%%%%%%%%%%%%%%%%%%%%%%%%%%%%%%

\sectionnewpage
\section{Analisis dengan ruang vektor}
\label{sec:normsmatsdets}

%mbxINTROSUBSECTION

\sectionnotes{3 kuliah}

\subsection{Norma}

Mari kita mulai dengan mengukur besar vektor, dan dengan demikian juga jarak.

\begin{defn}
Jika $X$ adalah ruang vektor, maka suatu fungsi
$\snorm{\cdot} \colon X \to \R$ disebut
\emph{\myindex{norma}}\glsadd{not:norm} jika
\begin{enumerate}[(i)]
\item \label{defn:norm:i} $\snorm{x} \geq 0$, dengan $\snorm{x}=0$ jika dan hanya jika $x=0$.
\item \label{defn:norm:ii} $\snorm{cx} = \sabs{c} \, \snorm{x}$ untuk setiap $c \in \R$ dan $x \in X$.
\item \label{defn:norm:iii} $\snorm{x+y} \leq \snorm{x}+\snorm{y}$ untuk setiap
$x,y \in X$
\qquad (\emph{ketaksamaan segitiga})\index{ketaksamaan segitiga!norma}.
\end{enumerate}
Ruang vektor yang dilengkapi dengan suatu norma disebut
\emph{\myindex{ruang vektor bernorma}}.
\end{defn}

Untuk suatu norma (sembarang norma) pada ruang vektor $X$,
definisikan jarak $d(x,y) \coloneqq \snorm{x-y}$, yang
menjadikan $X$ suatu ruang metrik (latihan).
Jadi, semua yang telah Anda ketahui tentang ruang metrik berlaku bagi ruang vektor bernorma.
Sebelum mendefinisikan norma baku pada $\R^n$, kita
mendefinisikan 
\emph{\myindex{hasil kali titik}} skalar baku pada $\R^n$.
Untuk 
$x=(x_1,x_2,\ldots,x_n) \in \R^n$
dan $y=(y_1,y_2,\ldots,y_n) \in \R^n$, definisikan\glsadd{not:dotprod}
\begin{equation*}
x \cdot y \coloneqq \sum_{k=1}^n x_k\, y_k .
\end{equation*}
Hasil kali titik bersifat linear terhadap setiap variabel
secara terpisah---dalam bahasa yang lebih canggih, hasil kali ini \emph{\myindex{bilinear}}.
Artinya,
jika $y$ ditetapkan, pemetaan $x \mapsto x \cdot y$ merupakan pemetaan linear dari
$\R^n$ ke $\R$.  Demikian pula, jika $x$ ditetapkan,
$y \mapsto x \cdot y$ bersifat linear.
Hasil kali ini \emph{\myindex{simetris}}, dalam arti $x \cdot y = y \cdot x$.
Definisikan \emph{\myindex{norma Euklides}} sebagai\glsadd{not:normRn}
\begin{equation*}
\snorm{x} \coloneqq \snorm{x}_{\R^n} \coloneqq \sqrt{x \cdot x} = \sqrt{(x_1)^2+(x_2)^2 + \cdots + (x_n)^2}.
\end{equation*}
Biasanya kita cukup menulis $\snorm{x}$.
Hanya dalam keadaan yang jarang, ketika perlu
ditekankan bahwa yang dimaksud adalah norma Euklides, kita akan menulis
$\snorm{x}_{\R^n}$.
Kecuali dinyatakan lain, jika kita membahas $\R^n$ sebagai ruang vektor bernorma,
yang dimaksud adalah norma Euklides baku.
Mudah dilihat bahwa norma Euklides memenuhi \ref{defn:norm:i} dan
\ref{defn:norm:ii}.  Untuk membuktikan
bahwa \ref{defn:norm:iii} berlaku, ketaksamaan kuncinya
adalah ketaksamaan Cauchy--Schwarz
yang telah kita jumpai sebelumnya.  Karena ketaksamaan ini sangat penting, kita menyatakan dan
membuktikan versi yang sedikit lebih kuat dengan menggunakan notasi bab ini.

\begin{thm}[\myindex{Ketaksamaan Cauchy--Schwarz}]
Jika $x, y \in \R^n$, maka
\begin{equation*}
\sabs{x \cdot y} \leq \snorm{x} \, \snorm{y} = \sqrt{x\cdot x}\, \sqrt{y\cdot y},
\avoidbreak
\end{equation*}
dengan kesamaan jika dan hanya jika $x = \lambda y$ atau $y = \lambda x$ untuk suatu
$\lambda \in \R$.
\pagebreak[2]
\end{thm}

\begin{proof}
Jika $x=0$ atau $y = 0$, maka teorema berlaku secara trivial.
Jadi, andaikan $x \neq 0$ dan $y \neq 0$.

Jika $x$ merupakan kelipatan skalar dari $y$, yaitu $x = \lambda y$ untuk suatu
$\lambda \in \R$, maka teorema berlaku dengan kesamaan:
\begin{equation*}
\sabs{ x \cdot y } =
\sabs{\lambda y \cdot y} = \sabs{\lambda} \, \sabs{y\cdot y} =
\sabs{\lambda} \, \snorm{y}^2 = \snorm{\lambda y} \, \snorm{y}
= \snorm{x} \, \snorm{y} .
\end{equation*}

Dengan menetapkan $x$ dan $y$,
$\snorm{x+ty}^2$
merupakan polinom kuadrat
sebagai fungsi dari $t$:
\begin{equation*}
%mbxlatex \begin{aligned}
\snorm{x+ty}^2
=
(x+ty) \cdot (x+ty)
%mbxlatex &
=
x \cdot x + x \cdot ty + ty \cdot x + ty \cdot ty
%mbxlatex \\
%mbxlatex &
=
\snorm{x}^2 + 2t(x \cdot y) + t^2 \snorm{y}^2 .
%mbxlatex \end{aligned}
\end{equation*}
Jika $x$ bukan kelipatan skalar dari $y$, maka 
$\snorm{x+ty}^2 > 0$ untuk setiap $t$.  Jadi, polinom $\snorm{x+ty}^2$
tidak pernah bernilai nol.
Aljabar elementer menyatakan bahwa diskriminannya harus negatif:
\begin{equation*}
4 {(x \cdot y)}^2 - 4 \snorm{x}^2\snorm{y}^2 < 0.
\end{equation*}
Dengan kata lain, ${(x \cdot y)}^2 < \snorm{x}^2\snorm{y}^2$.
\end{proof}

Butir \ref{defn:norm:iii}, yaitu ketaksamaan segitiga dalam $\R^n$,
mengikuti perhitungan berikut:
\begin{equation*}
\snorm{x+y}^2 
=
x \cdot x + y \cdot y + 2 (x \cdot y)
\leq
\snorm{x}^2 + \snorm{y}^2 + 2 \bigl(\snorm{x} \,\snorm{y}\bigr)
=
{\bigl(\snorm{x} + \snorm{y}\bigr)}^2 .
\end{equation*}

Jarak
$d(x,y) \coloneqq \snorm{x-y}$ merupakan jarak baku
(metrik baku) pada $\R^n$ yang
kita gunakan ketika membahas ruang metrik.

\begin{defn}
Misalkan $A \in L(X,Y)$.  Definisikan
\begin{equation*}
\snorm{A} \coloneqq
\sup \bigl\{ \snorm{Ax} : x \in X \text{ dengan } \snorm{x} = 1 \bigr\} .
\end{equation*}
Bilangan $\snorm{A}$ (yang mungkin bernilai $\infty$)
disebut \emph{\myindex{norma operator}}.  Di bawah ini akan kita lihat
bahwa pada ruang berdimensi hingga, besaran tersebut memang merupakan norma pada $L(X,Y)$.
Sekali lagi, jika perlu ditekankan norma mana yang dimaksud, kita dapat
menuliskannya sebagai $\snorm{A}_{L(X,Y)}$.\glsadd{not:normLXY}
\end{defn}

Sebagai contoh, jika $X=\R^1$ dengan norma $\snorm{x}=\sabs{x}$, elemen-elemen $L(X)$
merupakan pemetaan berupa perkalian dengan skalar, $x \mapsto ax$, dan kita mengidentifikasikan $a \in \R$ dengan
elemen $L(X)$ yang bersesuaian.  Jika $\snorm{x}=\sabs{x}=1$, maka $\sabs{a x} = \sabs{a}$, sehingga
norma operator dari $a$ adalah $\sabs{a}$.

Berdasarkan linearitas,
$\norm{A \frac{x}{\snorm{x}}} = \frac{\snorm{Ax}}{\snorm{x}}$
untuk setiap $x \in X$ yang tak nol.
Vektor $\frac{x}{\snorm{x}}$ memiliki norma 1.
Oleh karena itu,
\begin{equation*}
\snorm{A} =
\sup \bigl\{ \snorm{Ax} : x \in X \text{ dengan } \snorm{x} = 1 \bigr\}
=
\sup_{\substack{x \in X\\x\neq 0}} \frac{\snorm{Ax}}{\snorm{x}} .
\end{equation*}
Hal ini menyiratkan, dengan mengandaikan $\snorm{A} \neq \infty$ untuk menghindari persoalan teknis ketika $x=0$,
bahwa untuk setiap $x \in X$,
\begin{equation*}
\snorm{Ax} \leq \snorm{A}  \snorm{x} .
\avoidbreak
\end{equation*}
%We will use this inequality a lot.
Sebaliknya, jika dapat ditunjukkan bahwa $\snorm{Ax} \leq C \snorm{x}$ untuk setiap $x$,
maka $\snorm{A} \leq C$.

Dari definisi, tidak sulit untuk melihat bahwa $\snorm{A} = 0$ jika dan
hanya jika $A = 0$, dengan $A=0$ berarti bahwa $A$ memetakan setiap vektor ke vektor nol.
Norma operator dari operator identitas juga tidak sulit dihitung:
\begin{equation*}
\snorm{I} =
\sup_{\substack{x \in X\\x\neq 0}} \frac{\snorm{Ix}}{\snorm{x}} 
=
\sup_{\substack{x \in X\\x\neq 0}} \frac{\snorm{x}}{\snorm{x}} 
= 1.
\end{equation*}
Norma operator tidak selalu mudah dihitung hanya dengan menggunakan definisinya,
dan tidak pula mudah dibaca langsung dari bentuk operatornya.
Perhatikan $\R^2$ dan operator $A \in L(\R^2)$ yang memetakan
$(x,y)$ ke $(x+y,2x)$.  
Vektor-vektor bernorma satu dapat ditulis sebagai
$\bigl(\pm t, \pm \sqrt{1-t^2}\bigr)$
untuk $t \in [0,1]$
(atau dapat pula ditulis $\bigl(\cos(\theta), \sin(\theta) \bigr)$).
Kemudian kita memaksimumkan
\begin{equation*}
\bnorm{A(x,y)} = \sqrt{{\left(t \pm \sqrt{1-t^2}\right)}^2+4t^2}
\end{equation*}
untuk memperoleh $\snorm{A} = \sqrt{3+\sqrt{5}}$.
Secara umum, kita sering menempuh dua langkah.  Sebagai contoh, perhatikan
operator
$B \in L\bigl(C([0,1],\R),\R\bigr)$
yang memetakan fungsi kontinu $f$ ke $f(0)$.  Jika $\snorm{f} = 1$ (norma seragam),
maka jelas bahwa $\babs{f(0)} \leq 1$, sehingga $\sabs{Bf} \leq 1$, yang berarti
$\snorm{B} \leq 1$.
Untuk membuktikan bahwa nilainya sama dengan $1$, perhatikan bahwa fungsi konstan $1$ memiliki
norma $1$, sehingga $B1 = 1$, yang berarti $\snorm{B} \geq 1$.
Jadi, $\snorm{B}=1$.

Norma operator tidak selalu merupakan norma pada $L(X,Y)$.
Secara khusus,
$\snorm{A}$ tidak selalu berhingga untuk $A \in L(X,Y)$.
Di bawah ini kita membuktikan bahwa $\snorm{A}$ berhingga jika $X$ berdimensi hingga.
Norma operator berhingga jika dan hanya jika $A$ kontinu.
Untuk ruang berdimensi tak berhingga, kedua pernyataan tersebut tidak selalu benar.
Sebagai contoh,
perhatikan ruang vektor yang terdiri atas fungsi-fungsi yang diferensiabel dan memiliki turunan kontinu pada
$[0,2\pi]$ dengan norma seragam.  Fungsi-fungsi
$t \mapsto \sin(n t)$ memiliki norma 1, tetapi turunannya memiliki norma $n$.  Jadi,
diferensiasi, yang merupakan operator linear bernilai di ruang fungsi-fungsi
kontinu, memiliki norma operator tak berhingga pada ruang ini.
Kita akan membatasi pembahasan pada ruang berdimensi hingga.

Untuk suatu ruang vektor berdimensi hingga $X$, kita sering membayangkan
$\R^n$, meskipun jika $X$ dilengkapi dengan suatu norma, norma tersebut belum tentu
merupakan norma Euklides baku.  Dalam latihan, Anda dapat membuktikan bahwa
setiap norma pada $\R^n$ \myquote{ekuivalen} dengan norma Euklides, dalam arti bahwa
topologi yang dihasilkannya sama.  Demi kesederhanaan, kita hanya membuktikan
proposisi berikut untuk ruang Euklides, sedangkan pembuktian untuk ruang
berdimensi hingga secara umum diserahkan sebagai latihan.

\begin{prop} \label{prop:finitedimpropnormfin}
Misalkan $X$ dan $Y$ adalah ruang vektor bernorma dengan
$X$ berdimensi hingga, dan misalkan
$A \in L(X,Y)$.
Maka $\snorm{A} < \infty$, dan
$A$ kontinu seragam (Lipschitz dengan konstanta $\snorm{A}$).
\end{prop}

\begin{proof}
Seperti telah dinyatakan, kita hanya membuktikan proposisi ini untuk ruang Euklides, jadi andaikan
bahwa $X = \R^n$ dan normanya adalah norma Euklides baku.
Kasus umum diserahkan sebagai latihan.

Misalkan $\{ e_1,e_2,\ldots,e_n \}$ adalah basis baku dari $\R^n$.
Tuliskan $x \in \R^n$, dengan $\snorm{x} = 1$, sebagai
\begin{equation*}
x = \sum_{k=1}^n c_k \, e_k .
\end{equation*}
Karena $e_k \cdot e_\ell = 0$ apabila $k \neq \ell$ dan $e_k \cdot e_k = 1$,
kita memperoleh $c_k = x \cdot e_k$.  Berdasarkan Cauchy--Schwarz,
\begin{equation*}
\sabs{c_k} = \sabs{ x \cdot e_k }
\leq \snorm{x} \, \snorm{e_k} = 1 .
\end{equation*}
Dengan demikian,
\begin{equation*}
\snorm{Ax} =
\norm{\sum_{k=1}^n c_k \, Ae_k}
\leq
\sum_{k=1}^n \sabs{c_k} \, \snorm{Ae_k} 
\leq
\sum_{k=1}^n \snorm{Ae_k} .
\end{equation*}
Ruas kanan tidak bergantung pada $x$.  Kita telah memperoleh
batas atas berhingga bagi $\snorm{Ax}$ yang tidak bergantung pada $x$, sehingga $\snorm{A} < \infty$.

Sekarang ambil ruang vektor bernorma $X$ dan $Y$, serta $A \in L(X,Y)$ dengan
$\snorm{A} < \infty$.
Untuk $v,w \in X$,
\begin{equation*}
\snorm{Av - Aw} =
\bnorm{A(v-w)} \leq \snorm{A} \, \snorm{v-w} .
\end{equation*}
Karena $\snorm{A} < \infty$, ketaksamaan di atas menyatakan bahwa
$A$ bersifat Lipschitz dengan konstanta $\snorm{A}$.
\end{proof}

\begin{prop} \label{prop:finitedimpropnorm}
\pagebreak[3]
Misalkan $X$, $Y$, dan $Z$ adalah ruang vektor bernorma berdimensi
hingga\footnote{Jika bagian \myquote{Secara khusus} dihapus dan aljabar
dengan norma operator tak berhingga ditafsirkan secara tepat, yakni dengan menetapkan bahwa $0$ dikali
$\infty$ sama dengan 0, maka hasil ini juga berlaku untuk ruang berdimensi tak berhingga.}.
\begin{enumerate}[(i)]
\item \label{item:finitedimpropnorm:i}
Jika $A,B \in L(X,Y)$ dan $c \in \R$, maka
\begin{equation*}
\snorm{A+B} \leq \snorm{A}+\snorm{B}, \qquad \snorm{cA} = \sabs{c} \, \snorm{A} .
\end{equation*}
Secara khusus, norma operator merupakan suatu norma pada ruang vektor $L(X,Y)$.
\item \label{item:finitedimpropnorm:ii}
Jika $A \in L(X,Y)$ dan $B \in L(Y,Z)$, maka
\begin{equation*}
\snorm{BA} \leq \snorm{B} \, \snorm{A} .
\end{equation*}
\end{enumerate}
\end{prop}

\begin{proof}
Pertama, karena semua ruang berdimensi hingga, semua norma
operatornya berhingga, sehingga pernyataan-pernyataan tersebut memang bermakna.

Untuk \ref{item:finitedimpropnorm:i}, ambil sebarang $x \in X$.  Maka
\begin{equation*}
\bnorm{(A+B)x} =
\snorm{Ax+Bx} \leq
\snorm{Ax}+\snorm{Bx} \leq
\snorm{A} \, \snorm{x}+\snorm{B} \,\snorm{x} =
\bigl(\snorm{A}+\snorm{B}\bigr) \snorm{x} .
\end{equation*}
Jadi, $\snorm{A+B} \leq \snorm{A}+\snorm{B}$.
Serupa dengan itu,
\begin{equation*}
\bnorm{(cA)x} =
\sabs{c} \, \snorm{Ax} \leq \bigl(\sabs{c} \,\snorm{A}\bigr) \snorm{x} .
\end{equation*}
Dengan demikian, $\snorm{cA} \leq \sabs{c} \, \snorm{A}$.  Selanjutnya,
\begin{equation*}
\sabs{c} \,  \snorm{Ax}
=
\snorm{cAx} \leq \snorm{cA} \, \snorm{x} .
\end{equation*}
Oleh sebab itu, $\sabs{c} \, \snorm{A} \leq \snorm{cA}$.


Untuk \ref{item:finitedimpropnorm:ii}, tuliskan
\begin{equation*}
\snorm{BAx} \leq \snorm{B} \, \snorm{Ax} \leq \snorm{B} \, \snorm{A} \, \snorm{x} .
\qedhere
\end{equation*}
\end{proof}

Suatu norma mendefinisikan metrik,
yang memberikan topologi ruang metrik pada $L(X,Y)$ untuk ruang vektor berdimensi hingga.
Jadi, kita dapat membahas
himpunan terbuka/tertutup, kekontinuan, kekonvergenan, dan sebagainya.

\begin{prop} \label{prop:finitedimpropinv}
\pagebreak[2]
Misalkan $X$ adalah ruang vektor bernorma berdimensi hingga.
Misalkan $GL(X) \subset L(X)$\glsadd{not:GLX} adalah himpunan operator linear
invertibel.\footnote{$GL(X)$ disebut
\emph{\myindex{grup linear umum}}, dan singkatan GL berasal dari nama
bahasa Inggrisnya.}
\begin{enumerate}[(i)]
\item \label{finitedimpropinv:i}
Jika $A \in GL(X)$, $B \in L(X)$, dan
\begin{equation} \label{eqcontineq}
\snorm{A-B} <  \frac{1}{\snorm{A^{-1}}},
\end{equation}
maka $B \in GL(X)$, yaitu $B$ invertibel.
Secara khusus, $GL(X)$ terbuka.
\item \label{finitedimpropinv:ii}
$A \mapsto A^{-1}$ merupakan fungsi
kontinu pada $GL(X)$.
\end{enumerate}
\end{prop}

Kita ilustrasikan proposisi ini dengan suatu contoh sederhana.
Perhatikan $X=\R^1$.  Dalam ruang ini, operator linear tidak lain adalah
bilangan $a$, dan norma operator dari $a$ adalah $\sabs{a}$.
Operator $a$ invertibel ($a^{-1} = \nicefrac{1}{a}$)
apabila $a \neq 0$.  Syarat $\sabs{a-b} < \frac{1}{\sabs{a^{-1}}}$
memang menyiratkan bahwa $b$ tidak nol.  Selain itu, $a \mapsto \nicefrac{1}{a}$
merupakan fungsi kontinu.
Jika dimensinya $2$ atau lebih,
terdapat operator tak invertibel selain operator nol,
dan persoalannya menjadi sedikit lebih rumit.

\begin{proof}
Mari kita buktikan \ref{finitedimpropinv:i}.   Kita memiliki informasi tentang $A^{-1}$
dan $A-B$; keduanya merupakan operator linear.
Jadi, terapkan keduanya pada suatu vektor:
\begin{equation*}
A^{-1}(A-B)x
=
x-A^{-1}Bx .
\end{equation*}
Oleh karena itu,
\begin{equation*}
\begin{split}
\snorm{x} 
& =
\bnorm{A^{-1} (A-B)x + A^{-1}Bx}
\\
& \leq
\snorm{A^{-1}}\,\snorm{A-B}\, \snorm{x} + \snorm{A^{-1}}\,\snorm{Bx} .
\end{split}
\end{equation*}
Andaikan $x \neq 0$, sehingga $\snorm{x} \neq 0$.
Dengan menggunakan \eqref{eqcontineq}, kita memperoleh
\begin{equation*}
\snorm{x} < \snorm{x} + \snorm{A^{-1}} \, \snorm{Bx} .
\end{equation*}
Jadi, $\snorm{Bx} \neq 0$ untuk setiap $x \neq 0$, dan akibatnya
$Bx \neq 0$ untuk setiap $x \neq 0$.  Dengan demikian,
$B$ injektif; jika $Bx = By$, maka $B(x-y) = 0$, sehingga $x=y$.
Karena $B$ merupakan pemetaan linear injektif dari $X$ ke $X$, yang berdimensi hingga,
pemetaan ini juga surjektif berdasarkan \propref{mv:prop:lin11onto}.
Oleh karena itu, $B$ invertibel.
Secara khusus, hal ini menunjukkan bahwa $GL(X)$ terbuka.

Mari kita buktikan \ref{finitedimpropinv:ii}.
Kita harus menunjukkan bahwa operasi pengambilan invers bersifat kontinu.
Tetapkan $A \in GL(X)$.  Ambil $B$ yang dekat dengan $A$,
tepatnya $\snorm{A-B} < \frac{1}{2 \snorm{A^{-1}}}$.
Maka \eqref{eqcontineq} dipenuhi dan $B$ invertibel.
Perhitungan yang serupa dengan perhitungan di atas (dengan menggunakan $B^{-1}y$ sebagai pengganti $x$) memberikan
\begin{equation*}
\snorm{B^{-1}y} \leq 
\snorm{A^{-1}} \, \snorm{A-B} \,  \snorm{B^{-1}y} + \snorm{A^{-1}} \, \snorm{y}
\leq
\frac{1}{2} \snorm{B^{-1}y} + \snorm{A^{-1}}\,\snorm{y} ,
\end{equation*}
atau
\begin{equation*}
\snorm{B^{-1}y} \leq 
2\snorm{A^{-1}}\,\snorm{y} .
\end{equation*}
Jadi,
$
\snorm{B^{-1}} \leq 2 \snorm{A^{-1}}
$.

Sekarang,
\begin{equation*}
A^{-1}(A-B)B^{-1} = 
A^{-1}(AB^{-1}-I) = 
B^{-1}-A^{-1} ,
\end{equation*}
dan
\begin{equation*}
\snorm{B^{-1}-A^{-1}} =
\bnorm{A^{-1}(A-B)B^{-1}} \leq
\snorm{A^{-1}}\,\snorm{A-B}\,\snorm{B^{-1}}
\leq
2\snorm{A^{-1}}^2
\snorm{A-B} .
\end{equation*}
Oleh karena itu, ketika $B$ menuju $A$, $\snorm{B^{-1}-A^{-1}}$ menuju 0, dan
dengan demikian operasi pengambilan invers merupakan fungsi kontinu di $A$.
\end{proof}

\subsection{Matriks}

Setelah menetapkan suatu basis pada ruang vektor
berdimensi hingga $X$, kita dapat merepresentasikan sebuah vektor dalam $X$ sebagai
$n$-tupel bilangan---yakni sebuah vektor di $\R^n$.
Hal yang sama dapat dilakukan terhadap $L(X,Y)$, sehingga diperoleh matriks,
yang merupakan cara praktis untuk merepresentasikan
transformasi linear antarruang berdimensi hingga.
Misalkan $\{ x_1, x_2, \ldots, x_n \}$ dan $\{ y_1, y_2, \ldots, y_m \}$
masing-masing merupakan basis bagi ruang vektor $X$ dan $Y$.  Suatu operator linear 
ditentukan oleh nilai-nilainya pada basis.  Untuk $A \in L(X,Y)$,
$A x_j$ merupakan elemen $Y$.  Definisikan bilangan-bilangan
$a_{i,j}$ melalui
\begin{equation} \label{eq:matrixrepresent}
A x_j = \sum_{i=1}^m a_{i,j} \, y_i ,
\end{equation}
dan susun bilangan-bilangan tersebut sebagai sebuah \emph{\myindex{matriks}},
yang, dengan sedikit penyalahgunaan notasi, juga kita sebut $A$,
\glsadd{not:matrix}
\begin{equation*}
A =
\begin{bmatrix}
a_{1,1} & a_{1,2} & \cdots & a_{1,n} \\
a_{2,1} & a_{2,2} & \cdots & a_{2,n} \\
\vdots & \vdots & \ddots & \vdots \\
a_{m,1} & a_{m,2} & \cdots & a_{m,n}
\end{bmatrix} .
\end{equation*}
Kadang-kadang kita menulis $A$ sebagai $[a_{i,j}]$.
Kita menyebut $A$ matriks berukuran $m$-kali-$n$.
Pada matriks tersebut, \emph{\myindex{kolom}} ke-$j$ memuat tepat koefisien-koefisien
yang merepresentasikan $A x_j$ terhadap basis $\{ y_1,y_2,\ldots,y_m \}$.
Dari bilangan-bilangan $a_{i,j}$, melalui rumus
\eqref{eq:matrixrepresent}, kita memperoleh operator linear yang bersesuaian,
karena operator itu ditentukan oleh aksinya pada suatu basis.  Jadi, setelah menetapkan
basis pada $X$ dan $Y$, kita memperoleh korespondensi satu-ke-satu antara
$L(X,Y)$ dan matriks-matriks berukuran $m$-kali-$n$.
Jika
\begin{equation*}
z = \sum_{j=1}^n z_j \, x_j ,
\end{equation*}
maka
\begin{equation*}
A z =
\sum_{j=1}^n z_j \, A x_j 
=
\sum_{j=1}^n z_j \left( \sum_{i=1}^m  a_{i,j}\, y_i \right) 
=
\sum_{i=1}^m \left(\sum_{j=1}^n  a_{i,j}\, z_j \right) y_i ,
\end{equation*}
yang menghasilkan aturan perkalian matriks yang lazim,
dengan memandang $z$ sebagai vektor kolom, yaitu matriks berukuran $n$-kali-1.
%
Secara lebih umum, jika $B$ adalah matriks berukuran $n$-kali-$r$ dengan entri $b_{j,k}$, maka 
matriks $C = AB$ berukuran $m$-kali-$r$. Entri $(i,k)$ pada matriks ini,
yaitu $c_{i,k}$, adalah
\begin{equation*}
c_{i,k} =
\sum_{j=1}^n a_{i,j}\,b_{j,k} .
\end{equation*}
Salah satu cara mengingatnya ialah dengan mengurutkan indeks
seperti di atas---\emph{baris, kolom}---dan
menempatkan elemen-elemen dalam urutan yang sama dengan matriksnya;
dengan demikian, \myquote{indeks tengah} \myquote{dihilangkan melalui penjumlahan.}

Terdapat korespondensi satu-ke-satu antara matriks dan operator linear dalam
$L(X,Y)$, setelah kita menetapkan basis pada $X$ dan $Y$.  Jika kita 
memilih basis yang berbeda, kita memperoleh matriks yang berbeda.
Pembedaan ini penting.
Operator $A$ bekerja pada elemen-elemen $X$, sedangkan matriks
bekerja dengan $n$-tupel bilangan, yaitu vektor-vektor di
$\R^n$.  Sesuai konvensi, kita menggunakan basis standar di $\R^n$ kecuali dinyatakan
lain, dan kita mengidentifikasi $L(\R^n,\R^m)$ dengan himpunan matriks berukuran $m$-kali-$n$.

Pemetaan linear yang mengubah satu basis menjadi basis lain direpresentasikan oleh
matriks persegi yang kolom-kolomnya merepresentasikan vektor-vektor
basis kedua terhadap basis pertama.  Pemetaan linear semacam itu kita
sebut \emph{\myindex{perubahan basis}}.  Jadi, untuk dua pilihan
basis dalam ruang vektor berdimensi $n$, terdapat suatu pemetaan linear
(suatu perubahan basis)
yang memetakan satu basis ke basis lainnya, dan pemetaan ini bersesuaian dengan matriks
berukuran $n$-kali-$n$ yang melakukan operasi terkait pada $\R^n$.

\medskip

Misalkan
$X=\R^n$, $Y=\R^m$, dan
semua basis yang digunakan adalah basis standar.
Dengan menggunakan ketaksamaan Cauchy--Schwarz untuk $c=(c_1,c_2,\ldots,c_n) \in \R^n$, kita peroleh
\begin{equation*}
\snorm{Ac}^2
=
\sum_{i=1}^m { \left(\sum_{j=1}^n a_{i,j} \, c_j \right)}^2
\leq
\sum_{i=1}^m
\left(
\left(\sum_{j=1}^n {(a_{i,j})}^2 \right)
\left(\sum_{j=1}^n {(c_j)}^2 \right)
\right)
=
\left(
\sum_{i=1}^m \sum_{j=1}^n {(a_{i,j})}^2 \right)
\snorm{c}^2 .
\end{equation*}
Dengan kata lain, kita memperoleh batas bagi norma operator (perhatikan bahwa kesamaan
jarang terjadi)
\begin{equation*}
\snorm{A} \leq
\sqrt{\sum_{i=1}^m \sum_{j=1}^n {(a_{i,j})}^2} .
\end{equation*}
Ruas kanan merupakan norma Euklides pada $\R^{nm}$, yang
diperoleh dengan memandang semua entri matriks sebagai koordinat.
Jika entri-entri tersebut menuju nol, maka $\snorm{A}$ menuju nol.
Sebaliknya,
\begin{equation*}
\sum_{i=1}^m \sum_{j=1}^n {(a_{i,j})}^2
=
\sum_{j=1}^n \snorm{A e_j}^2
\leq
\sum_{j=1}^n \snorm{A}^2
= n \snorm{A}^2 .
\end{equation*}
Jadi, jika norma operator $A$ menuju nol, entri-entrinya juga menuju nol.
Secara khusus, jika $A$ tetap dan $B$ berubah,
maka entri-entri $B$ menuju entri-entri $A$ jika dan hanya jika $B$ menuju $A$ dalam
norma operator ($\snorm{A-B}$ menuju nol).  Kita telah membuktikan:

\begin{prop} \label{prop:matrixcont}
Topologi (himpunan semua himpunan terbuka) pada $L(\R^n,\R^m)$ tetap sama,
baik ketika $L(\R^n,\R^m)$ dipandang sebagai ruang metrik dengan norma
operator maupun dengan metrik Euklides pada $\R^{nm}$.

Secara khusus, misalkan $S$ merupakan ruang metrik dan
$\pi \colon L(\R^n,\R^m) \to \R^{nm}$ mengidentifikasi
suatu operator dengan $nm$-tupel yang terdiri atas entri-entri matriks bersesuaian.
Maka
$f \colon S \to L(\R^n,\R^m)$
kontinu
jika dan hanya jika
$\pi \circ f \colon S \to \R^{nm}$
kontinu.
Demikian pula halnya dengan
$g \colon L(\R^n,\R^m) \to S$ dan
$g \circ \pi^{-1} \colon \R^{nm} \to S$.
\end{prop}

\subsection{Determinan}

Suatu bilangan tertentu dapat dikaitkan dengan matriks persegi untuk mengukur
bagaimana pemetaan linear yang bersesuaian meregangkan ruang.  Secara khusus,
bilangan yang disebut determinan ini dapat digunakan untuk menguji invertibilitas matriks.

Definisikan simbol
$\operatorname{sgn}(x)$\glsadd{not:sgn} (dibaca \myquote{tanda dari $x$}) untuk suatu bilangan $x$ dengan
\begin{equation*}
\operatorname{sgn}(x)
\coloneqq
\begin{cases}
-1 & \text{jika } x < 0 , \\
0  & \text{jika } x = 0 , \\
1  & \text{jika } x > 0 .
\end{cases}
\end{equation*}
Suatu \emph{\myindex{permutasi}}
$\sigma = (\sigma_1,\sigma_2,\ldots,\sigma_n)$ adalah
penyusunan ulang dari $(1,2,\ldots,n)$.
Definisikan
\begin{equation} \label{eq:sgndef}
\operatorname{sgn}(\sigma) = \operatorname{sgn}(\sigma_1,\ldots,\sigma_n)
\coloneqq
\prod_{p < q} \operatorname{sgn}(\sigma_q-\sigma_p) .
\avoidbreak
\end{equation}
Di sini $\prod$ menyatakan perkalian, sebagaimana $\sum$ menyatakan
penjumlahan.\glsadd{not:prod}

Setiap permutasi dapat diperoleh melalui
suatu barisan transposisi (pertukaran dua unsur).
Suatu permutasi disebut
\emph{genap}\index{permutasi genap}
(atau\ \emph{ganjil})\index{permutasi ganjil}
jika diperlukan sejumlah genap (atau\ ganjil)
transposisi untuk mengubah $(1,2,\ldots,n)$ menjadi $\sigma$.
Sebagai contoh, $(2,4,3,1)$ berjarak dua transposisi dari
$(1,2,3,4)$ dan karena itu genap: $(1,2,3,4) \to (2,1,3,4) \to (2,4,3,1)$.
Sifat genap atau ganjil ini terdefinisi dengan baik:
$\operatorname{sgn}(\sigma)$
bernilai $1$ jika $\sigma$ genap dan $-1$ jika $\sigma$ ganjil (latihan).
Fakta ini mengikuti karena penerapan
suatu transposisi mengubah tanda, sedangkan
$\operatorname{sgn}(1,2,\ldots,n) = 1$.

Misalkan $S_n$ adalah himpunan semua permutasi pada $n$ unsur (yakni
\emph{\myindex{grup simetris}}).
Misalkan $A= [a_{i,j}]$ suatu matriks persegi berukuran $n$-kali-$n$.
Definisikan \emph{\myindex{determinan}} dari $A$ sebagai
\glsadd{not:det}
\begin{equation*}
\det(A) \coloneqq
\sum_{\sigma \in S_n}
\operatorname{sgn} (\sigma) \prod_{i=1}^n a_{i,\sigma_i} .
\end{equation*}

\begin{prop}
\pagebreak[0]
\leavevmode
\begin{enumerate}[(i)]
\item \label{prop:det:i} $\det(I) = 1$.
\item \label{prop:det:ii} Untuk setiap $j=1,2,\ldots,n$, fungsi $x_j
\mapsto \det\bigl([x_1 ~~~ x_2 ~~~ \cdots ~~~ x_n ]\bigr)$
bersifat linear.
\item \label{prop:det:iii} Jika dua kolom suatu matriks dipertukarkan, tanda determinannya
berubah.
\item \label{prop:det:iv} Jika dua kolom $A$ sama, maka $\det(A) = 0$.
\item \label{prop:det:v} Jika suatu kolom bernilai nol, maka $\det(A) = 0$.
\item \label{prop:det:vi} $A \mapsto \det(A)$ merupakan fungsi kontinu pada
$L(\R^n)$.
\item \label{prop:det:vii} $\det\left( \left[\begin{smallmatrix} a & b \\ c
&d \end{smallmatrix}\right] \right)
= ad-bc$, dan $\det \bigl( [a] \bigr) = a$.
\end{enumerate}
\end{prop}

Bahkan, determinan merupakan satu-satunya fungsi yang memenuhi
\ref{prop:det:i},
\ref{prop:det:ii}, dan
\ref{prop:det:iii},
tetapi hal itu menyimpang dari pembahasan kita.
Maksud \ref{prop:det:ii} ialah: jika semua vektor
$x_1,\ldots,x_n$ kecuali $x_j$ ditetapkan, lalu
$v,w \in \R^n$ adalah dua vektor,
dan $a,b \in \R$ adalah skalar, maka
\begin{multline*}
\det\bigl([x_1 ~~~ \cdots ~~~ x_{j-1} ~~~ (av+bw) ~~~ x_{j+1} ~~~
\cdots ~~~ x_n]\bigr) =
\\
a \det\bigl([x_1 ~~~ \cdots ~~~ x_{j-1} ~~~ v ~~~ x_{j+1} ~~~
\cdots ~~~ x_n]\bigr)
+
b
\det\bigl([x_1 ~~~ \cdots ~~~ x_{j-1} ~~~ w ~~~ x_{j+1} ~~~ \cdots
~~~ x_n]\bigr) .
\end{multline*}

\begin{proof}
Kita membahas pembuktiannya dengan singkat karena kemungkinan besar Anda pernah melihatnya.
%
Butir \ref{prop:det:i} bersifat langsung.  Untuk \ref{prop:det:ii}, perhatikan bahwa setiap suku dalam definisi
determinan memuat tepat satu faktor dari setiap kolom.
%
Butir \ref{prop:det:iii} mengikuti karena menukar dua kolom berarti menukar
dua bilangan yang bersesuaian dalam setiap unsur $S_n$.  Akibatnya, semua tanda
berubah.
Butir \ref{prop:det:iv} mengikuti karena jika dua kolom sama dan kita
menukarnya, kita memperoleh kembali
matriks yang sama.  Jadi butir \ref{prop:det:iii} menyatakan bahwa determinannya harus 0.
%
Butir \ref{prop:det:v} mengikuti karena hasil kali pada setiap suku dalam definisi memuat
satu unsur dari kolom nol.
Butir \ref{prop:det:vi} mengikuti karena $\det$ merupakan polinom dalam entri-entri matriks
dan karenanya kontinu (sebagai fungsi dari entri-entri matriks).
Suatu fungsi pada ruang matriks bersifat kontinu terhadap norma operator jika dan hanya jika
fungsi tersebut kontinu sebagai fungsi dari entri-entri matriks
(\propref{prop:matrixcont}).
Terakhir, butir \ref{prop:det:vii} diperoleh melalui perhitungan langsung.
\end{proof}

Determinan memberi tahu kita tentang luas dan volume serta perubahannya.
Sebagai contoh, dalam kasus $1$-kali-$1$, suatu matriks hanyalah sebuah bilangan dan
determinannya tepat sama dengan bilangan tersebut.  Determinan menyatakan bagaimana pemetaan linear
\myquote{meregangkan} ruang.  Demikian pula,
misalkan $A \in L(\R^2)$ suatu transformasi linear.
Dapat diperiksa secara langsung bahwa
luas citra persegi satuan $A\bigl([0,1]^2\bigr)$ adalah
$\babs{\det(A)}$; lihat \figureref{fig:imagesquare} untuk sebuah contoh.
Hal ini berlaku untuk bangun sebarang, bukan hanya persegi
satuan:
Nilai mutlak determinan menyatakan faktor peregangan luas.
Tanda determinan menyatakan apakah citranya
terbalik (mengubah orientasi) atau tidak.
Dalam $\R^3$, determinan
menyatakan perubahan volume tiga dimensi, dan dalam $n$ dimensi menyatakan perubahan
volume $n$ dimensi.  Pernyataan ini kita berikan tanpa pembuktian.
\begin{myfigureht}
\myincludepdft{linalg-imagesquare}{%
Terdapat dua diagram bidang.  Pada diagram kiri,
terdapat persegi dengan titik sudut (0,0), (1,0), (1,1), dan (0,1).
Pada diagram kanan, terdapat persegi dengan titik sudut yang bersesuaian
(0,0), (1, minus satu), (2,0), dan (1,1), dengan korespondensi yang ditunjukkan
oleh anak panah.}
\caption{Citra persegi satuan ${[0,1]}^2$ melalui matriks
$\left[\begin{smallmatrix}1 & 1 \\ -1 & 1\end{smallmatrix}\right]$.
Citranya merupakan persegi dengan panjang sisi $\sqrt{2}$, sehingga luasnya 2, dan determinan
matriks tersebut adalah 2.\label{fig:imagesquare}}
\end{myfigureht}

\begin{prop}
Jika $A$ dan $B$ matriks berukuran $n$-kali-$n$, maka $\det(AB) = \det(A)\det(B)$.
Lebih lanjut, $A$ invertibel jika dan hanya jika $\det(A) \neq 0$, dan dalam
kasus ini, $\det(A^{-1}) = \frac{1}{\det(A)}$.
\end{prop}

\begin{proof}
Misalkan $b_1,b_2,\ldots,b_n$ kolom-kolom dari $B$.  Maka
\begin{equation*}
AB = [ Ab_1 \quad Ab_2 \quad  \cdots \quad  Ab_n ] .
\end{equation*}
Artinya, kolom-kolom $AB$ adalah
$Ab_1,Ab_2,\ldots,Ab_n$.

Misalkan $b_{j,k}$ menyatakan unsur-unsur $B$ dan
$a_j$ menyatakan kolom-kolom $A$.
Berdasarkan linearitas determinan,
\begin{equation*}
\begin{split}
\det(AB) & =
\det \bigl([ Ab_1 \quad Ab_2 \quad  \cdots \quad  Ab_n ] \bigr) =
\det \left(\left[ \sum_{j=1}^n b_{j,1} a_j \quad Ab_2 \quad  \cdots \quad  Ab_n \right]\right) \\
& =
\sum_{j=1}^n
b_{j,1}
\det \bigl([ a_j \quad Ab_2 \quad  \cdots \quad  Ab_n ]\bigr) \\
& =
\sum_{1 \leq j_1,j_2,\ldots,j_n \leq n}
b_{j_1,1}
b_{j_2,2}
\cdots
b_{j_n,n}
\det \bigl([ a_{j_1} \quad a_{j_2} \quad  \cdots \quad  a_{j_n} ]\bigr) \\
& =
\left(
\sum_{(j_1,j_2,\ldots,j_n) \in S_n}
b_{j_1,1}
b_{j_2,2}
\cdots
b_{j_n,n}
\operatorname{sgn}(j_1,j_2,\ldots,j_n)
\right)
\det \bigl([ a_{1} \quad a_{2} \quad  \cdots \quad  a_{n} ]\bigr) .
\end{split}
\end{equation*}
Pada kesamaan terakhir, kita menjumlahkan atas unsur-unsur $S_n$
alih-alih semua tupel-$n$ bilangan bulat antara 1 dan $n$,
karena
jika dua kolom dalam determinan sama, maka
determinannya nol.  Menyusun ulang kolom-kolom ke
urutan semula menghasilkan nilai sgn.

Kesimpulan bahwa $\det(AB) = \det(A)\det(B)$
mengikuti dengan mengenali bahwa ekspresi dalam tanda kurung di atas adalah determinan $B$.
Hal ini kita peroleh dengan menyubstitusikan $A=I$.
Ekspresi yang diperoleh untuk determinan $B$ memiliki baris dan kolom
yang tertukar; sebagai hasil tambahan, kita juga baru saja membuktikan bahwa determinan
suatu matriks sama dengan determinan transposnya.

Mari kita buktikan bagian
\myquote{Lebih lanjut.}
Jika $A$ invertibel,
maka $A^{-1}A = I$.
Akibatnya, $\det(A^{-1})\det(A) = \det(A^{-1}A) = \det(I) = 1$.
Jika $A$ tidak invertibel, maka matriks tersebut tidak injektif, sehingga $A$ memetakan
suatu vektor tak nol ke nol.
Dengan kata lain, kolom-kolom $A$ bergantung linear.
Misalkan
\begin{equation*}
\sum_{k=1}^n \gamma_k\, a_k = 0 ,
\end{equation*}
dengan tidak semua $\gamma_k$ sama dengan 0.
Tanpa mengurangi keumuman, misalkan $\gamma_1\neq 0$.
Ambil
\begin{equation*}
B \coloneqq
\begin{bmatrix}
\gamma_1 & 0 & 0 & \cdots & 0 \\
\gamma_2 & 1 & 0 & \cdots & 0 \\
\gamma_3 & 0 & 1 & \cdots & 0 \\
\vdots & \vdots & \vdots & \ddots & \vdots \\
\gamma_n & 0 & 0 & \cdots & 1
\end{bmatrix} .
\end{equation*}
Dengan menggunakan definisi determinan (hanya ada satu
permutasi $\sigma$ yang membuat $\prod_{i=1}^n b_{i,\sigma_i}$ tak nol),
kita memperoleh $\det(B) = \gamma_1 \neq 0$.
Kemudian
$\det(AB) = \det(A)\det(B) = \gamma_1\det(A)$.
Kolom pertama $AB$ bernilai nol, sehingga $\det(AB) = 0$.  Kita menyimpulkan
$\det(A) = 0$.
\end{proof}

\begin{prop}
Determinan tidak bergantung pada basis: Jika $A$ dan $B$ adalah
matriks berukuran $n$-kali-$n$ dan $B$
invertibel, maka
\begin{equation*}
\det(A) = \det(B^{-1}AB) .
\end{equation*}
\end{prop}

\begin{proof}
$\det(B^{-1}AB) =
\det(B^{-1})\det(A)\det(B) =
\frac{1}{\det(B)}\det(A)\det(B) = \det(A)$.
\end{proof}

Jika pada suatu basis $A$ adalah matriks yang merepresentasikan suatu
operator linear, maka untuk basis lain kita dapat menemukan matriks $B$ sedemikian
sehingga matriks $B^{-1}AB$ membawa kita ke basis pertama, menerapkan $A$ dalam
basis pertama, lalu membawa kita kembali ke basis awal.
Misalkan $X$ ruang vektor berdimensi hingga.
Misalkan $\Phi \in L(X,\R^n)$ memetakan basis $\{ x_1,\ldots,x_n \}$ ke
basis standar $\{ e_1,\ldots,e_n \}$ dan $\Psi \in L(X,\R^n)$
memetakan basis lain $\{ y_1, \ldots, y_n \}$ ke basis standar.
Misalkan $T \in L(X)$ operator linear dan misalkan matriks $A$
merepresentasikan operator tersebut
dalam basis $\{ x_1,\ldots,x_n \}$.
Maka $B$ dipilih sedemikian sehingga diperoleh diagram berikut\footnote{%
Gambar semacam ini disebut \emph{\myindex{diagram komutatif}}.
Dalam diagram komutatif,
mengikuti anak panah melalui jalur mana pun harus memberikan hasil yang sama.}:
%mbxSTARTIGNORE
\begin{equation*}
\begin{tikzcd}[row sep=normal,column sep=large]
\R^n \ar[r, "B^{-1}AB"]  \ar[dd, bend right=50, "B"'] &
  \R^n \ar[d, "\Psi^{-1}"'] \\
X \ar[r, "T"] \ar[u, "\Psi"'] \ar[d, "\Phi"]        &
  X \\
\R^n \ar[r, "A"]                                      &
  \R^n \ar[u, "\Phi^{-1}"] \ar[uu, bend right=50, "B^{-1}"']
\end{tikzcd}
\end{equation*}
%mbxENDIGNORE
%mbxlatex \begin{center}
%mbxlatex \myincludepdft{change_of_basis_cd}{Diagram himpunan
%mbxlatex dan pemetaan di antaranya yang direpresentasikan oleh anak panah.
%mbxlatex Baris pertama memuat dua himpunan yang keduanya merupakan ruang Euklides n dimensi, R pangkat n.
%mbxlatex Baris kedua memuat dua himpunan yang keduanya adalah X.
%mbxlatex Baris ketiga memuat dua himpunan yang keduanya merupakan ruang Euklides n dimensi, R pangkat n.
%mbxlatex R pangkat n kiri atas memiliki anak panah berlabel B invers A B
%mbxlatex menuju R pangkat n kanan atas.
%mbxlatex R pangkat n kiri atas memiliki anak panah berlabel B menuju R pangkat n kiri bawah.
%mbxlatex R pangkat n kanan bawah memiliki anak panah berlabel B invers
%mbxlatex menuju R pangkat n kanan atas.
%mbxlatex X kiri memiliki anak panah berlabel Psi
%mbxlatex menuju R pangkat n kiri atas,
%mbxlatex anak panah berlabel Phi menuju R pangkat n kiri bawah,
%mbxlatex dan anak panah berlabel T menuju X kanan.
%mbxlatex R pangkat n kanan atas memiliki anak panah berlabel
%mbxlatex Psi invers menuju X kanan, dan
%mbxlatex R pangkat n kanan bawah memiliki anak panah
%mbxlatex berlabel Phi invers menuju X kanan.
%mbxlatex R pangkat n kiri bawah memiliki anak panah berlabel A
%mbxlatex menuju R pangkat n kanan bawah.}
%mbxlatex \end{center}
Kedua $\R^n$ pada baris bawah merepresentasikan
$X$ dalam basis pertama, sedangkan kedua $\R^n$ pada baris atas merepresentasikan $X$ dalam
basis
kedua.

Jika kita menghitung determinan matriks $A$, kita memperoleh
determinan yang sama ketika menggunakan basis lain;
dalam basis lain matriksnya adalah $B^{-1}AB$.
Akibatnya,
\begin{equation*}
\det \colon L(X) \to \R
\avoidbreak
\end{equation*}
merupakan fungsi yang terdefinisi dengan baik tanpa perlu menetapkan suatu basis.
Artinya, $\det$ didefinisikan pada $L(X)$, bukan hanya pada matriks.

\medskip

Terdapat tiga jenis matriks yang kita sebut
\emph{matriks elementer}\index{matriks elementer}.
Misalkan
$e_1,e_2,\ldots,e_n$ adalah basis standar pada $\R^n$ seperti biasa.
Pertama,
untuk $j =
1,2,\ldots,n$ dan
$\lambda \in \R$, $\lambda \neq 0$, definisikan
jenis pertama matriks elementer,
yaitu matriks
berukuran $n$-kali-$n$, $E$, dengan
\begin{equation*}
Ee_i \coloneqq
\begin{cases}
e_i & \text{jika } i \neq j , \\
\lambda e_i & \text{jika } i = j .
\end{cases}
\end{equation*}
Untuk sebarang matriks berukuran $n$-kali-$m$ yang dinyatakan dengan $M$, matriks $EM$ sama dengan matriks $M$
kecuali baris ke-$j$ dikalikan dengan $\lambda$.
Melalui perhitungan mudah (latihan), diperoleh $\det(E) = \lambda$.

Selanjutnya, untuk $j,k$ dengan $j\neq k$ dan $\lambda \in \R$,
definisikan jenis kedua matriks elementer $E$ dengan
\begin{equation*}
Ee_i \coloneqq
\begin{cases}
e_i               & \text{jika } i \neq j , \\
e_i + \lambda e_k & \text{jika } i = j .
\end{cases}
\end{equation*}
Untuk sebarang matriks berukuran $n$-kali-$m$ yang dinyatakan dengan $M$, matriks $EM$ sama dengan matriks $M$
kecuali $\lambda$ kali baris ke-$j$ ditambahkan pada baris ke-$k$.
Melalui perhitungan mudah (latihan), diperoleh $\det(E) = 1$.

Terakhir, untuk $j$ dan $k$ dengan $j\neq k$, definisikan
jenis ketiga matriks elementer $E$ dengan
\begin{equation*}
Ee_i \coloneqq
\begin{cases}
e_i & \text{jika } i \neq j \text{ dan } i \neq k , \\
e_k & \text{jika } i = j , \\
e_j & \text{jika } i = k .
\end{cases}
\end{equation*}
Untuk sebarang matriks berukuran $n$-kali-$m$ yang dinyatakan dengan $M$, matriks $EM$ adalah matriks yang sama dengan
baris ke-$j$ dan ke-$k$ dipertukarkan.
Melalui perhitungan mudah (latihan), diperoleh $\det(E) = -1$.

\begin{prop} \label{prop:elemmatrixdecomp}
Misalkan $T$ matriks invertibel berukuran $n$-kali-$n$.  Maka terdapat suatu
barisan berhingga matriks elementer $E_1, E_2, \ldots, E_k$ sedemikian sehingga
\begin{equation*}
T = E_1 E_2 \cdots E_k ,
\end{equation*}
dan
\begin{equation*}
\det(T) = \det(E_1)\det(E_2)\cdots \det(E_k) .
\end{equation*}
\end{prop}

Pembuktiannya diserahkan sebagai latihan.
Proposisi tersebut menyatakan bahwa determinan dapat dihitung melalui operasi baris elementer.
Kita tidak perlu memfaktorkan matriks sepenuhnya menjadi hasil kali
matriks-matriks elementer.  Cukup lakukan operasi baris
sampai diperoleh \emph{\myindex{matriks segitiga atas}}, yaitu matriks
$[a_{i,j}]$ dengan $a_{i,j} = 0$ jika $i > j$.  Menghitung determinan matriks
semacam ini tidak sulit (latihan).

Faktorisasi menjadi matriks-matriks elementer (atau variasi matriks elementer)
berguna dalam pembuktian yang melibatkan operator linear sebarang, dengan
mereduksi pembuktian menjadi kasus matriks elementer, serupa dengan perhitungan
determinan.


\subsection{Latihan}

\begin{exercise}
Untuk ruang vektor $X$ dengan norma $\snorm{\cdot}$, tunjukkan bahwa
$d(x,y) \coloneqq \snorm{x-y}$ menjadikan $X$ ruang metrik.
\end{exercise}

\begin{exercise}[Mudah]
Tunjukkan bahwa untuk matriks persegi $A$ dan $B$ yang berukuran sama,
$\det(AB) = \det(BA)$.
\end{exercise}

\begin{exercise}
Untuk $x \in \R^n$, definisikan
\begin{equation*}
\snorm{x}_{\infty} \coloneqq
\max \bigl\{ \sabs{x_1}, \sabs{x_2}, \ldots, \sabs{x_n} \bigr\} ,
\end{equation*}
yang kadang-kadang disebut norma supremum atau norma maksimum.
\begin{enumerate}[a)]
\item
Tunjukkan bahwa $\snorm{\cdot}_\infty$ merupakan norma pada $\R^n$ (yang mendefinisikan
jarak berbeda).
\item
Seperti apa bola satuan $B(0,1)$ menurut norma ini?
\end{enumerate}
\end{exercise}

\begin{exercise}
Untuk $x \in \R^n$, definisikan
\begin{equation*}
\snorm{x}_{1} \coloneqq \sum_{k=1}^n \sabs{x_k},
\end{equation*}
yang kadang-kadang disebut norma-$1$ (atau norma $L^1$).
\begin{enumerate}[a)]
\item
Tunjukkan bahwa $\snorm{\cdot}_1$ merupakan norma pada $\R^n$ (yang mendefinisikan
jarak berbeda, kadang-kadang disebut jarak taksi).
\item
Seperti apa bola satuan $B(0,1)$ menurut norma ini?
Pikirkan bentuknya dalam ${\mathbb{R}}^2$ dan $\mathbb{R}^3$.
Petunjuk: Bola itu merupakan interior selubung konveks dari sejumlah berhingga titik.
\end{enumerate}
\end{exercise}

\begin{exercise}
\pagebreak[3]
Dengan menggunakan norma Euklides pada $\R^2$, hitung norma operator dari
operator-operator dalam $L(\R^2)$ yang diberikan oleh matriks berikut:

\medskip
\noindent
a)
$\left[
\begin{smallmatrix}
1 & 0 \\
0 & 2
\end{smallmatrix}
\right]$
\quad
b)
$\left[
\begin{smallmatrix}
0 & 1 \\
-1 & 0
\end{smallmatrix}
\right]$
\quad
c)
$\left[
\begin{smallmatrix}
1 & 1 \\
0 & 1
\end{smallmatrix}
\right]$
\quad
d)
$\left[
\begin{smallmatrix}
0 & 1 \\
0 & 0
\end{smallmatrix}
\right]$
\end{exercise}


\begin{exercise} \label{exercise:normonedim}
\pagebreak[2]
Dengan menggunakan norma Euklides standar pada $\R^n$, tunjukkan:
\begin{enumerate}[a)]
\item
Misalkan $A \in L(\R,\R^n)$ didefinisikan untuk $x \in \R$ oleh $Ax \coloneqq xa$
dengan suatu vektor $a \in \R^n$.
Maka norma operator $\snorm{A}_{L(\R,\R^n)} = \snorm{a}_{\R^n}$.
(Artinya, norma operator $A$ adalah norma Euklides dari $a$.)
\item
Misalkan $B \in L(\R^n,\R)$ didefinisikan untuk $x \in \R^n$ oleh $Bx \coloneqq b \cdot x$
dengan suatu vektor $b \in \R^n$.
Maka norma operator $\snorm{B}_{L(\R^n,\R)} = \snorm{b}_{\R^n}$.
\end{enumerate}
\end{exercise}

\begin{exercise}
Misalkan $\sigma = (\sigma_1,\sigma_2,\ldots,\sigma_n)$ merupakan permutasi dari
$(1,2,\ldots,n)$.
\begin{enumerate}[a)]
\item
Tunjukkan bahwa kita dapat melakukan sejumlah berhingga transposisi (pertukaran dua
elemen) untuk memperoleh $(1,2,\ldots,n)$.
\item
Dengan menggunakan definisi \eqref{eq:sgndef},
tunjukkan bahwa $\sigma$ genap jika $\operatorname{sgn}(\sigma) = 1$ dan $\sigma$
ganjil jika $\operatorname{sgn}(\sigma) = -1$.  Secara khusus, tunjukkan bahwa
sifat ganjil atau genap terdefinisi dengan baik.
\end{enumerate}
\end{exercise}

\begin{exercise}
Verifikasi perhitungan determinan untuk ketiga jenis 
matriks elementer.
\end{exercise}

\begin{exercise}
Buktikan \propref{prop:elemmatrixdecomp}.
\end{exercise}

\begin{exercise}
\leavevmode
\begin{enumerate}[a)]
\item
Misalkan $D = [d_{i,j}]$ merupakan \emph{\myindex{matriks diagonal}} berukuran $n$-kali-$n$, yakni $d_{i,j} = 0$ jika $i
\neq j$.  Tunjukkan bahwa $\det(D) = d_{1,1}d_{2,2} \cdots d_{n,n}$.
\item
Misalkan $A$ merupakan matriks yang dapat didiagonalisasi.  Artinya, terdapat matriks
$B$ sedemikian sehingga $B^{-1}AB = D$ untuk suatu matriks diagonal $D = [d_{i,j}]$.  Tunjukkan
bahwa $\det(A) = d_{1,1}d_{2,2} \cdots d_{n,n}$.
\end{enumerate}
\end{exercise}

\begin{exercise}
\pagebreak[1]
Ambil ruang vektor polinom $\R[t]$ dan misalkan $D \in L(\R[t])$ adalah
operator turunan (dalam latihan sebelumnya telah kita buktikan bahwa $D$ merupakan operator
linear).  Untuk $P(t) = c_0 + c_1 t + \cdots + c_n
t^n \in \R[t]$, definisikan $\snorm{P} \coloneqq \sup \bigl\{ \sabs{c_j} : j = 0,1,2,\ldots,n \bigr\}$.
\begin{enumerate}[a)]
\item
Tunjukkan bahwa $\snorm{\cdot}$ merupakan norma pada $\R[t]$.
\item
Buktikan $\snorm{D} = \infty$.  Petunjuk: Tinjau polinom $t^n$ ketika $n$ menuju tak berhingga.
\end{enumerate}
\end{exercise}

\begin{exercise}
\pagebreak[1]
Kita selesaikan pembuktian \propref{prop:finitedimpropnormfin}.
Misalkan $X$ ruang vektor bernorma berdimensi hingga dengan basis
$\{ x_1,x_2,\ldots,x_n \}$, dan misalkan $Y$ suatu ruang vektor bernorma. Gunakan $\snorm{\cdot}_X$ untuk norma pada $X$,
$\snorm{\cdot}_{\R^n}$ untuk norma Euklides standar pada $\R^n$, dan
$\snorm{\cdot}_{L(X,Y)}$ untuk norma operator.
\begin{enumerate}[a)]
\item
Definisikan $f \colon \R^n \to \R$,
\begin{equation*}
f(c_1,c_2,\ldots,c_n) \coloneqq 
\snorm{c_1 x_1 + c_2 x_2 + \cdots + c_n x_n}_X .
\end{equation*}
Tunjukkan bahwa $f$ 
kontinu.
\item
Tunjukkan bahwa terdapat bilangan positif $m$ dan $M$ sedemikian
sehingga jika $c = (c_1,c_2,\ldots,c_n) \in \R^n$ dengan
$\snorm{c}_{\R^n} = 1$, maka 
$m \leq \snorm{c_1 x_1 + c_2 x_2 + \cdots + c_n x_n}_X \leq M$.
\item
Tunjukkan bahwa terdapat bilangan $B$ sedemikian sehingga jika
$\snorm{c_1 x_1 + c_2 x_2 + \cdots + c_n x_n}_X = 1$,
maka $\sabs{c_j} \leq B$.
\item
Gunakan bagian c) untuk menunjukkan bahwa jika $X$ merupakan ruang vektor 
berdimensi hingga dan $A \in L(X,Y)$, maka $\snorm{A}_{L(X,Y)} < \infty$.
\end{enumerate}
\end{exercise}

\begin{exercise} \label{exercise:allnormsequiv}
\pagebreak[3]
Misalkan $X$ ruang vektor berdimensi hingga dengan basis
$\{ x_1,x_2,\ldots,x_n \}$.
\begin{enumerate}[a)]
\item
Misalkan $\snorm{\cdot}_X$ merupakan suatu norma pada $X$,
$c = (c_1,c_2,\ldots,c_n) \in \R^n$,
dan $\snorm{\cdot}_{\R^n}$ adalah
norma Euklides standar pada $\R^n$.
Buktikan bahwa terdapat bilangan $m,M > 0$ sedemikian sehingga
untuk setiap $c \in \R^n$,
\begin{equation*}
m \snorm{c}_{\R^n}
\leq
\snorm{c_1 x_1 + c_2 x_2 + \cdots + c_n x_n}_X
\leq
M \snorm{c}_{\R^n} .
\end{equation*}
Petunjuk: Lihat latihan sebelumnya.
\item
Gunakan bagian a) untuk menunjukkan bahwa jika
$\snorm{\cdot}_1$ dan
$\snorm{\cdot}_2$ merupakan dua norma pada $X$, maka terdapat
bilangan $m,M > 0$ (mungkin berbeda dari yang di atas) sedemikian sehingga
untuk setiap $x \in X$,
\begin{equation*}
m \snorm{x}_1
\leq
\snorm{x}_2
\leq
M \snorm{x}_1 .
\end{equation*}
\item
Tunjukkan bahwa $U \subset X$ terbuka dalam metrik yang didefinisikan oleh
$\norm{x-y}_1$ jika dan hanya jika $U$ terbuka dalam metrik yang didefinisikan oleh
$\norm{x-y}_2$.  Jadi, kekonvergenan barisan dan kekontinuan
fungsi sama untuk kedua norma.
\end{enumerate}
\end{exercise}

\begin{exercise}
Misalkan $A$ matriks segitiga atas.  Temukan rumus determinan
$A$ berdasarkan entri-entri diagonalnya, dan buktikan bahwa rumus tersebut benar.
\end{exercise}

\begin{exercise}
Diberikan matriks berukuran $n$-kali-$n$ yang dilambangkan dengan $A$, buktikan bahwa
$\babs{\det(A)} \leq \snorm{A}^n$ (norma yang digunakan untuk $A$ adalah norma operator).
Petunjuk: Salah satu caranya ialah terlebih dahulu membuktikannya untuk kasus $\snorm{A}=1$, yang
berarti semua kolom bernorma paling besar 1, lalu buktikan bahwa hal ini mengakibatkan
$\babs{\det(A)} \leq 1$ dengan menggunakan linearitas.
%Petunjuk: Perhatikan bahwa cukup membuktikan hal ini untuk matriks invertibel.  Kemudian
%jika perlu, susun ulang kolom-kolomnya dan faktorkan $A$ menjadi $n$ matriks, yang masing-masing
%berbeda dari identitas hanya pada satu kolom.
\end{exercise}

\begin{exercise}
Tinjau \propref{prop:finitedimpropinv}
dengan $X=\R^n$ (untuk setiap $n$) menggunakan norma Euklides.
\begin{enumerate}[a)]
\item
Buktikan bahwa taksiran $\snorm{A-B} < \frac{1}{\snorm{A^{-1}}}$ merupakan
taksiran terbaik yang mungkin:
Untuk setiap 
$A \in GL(\R^n)$,
carilah $B$ yang memenuhi kesamaan dan $B$ tidak invertibel.
Petunjuk:
Kesulitannya ialah $\snorm{A}\snorm{A^{-1}}$ tidak selalu sama dengan $1$.
Buktikan bahwa suatu vektor $x_1$ dapat dilengkapi menjadi basis
$\{ x_1,\ldots,x_n \}$ sedemikian sehingga $x_1 \cdot x_j = 0$ untuk $j \geq 2$.
Dengan pilihan $x_1$ yang tepat,
buat agar $(A-B)x_j = 0$ untuk $j \geq 2$.
\item
Untuk setiap $A \in GL(\R^n)$ yang tetap,
misalkan $\sM$ menyatakan himpunan matriks
$B$ sedemikian sehingga
$\snorm{A-B} < \frac{1}{\snorm{A^{-1}}}$.
Buktikan bahwa meskipun setiap $B \in \sM$ 
invertibel, $\snorm{B^{-1}}$ tidak terbatas
sebagai fungsi dari $B$ pada $\sM$.
\end{enumerate}
\end{exercise}

\begin{exnote}
Misalkan $A$ matriks berukuran $n$-kali-$n$.
Suatu $\lambda \in \C$ (mungkin kompleks bahkan untuk matriks real)
merupakan \emph{\myindex{nilai eigen}} dari $A$
jika terdapat vektor tak nol (mungkin kompleks) $x \in \C^n$ sedemikian sehingga
$Ax = \lambda x$
(perkalian dengan vektor kompleks didefinisikan sama seperti perkalian dengan vektor real;
jika $x = a+ib$ untuk vektor real $a$ dan $b$, dan $A$ matriks real, maka $Ax = Aa + i Ab$).
Bilangan
\begin{equation*}
\rho(A) \coloneqq
\sup \bigl\{ \sabs{\lambda} : \lambda \text{ adalah nilai eigen dari } A \bigr\}
\end{equation*}
merupakan \emph{\myindex{radius spektral}} dari $A$.  Di sini $\sabs{\lambda}$ adalah
modulus kompleks.  Kita menyatakan tanpa bukti bahwa setidaknya satu nilai eigen selalu ada,
dan tidak terdapat lebih dari $n$ nilai eigen $A$ yang berbeda.  Oleh karena itu, Anda dapat
mengasumsikan bahwa $0 \leq \rho(A) < \infty$.
Latihan-latihan di bawah berlaku untuk matriks kompleks, tetapi Anda boleh menganggap
matriks-matriks tersebut real.
\end{exnote}

\begin{exercise}
Misalkan $A,S$ matriks berukuran $n$-kali-$n$, dengan $S$ invertibel.
Buktikan bahwa $\lambda$ merupakan nilai eigen dari $A$ jika dan hanya
jika bilangan tersebut merupakan nilai eigen dari $S^{-1}AS$.  Kemudian buktikan bahwa
$\rho(S^{-1}AS) = \rho(A)$.
Secara khusus, $\rho$ merupakan fungsi yang terdefinisi dengan baik pada $L(X)$ untuk
setiap ruang vektor berdimensi hingga $X$.
\end{exercise}

\begin{exercise}
\pagebreak[2]
Misalkan $A$ matriks berukuran $n$-kali-$n$.
\begin{enumerate}[a)]
\item
Buktikan $\rho(A) \leq \snorm{A}$.  (Lihat definisi $\rho$ di atas.)
\item
Untuk setiap $k \in \N$, buktikan
$\rho(A) \leq \snorm{A^k}^{1/k}$.
\item
Misalkan $\displaystyle \lim_{k\to\infty} A^k = 0$ (limit diambil dalam norma operator).
Buktikan bahwa $\rho(A) < 1$.
\end{enumerate}
\end{exercise}

\begin{exercise}
Kita menyebut himpunan $C \subset \R^n$ \emph{simetris} jika
$x \in C$ mengakibatkan $-x \in C$.
\begin{enumerate}[a)]
\item
Misalkan $\snorm{\cdot}$ sebarang norma pada $\R^n$.  Tunjukkan bahwa bola satuan tertutup
$C(0,1)$ (dengan menggunakan metrik yang diinduksi oleh norma ini)
merupakan himpunan kompak, simetris, dan konveks.
\item (Menantang)
Misalkan $C \subset \R^n$ merupakan himpunan kompak,
simetris, dan konveks yang merentang $\R^n$, serta $0 \in C$.  Tunjukkan bahwa
\begin{equation*}
\snorm{x} \coloneqq \inf \left\{ \lambda : \lambda > 0 \text{ dan } \frac{x}{\lambda} \in C \right\}
\end{equation*}
merupakan norma pada $\R^n$, dan $C = C(0,1)$ (bola satuan tertutup) dalam metrik yang diinduksi oleh norma ini.
\end{enumerate}
Petunjuk: Anda boleh menggunakan hasil \exerciseref{exercise:allnormsequiv}
bagian c).  Secara khusus, sifat \myquote{kompak} suatu himpunan tidak bergantung pada
norma.
\end{exercise}

%%%%%%%%%%%%%%%%%%%%%%%%%%%%%%%%%%%%%%%%%%%%%%%%%%%%%%%%%%%%%%%%%%%%%%%%%%%%%%

\sectionnewpage
\section{Turunan}
\label{sec:svtheder}

%mbxINTROSUBSECTION

\sectionnotes{2--3 kuliah}

\subsection{Turunan}

Untuk fungsi $f \colon \R \to \R$, kita mendefinisikan
turunan di $x$ sebagai
\begin{equation*}
\lim_{h \to 0} \frac{f(x+h)-f(x)}{h} .
\end{equation*}
Dengan kata lain, terdapat suatu bilangan $a$ (turunan $f$ di $x$) sedemikian sehingga
\begin{equation*}
%mbxlatex \begin{aligned}
\lim_{h \to 0} \abs{\frac{f(x+h)-f(x)}{h} - a}
%mbxlatex &
=
\lim_{h \to 0} \abs{\frac{f(x+h)-f(x) - ah}{h}}
%mbxlatex \\
%mbxlatex &
=
\lim_{h \to 0} \frac{\babs{f(x+h)-f(x) - ah}}{\sabs{h}}
= 0.
%mbxlatex \end{aligned}
\end{equation*}

Perkalian dengan $a$ merupakan pemetaan linear dalam satu dimensi:
$h \mapsto ah$.
Dengan kata lain,
kita memandang $a \in L(\R^1,\R^1)$ sebagai
hampiran linear terbaik bagi
perubahan $f$ di sekitar $x$.  Kita menggunakan penafsiran ini
untuk memperluas diferensiasi ke lebih banyak variabel.

\begin{defn}
Misalkan $U \subset \R^n$ terbuka dan $f \colon U \to \R^m$ suatu fungsi.  Kita
mengatakan bahwa $f$ \emph{\myindex{diferensiabel}} di $x \in U$ jika terdapat
$A \in L(\R^n,\R^m)$ sedemikian sehingga
\begin{equation*}
\lim_{\substack{h \to 0\\h\in \R^n}}
\frac{\bnorm{f(x+h)-f(x) - Ah}}{\snorm{h}} = 0 .
\end{equation*}
Sebentar lagi akan kita tunjukkan bahwa $A$, jika ada, bersifat unik.
Kita menulis $Df(x) \coloneqq A$, atau $f'(x) \coloneqq A$,\glsadd{not:mvder} dan
mengatakan bahwa $A$ merupakan \emph{\myindex{turunan}} $f$ di $x$.
Jika $f$ diferensiabel di
setiap $x \in U$, kita cukup mengatakan bahwa $f$ \emph{diferensiabel}.  Lihat
\figureref{fig:svder} untuk ilustrasinya.
\end{defn}

\begin{myfigureht}
\myincludegraphics{svder}{%
Grafik contoh fungsi y sama dengan f dari x sub 1, x sub 2.
Sebuah titik tertentu ditandai pada bidang x sub 1 x sub 2.
Terdapat anak panah berlabel h pada bidang x sub 1 x sub 2.
Pada grafik ditampilkan bagian bidang singgung (yakni bidang yang menyinggung
grafik) di atas titik yang ditandai. Sebuah garis horizontal yang menyatakan h
kini digambar berawal pada grafik. Kemudian, dari ujung garis itu,
digambar anak panah vertikal yang membawa kita ke bidang singgung dan
diberi label A h.}
\caption{Ilustrasi turunan untuk fungsi $f \colon \R^2 \to \R$.  Vektor $h$ ditampilkan
pada bidang $x_1x_2$ dengan pangkal di $(x_1,x_2)$, dan vektor
$Ah \in \R^1$ ditampilkan sepanjang arah $y$.\label{fig:svder}}
\end{myfigureht}

Untuk fungsi diferensiabel,
turunan $f$ merupakan fungsi dari $U$ ke $L(\R^n,\R^m)$.  Bandingkan
dengan kasus satu dimensi, di mana turunannya merupakan fungsi
dari $U$ ke $\R$, tetapi sesungguhnya kita ingin memandang $\R$ di sini sebagai
$L(\R^1,\R^1)$.  Seperti dalam satu dimensi, gagasannya ialah bahwa suatu
pemetaan diferensiabel \myquote{dekat secara infinitesimal} dengan pemetaan linear, dan
pemetaan linear inilah turunannya.

Perhatikan norma-norma dalam definisi tersebut.
Norma pada
pembilang adalah norma pada $\R^m$, sedangkan norma pada penyebut adalah norma pada $\R^n$, tempat $h$
berada.
Biasanya, dari konteks dipahami bahwa $h \in \R^n$
(rumusnya tidak bermakna jika tidak demikian);
hal ini tidak selalu perlu dinyatakan secara eksplisit.
Seperti yang dijanjikan, mari kita buktikan bahwa turunannya unik.

\begin{prop}
Misalkan $U \subset \R^n$ suatu himpunan terbuka dan $f \colon U \to \R^m$ suatu fungsi.  Andaikan
$x \in U$ dan terdapat
$A,B \in L(\R^n,\R^m)$ sedemikian sehingga
\begin{equation*}
\lim_{h \to 0}
\frac{\bnorm{f(x+h)-f(x) - Ah}}{\snorm{h}} = 0
\qquad \text{dan} \qquad
\lim_{h \to 0}
\frac{\bnorm{f(x+h)-f(x) - Bh}}{\snorm{h}} = 0 .
\end{equation*}
Maka $A=B$.
\end{prop}

\begin{proof}
Misalkan $h \in \R^n$, $h \neq 0$.  Hitung
\begin{equation*}
\begin{split}
\frac{\bnorm{(A-B)h}}{\snorm{h}} & =
\frac{\bnorm{-\bigl(f(x+h)-f(x) - Ah\bigr) + f(x+h)-f(x) - Bh}}{\snorm{h}} \\
& \leq
\frac{\bnorm{f(x+h)-f(x) - Ah}}{\snorm{h}} + \frac{\bnorm{f(x+h)-f(x) -
Bh}}{\snorm{h}} .
\end{split}
\end{equation*}
Jadi
$\frac{\snorm{(A-B)h}}{\snorm{h}} \to 0$ ketika $h \to 0$.  Untuk
$\epsilon > 0$, bagi setiap $h$ tak nol dalam suatu bola-$\delta$ di sekitar
titik asal berlaku
\begin{equation*}
\epsilon > 
\frac{\bnorm{(A-B)h}}{\snorm{h}}
=
\norm{(A-B)\frac{h}{\snorm{h}}} .
\end{equation*}
Untuk sebarang $v \in \R^n$ dengan $\snorm{v}=1$,
jika $h = (\nicefrac{\delta}{2}) \, v$, maka $\snorm{h} < \delta$
dan $\frac{h}{\snorm{h}} = v$.
Jadi, $\bnorm{(A-B)v} < \epsilon$.  Dengan mengambil supremum atas semua $v$ dengan
$\snorm{v} = 1$, kita memperoleh norma operator
$\snorm{A-B} \leq \epsilon$.  Karena $\epsilon > 0$
dapat dipilih sebarang, $\snorm{A-B} = 0$, atau dengan kata lain $A = B$.
\end{proof}

\begin{example}
Jika $f(x) = Ax$ untuk suatu pemetaan linear $A$, maka
$f'(x) = A$:
\begin{equation*}
\frac{\bnorm{f(x+h)-f(x) - Ah}}{\snorm{h}}
=
\frac{\bnorm{A(x+h)-Ax - Ah}}{\snorm{h}}
=
\frac{0}{\snorm{h}} = 0 .
\end{equation*}
\end{example}

\begin{example}
Definisikan $f \colon \R^2 \to \R^2$ dengan
\begin{equation*}
f(x,y) = \bigl(f_1(x,y),f_2(x,y)\bigr) \coloneqq (1+x+2y+x^2,2x+3y+xy).
\end{equation*}
Mari kita tunjukkan bahwa $f$ diferensiabel di titik asal dan
hitung turunannya secara langsung menggunakan definisi.
Jika turunannya ada, turunan itu berada dalam $L(\R^2,\R^2)$, sehingga dapat
dinyatakan oleh matriks berukuran $2$-kali-$2$
$\left[\begin{smallmatrix}a&b\\c&d\end{smallmatrix}\right]$.  Misalkan $h =
(h_1,h_2)$.  Ungkapan berikut harus menuju nol.
\begin{multline*}
\frac{\bnorm{
f(h_1,h_2)-f(0,0)
-
(ah_1 +bh_2 , ch_1+dh_2)}
}{\bnorm{(h_1,h_2)}}
=
\\
\frac{\sqrt{
{\bigl((1-a)h_1 + (2-b)h_2 + h_1^2\bigr)}^2
+
{\bigl((2-c)h_1 + (3-d)h_2 + h_1h_2\bigr)}^2}}{\sqrt{h_1^2+h_2^2}} .
\end{multline*}
Jika kita memilih $a=1$, $b=2$, $c=2$, $d=3$, ungkapan tersebut menjadi
\begin{equation*}
\frac{\sqrt{
h_1^4 + h_1^2h_2^2}}{\sqrt{h_1^2+h_2^2}}
=
\sabs{h_1}
\frac{\sqrt{
h_1^2 + h_2^2}}{\sqrt{h_1^2+h_2^2}}
= \sabs{h_1} .
\end{equation*}
Ungkapan ini memang menuju nol ketika $h \to 0$.  Fungsi
$f$ diferensiabel di titik asal dan
turunan $f'(0)$ dinyatakan oleh matriks
$\left[\begin{smallmatrix}1&2\\2&3\end{smallmatrix}\right]$.
\end{example}

\begin{prop}
Misalkan $U \subset \R^n$ terbuka dan $f \colon U \to \R^m$
diferensiabel di $p \in U$.  Maka $f$ kontinu di $p$.
\end{prop}

\begin{proof}
Sifat diferensiabel $f$ di $p$ dapat dituliskan dengan cara lain, yaitu dengan meninjau
\begin{equation*}
r(h) \coloneqq f(p+h)-f(p) - f'(p) h .
\end{equation*}
Fungsi $f$ diferensiabel di $p$ jika
$\frac{\snorm{r(h)}}{\snorm{h}}$ menuju nol ketika $h \to 0$,
sehingga
$r(h)$ sendiri menuju nol.  Pemetaan $h \mapsto f'(p) h$
merupakan pemetaan linear antarruang berdimensi hingga, sehingga kontinu,
dan $f'(p) h \to 0$ ketika $h \to 0$.  Dengan demikian,
$f(p+h)$ harus menuju $f(p)$ ketika $h \to 0$.  Artinya, $f$ kontinu di $p$.
\end{proof}

Operasi diferensiasi merupakan operator linear pada ruang fungsi-fungsi
diferensiabel.

\begin{prop}
Andaikan $U \subset \R^n$ terbuka,
$f \colon U \to \R^m$ dan
$g \colon U \to \R^m$ diferensiabel di $p \in U$,
serta $\alpha \in \R$.  Maka fungsi $f+g$ dan $\alpha f$
diferensiabel di $p$, dan
\begin{equation*}
(f+g)'(p) = f'(p) + g'(p) , \qquad \text{dan} \qquad (\alpha f)'(p) = \alpha
f'(p) .
\end{equation*}
\end{prop}

\begin{proof}
Misalkan $h \in \R^n$, $h \neq 0$.  Maka
\begin{multline*}
\frac{\bnorm{f(p+h)+g(p+h)-\bigl(f(p)+g(p)\bigr) - \bigl(f'(p) + g'(p)\bigr)h}}{\snorm{h}}
\\
\leq
\frac{\bnorm{f(p+h)-f(p) - f'(p)h}}{\snorm{h}}
+
\frac{\bnorm{g(p+h)-g(p) - g'(p)h}}{\snorm{h}} ,
\end{multline*}
dan
\begin{equation*}
\frac{\bnorm{\alpha f(p+h) - \alpha f(p) - \alpha f'(p)h}}{\snorm{h}}
=
\sabs{\alpha} \frac{\bnorm{f(p+h)-f(p) - f'(p)h}}{\snorm{h}} .
\end{equation*}
Berdasarkan hipotesis, limit ruas-ruas kanan ketika $h$ menuju nol adalah nol.
Hasil yang diinginkan pun diperoleh.
\end{proof}

Jika $A \in L(\R^n,\R^m)$ dan $B \in L(\R^m,\R^k)$ merupakan pemetaan linear, maka
masing-masing merupakan turunannya sendiri.  Komposisi
$BA \in L(\R^n,\R^k)$ juga merupakan turunannya sendiri, sehingga
turunan komposisi merupakan komposisi
turunan-turunannya.  Karena pemetaan diferensiabel
\myquote{dekat secara infinitesimal}
dengan pemetaan linear, pemetaan semacam itu memiliki sifat yang sama:

\begin{thm}[Aturan rantai] \index{aturan rantai}
Misalkan $U \subset \R^n$ dan $V \subset \R^m$ himpunan terbuka, $f \colon U \to
\R^m$ diferensiabel di $p \in U$, $f(U) \subset V$,
dan $g \colon V \to \R^\ell$ diferensiabel
di $f(p)$.  Maka $F \colon U \to \R^{\ell}$ yang didefinisikan oleh
\begin{equation*}
F(x) \coloneqq g\bigl(f(x)\bigr)
\end{equation*}
diferensiabel di $p$, dan
\begin{equation*}
F'(p) = g'\bigl(f(p)\bigr) f'(p) .
\end{equation*}
\end{thm}

Tanpa menuliskan titik-titik tempat fungsi dievaluasi, kita menulis
$F' = {(g \circ f)}' = g' f'$.
Turunan komposisi $g \circ f$
merupakan komposisi turunan $g$ dan $f$:
Jika $f'(p) = A$ dan $g'\bigl(f(p)\bigr) = B$, maka $F'(p) = BA$,
sama seperti untuk pemetaan linear.

\begin{proof}
Misalkan $A \coloneqq f'(p)$ dan $B \coloneqq g'\bigl(f(p)\bigr)$.  Ambil $h \in \R^n$ yang tak nol
dan tulis $q \coloneqq f(p)$, $k \coloneqq f(p+h)-f(p)$.  Definisikan
\begin{equation*}
r(h) \coloneqq f(p+h)-f(p) - A h . %= k - Ah.
\end{equation*}
Maka $r(h) = k-Ah$ atau $Ah = k-r(h)$, dan $f(p+h) = q+k$.
Kita meninjau besaran yang perlu
menuju nol:
\begin{equation*}
\begin{split}
\frac{\bnorm{F(p+h)-F(p) - BAh}}{\snorm{h}}
& =
\frac{\bnorm{g\bigl(f(p+h)\bigr)-g\bigl(f(p)\bigr) - BAh}}{\snorm{h}}
\\
& =
\frac{\bnorm{g(q+k)-g(q) - B\bigl(k-r(h)\bigr)}}{\snorm{h}}
\\
& \leq
\frac
{\bnorm{g(q+k)-g(q) - Bk}}
{\snorm{h}}
+
\snorm{B}
\frac
{\bnorm{r(h)}}
{\snorm{h}}
.
\end{split}
\end{equation*}
Kedua suku di ruas kanan perlu menuju 0
ketika $h$ menuju 0.
Pertama, $\snorm{B}$ merupakan konstanta dan $f$ diferensiabel di $p$,
sehingga
suku $\snorm{B}\frac{\snorm{r(h)}}{\snorm{h}}$ menuju 0.
Selanjutnya,
jika $k=0$, maka
$\frac
{\snorm{g(q+k)-g(q) - Bk}}
{\snorm{h}} = 0$.
Jadi, andaikan $k \neq 0$.
Maka
\begin{equation*}
\frac
{\bnorm{g(q+k)-g(q) - Bk}}
{\snorm{h}}
 =
\frac
{\bnorm{g(q+k)-g(q) - Bk}}
{\snorm{k}}
\frac
{\bnorm{f(p+h)-f(p)}}
{\snorm{h}} .
\end{equation*}
Karena $f$ kontinu di $p$,
$k$ menuju 0 ketika $h$ menuju 0.
Dengan demikian,
$\frac
{\snorm{g(q+k)-g(q) - Bk}}
{\snorm{k}}$ menuju 0 karena $g$ diferensiabel di $q$.
Kita memiliki
\begin{equation*}
%mbxlatex \begin{aligned}
\frac
{\bnorm{f(p+h)-f(p)}}
{\snorm{h}}
%mbxlatex &
\leq
\frac
{\bnorm{f(p+h)-f(p)-Ah}}
{\snorm{h}}
+
\frac
{\snorm{Ah}}
{\snorm{h}}
%mbxlatex \\
%mbxlatex &
\leq
\frac
{\bnorm{f(p+h)-f(p)-Ah}}
{\snorm{h}}
+
\snorm{A} .
%mbxlatex \end{aligned}
\end{equation*}
Karena $f$ diferensiabel di $p$,
untuk $h$ yang cukup kecil, besaran
$\frac{\snorm{f(p+h)-f(p)-Ah}}{\snorm{h}}$ terbatas.  Oleh karena itu,
suku
$
\frac
{\snorm{f(p+h)-f(p)}}
{\snorm{h}}
$
tetap terbatas ketika $h$ menuju 0.
Akibatnya, suku
$\frac
{\snorm{g(q+k)-g(q) - Bk}}
{\snorm{h}}$ menuju nol ketika $h$ menuju 0.
Oleh karena itu,
$\frac{\snorm{F(p+h)-F(p) - BAh}}{\snorm{h}}$ menuju nol, dan
$F'(p) = BA$, seperti yang diklaim.
\end{proof}

\subsection{Turunan parsial}

Ada cara lain untuk memperumum turunan dari satu dimensi.
Kita mempertahankan semua variabel kecuali satu agar tetap konstan, lalu mengambil
turunan satu variabel seperti biasa.

\begin{defn}
Misalkan
$f \colon U \to \R$ suatu fungsi pada himpunan terbuka $U \subset \R^n$.
Jika limit berikut ada, kita tulis
\glsadd{not:partialder}
\begin{equation*}
\frac{\partial f}{\partial x_j} (x) \coloneqq 
\lim_{h\to 0}\frac{f(x_1,\ldots,x_{j-1},x_j+h,x_{j+1},\ldots,x_n)-f(x)}{h}
=
\lim_{h\to 0}\frac{f(x+h e_j)-f(x)}{h} .
\end{equation*}
Kita menyebut
$\frac{\partial f}{\partial x_j} (x)$ sebagai \emph{\myindex{turunan parsial}}
dari $f$
terhadap $x_j$.  %Kadang-kadang kita menulis $D_j f$ sebagai gantinya.
Lihat \figureref{fig:svpartder}.
Di sini $h$ adalah bilangan, bukan vektor.

Untuk pemetaan $f \colon U \to \R^m$, kita tulis
$f = (f_1,f_2,\ldots,f_m)$, dengan $f_k$ merupakan fungsi-fungsi
bernilai real.  Kemudian kita mengambil turunan parsial dari
komponen-komponennya,
$\frac{\partial f_k}{\partial x_j}$.
\end{defn}

\begin{myfigureht}
\myincludegraphics{svpartder}{%
Grafik suatu fungsi contoh y sama dengan f dari x subskrip 1, x subskrip 2.
Sebuah titik tertentu ditandai pada bidang x subskrip 1 x subskrip 2.
Sebuah bidang dalam arah x subskrip 2 dan y ditandai dengan garis titik-titik,
dan irisan grafik fungsi f dengan bidang ini
ditandai dengan garis yang lebih tebal.  Titik pada grafik di atas
titik x subskrip 1 x subskrip 2 tersebut juga ditandai.  Melalui titik ini
dan di dalam bidang bertitik-titik digambar garis singgung pada grafik.
Garis itu diberi label bahwa kemiringannya sama dengan turunan parsial f
terhadap x subskrip 2 di titik x subskrip 1, x subskrip 2.}
\caption{Ilustrasi turunan parsial untuk fungsi $f \colon \R^2
\to \R$.  Bidang $yx_2$ dengan $x_1$ tetap ditandai oleh garis titik-titik,
dan kemiringan garis singgung pada bidang $yx_2$ adalah
$\frac{\partial f}{\partial x_2}(x_1,x_2)$.\label{fig:svpartder}}
\end{myfigureht}

Turunan parsial lebih mudah dihitung dengan seluruh perangkat
kalkulus, dan menyediakan cara untuk menghitung turunan suatu
fungsi.

\begin{prop} \label{mv:prop:jacobianmatrix}
Misalkan $U \subset \R^n$ terbuka dan $f \colon U \to \R^m$
diferensiabel di $p \in U$.  Maka semua turunan parsial di $p$
ada dan, terhadap basis standar $\R^n$ dan $\R^m$,
$f'(p)$ direpresentasikan oleh matriks
\begin{equation*}
\begin{bmatrix}
\frac{\partial f_1}{\partial x_1}(p)
&
\frac{\partial f_1}{\partial x_2}(p)
& \ldots &
\frac{\partial f_1}{\partial x_n}(p)
\\[6pt]
\frac{\partial f_2}{\partial x_1}(p)
&
\frac{\partial f_2}{\partial x_2}(p)
& \ldots &
\frac{\partial f_2}{\partial x_n}(p)
\\
\vdots & \vdots & \ddots & \vdots
\\
\frac{\partial f_m}{\partial x_1}(p)
&
\frac{\partial f_m}{\partial x_2}(p)
& \ldots &
\frac{\partial f_m}{\partial x_n}(p)
\end{bmatrix} .
\end{equation*}
\end{prop}


Dengan kata lain,
\begin{equation*}
f'(p) \, e_j =
\sum_{k=1}^m
\frac{\partial f_k}{\partial x_j}(p) \,e_k .
\end{equation*}
Jika $v = \sum_{j=1}^n c_j\, e_j = (c_1,c_2,\ldots,c_n)$, maka
\begin{equation*}
f'(p) \, v =
\sum_{j=1}^n
\sum_{k=1}^m
 c_j
\frac{\partial f_k}{\partial x_j}(p) \,e_k
=
\sum_{k=1}^m
\left(
\sum_{j=1}^n
 c_j
\frac{\partial f_k}{\partial x_j}(p) \right) \,e_k .
\end{equation*}

\begin{proof}
Tetapkan suatu $j$ dan perhatikan bahwa untuk $h$ tak nol,
\begin{equation*}
\begin{split}
\norm{\frac{f(p+h e_j)-f(p)}{h} - f'(p) \, e_j} & = 
\norm{\frac{f(p+h e_j)-f(p) - f'(p) \, h e_j}{h}} \\
& =
\frac{\bnorm{f(p+h e_j)-f(p) - f'(p) \, h e_j}}{\snorm{h e_j}} .
\end{split}
\end{equation*}
Ketika $h$ menuju 0, ruas kanan menuju nol berdasarkan
kediferensiabelan $f$.  Oleh karena itu,
\begin{equation*}
\lim_{h \to 0}
\frac{f(p+h e_j)-f(p)}{h} = f'(p) \, e_j  .
\end{equation*}
Limit tersebut diambil dalam $\R^m$.
Nyatakan $f$ melalui komponen-komponennya,
$f = (f_1,f_2,\ldots,f_m)$.
Mengambil limit dalam $\R^m$
sama dengan mengambil limit pada setiap komponen secara terpisah.
Jadi, untuk setiap $k$,
turunan parsial
\begin{equation*}
\frac{\partial f_k}{\partial x_j} (p)
=
\lim_{h \to 0}
\frac{f_k(p+h e_j)-f_k(p)}{h}
\end{equation*}
ada dan sama dengan komponen ke-$k$ dari $f'(p)\, e_j$, yang merupakan
kolom ke-$j$ dari $f'(p)$,
dan pembuktian selesai.
\end{proof}

Kebalikan proposisi tersebut tidak benar.  Adanya turunan-turunan parsial
tidak berarti bahwa fungsi tersebut diferensiabel.  Lihat
latihan.
Namun, ketika turunan-turunan parsialnya kontinu, kita akan membuktikan bahwa
kebalikannya berlaku.
Salah satu akibat proposisi di atas ialah bahwa jika $f$
diferensiabel pada $U$, maka $f' \colon U \to
L(\R^n,\R^m)$ merupakan fungsi kontinu jika dan hanya jika
semua $\frac{\partial f_k}{\partial x_j}$ merupakan fungsi kontinu.

\subsection{Gradien, kurva, dan turunan arah}

Misalkan $U \subset \R^n$ terbuka dan $f \colon U \to \R$ suatu fungsi
diferensiabel.  Kita mendefinisikan
\emph{\myindex{gradien}} sebagai
\glsadd{not:gradient}
\begin{equation*}
\nabla f (x) \coloneqq \sum_{j=1}^n \frac{\partial f}{\partial x_j} (x)\, e_j .
\end{equation*}
Gradien menyediakan cara untuk merepresentasikan aksi
turunan sebagai hasil kali titik: $f'(x)\,v = \nabla f(x) \cdot v$.

Misalkan $\gamma \colon (a,b) \subset \R \to \R^n$ diferensiabel.
Fungsi semacam ini (dan citranya) kadang-kadang disebut \emph{\myindex{kurva}},
atau \emph{\myindex{kurva diferensiabel}}.
Tulis $\gamma =
(\gamma_1,\gamma_2,\ldots,\gamma_n)$.
Untuk keperluan perhitungan,
kita mengidentifikasi $L(\R^1)$ dan $\R$ sebagaimana ketika mendefinisikan
turunan satu variabel.
Kita juga mengidentifikasi $L(\R^1,\R^n)$ dengan $\R^n$.
Kita memperlakukan $\gamma'(t)$ sekaligus sebagai operator dalam
$L(\R^1,\R^n)$ dan sebagai vektor
$\bigl(\gamma_1'(t),\gamma_2'(t),\ldots,\gamma_n'(t)\bigr)$
dalam $\R^n$.
Berdasarkan \propref{mv:prop:jacobianmatrix},
jika $v\in \R^n$ adalah $\gamma'(t)$ yang diperlakukan sebagai vektor,
maka $h \mapsto h \, v$ (untuk $h \in \R^1 = \R$) adalah
$\gamma'(t)$ yang diperlakukan sebagai operator
dalam $L(\R^1,\R^n)$.
Kita sering menggunakan sedikit
penyalahgunaan notasi ini ketika membahas kurva.
Vektor $\gamma'(t)$ disebut \emph{\myindex{vektor singgung}}.
Lihat \figureref{fig:difcurveder}.
\begin{myfigureht}
\myincludepdft{diffcurveder}{%
Sebuah kurva digambar pada bidang.
Awal kurva ditandai sebagai gamma dari a,
dan ujungnya ditandai sebagai gamma dari b.  Seluruh kurva
ditandai gamma dari (a,b).
Sebuah titik pada kurva ditandai gamma dari t, dan sebuah anak panah
yang berawal di titik ini serta menunjuk sepanjang kurva
ditandai gamma aksen dari t.}
\caption{Kurva diferensiabel dan turunannya sebagai
vektor (agar jelas, diasumsikan $\gamma$ terdefinisi pada $[a,b]$).
Vektor singgung $\gamma'(t)$ menunjuk sepanjang kurva.\label{fig:difcurveder}}
\end{myfigureht}

Misalkan $\gamma\bigl((a,b)\bigr) \subset U$ dan tetapkan
\begin{equation*}
g(t) \coloneqq f\bigl(\gamma(t)\bigr) .
\end{equation*}
Fungsi
$g$ diferensiabel. Dengan memperlakukan $g'(t)$ sebagai bilangan,
\begin{equation*}
g'(t) =
f'\bigl(\gamma(t)\bigr) \gamma'(t)
=
\sum_{j=1}^n
\frac{\partial f}{\partial x_j} \bigl(\gamma(t)\bigr)
\frac{d\gamma_j}{dt} (t)
=
\sum_{j=1}^n
\frac{\partial f}{\partial x_j}
\frac{d\gamma_j}{dt} .
\end{equation*}
Demi kemudahan,
kita sering
menghilangkan titik tempat evaluasi dilakukan,
seperti pada ruas paling kanan di atas.
Dengan notasi gradien dan hasil kali titik,
persamaan tersebut menjadi
\begin{equation*}
g'(t) = (\nabla f) \bigl(\gamma(t)\bigr) \cdot \gamma'(t)
= \nabla f \cdot \gamma'.
\end{equation*}

Kita menggunakan gagasan ini untuk mendefinisikan turunan dalam arah tertentu.  Suatu arah
hanyalah vektor yang menunjuk ke arah tersebut.  Pilih vektor $u \in \R^n$
sedemikian sehingga $\snorm{u} = 1$, dan tetapkan $x \in U$.
Kita mendefinisikan
\emph{\myindex{turunan arah}} sebagai
\glsadd{not:mvdirder}
\begin{equation*}
D_u f (x) \coloneqq \frac{d}{dt}\Big|_{t=0} \bigl[ f(x+tu) \bigr] =
\lim_{h\to 0}
\frac{f(x+hu)-f(x)}{h} ,
\end{equation*}
di mana notasi
$\frac{d}{dt}\big|_{t=0}$ menyatakan turunan yang dievaluasi pada $t=0$.
Ketika $u=e_j$ merupakan vektor basis standar, kita memperoleh
$\frac{\partial f}{\partial x_j} = D_{e_j} f$.
Karena itu, notasi $\frac{\partial f}{\partial u}$ kadang-kadang
digunakan sebagai pengganti $D_u f$.

Definisikan $\gamma$ dengan
\begin{equation*}
\gamma(t) \coloneqq x + tu .
\end{equation*}
Maka $\gamma'(t) = u$ untuk setiap $t$.
Mari kita lihat apa yang terjadi pada $f$ ketika kita bergerak sepanjang $\gamma$:
\begin{equation*}
D_u f (x) =
\frac{d}{dt}\Big|_{t=0} \bigl[ f(x+tu) \bigr] =
(\nabla f) \bigl(\gamma(0)\bigr) \cdot \gamma'(0)
=
(\nabla f) (x) \cdot u .
\end{equation*}
Faktanya, perhitungan ini berlaku untuk sebarang kurva $\gamma$ sedemikian sehingga
$\gamma(0) = x$ dan $\gamma'(0) = u$.

Misalkan $(\nabla f)(x) \neq 0$.
Berdasarkan ketaksamaan Cauchy--Schwarz,
\begin{equation*}
\babs{D_u f(x)} \leq \bnorm{(\nabla f)(x)} .
\end{equation*}
Kesamaan tercapai ketika $u$ merupakan kelipatan skalar dari
$(\nabla f)(x)$.  Artinya, ketika
\begin{equation*}
u = 
\frac{(\nabla f)(x)}{\bnorm{(\nabla f)(x)}} ,
\end{equation*}
kita memperoleh $D_u f(x) = \bnorm{(\nabla f)(x)}$.
Gradien menunjuk ke arah tempat
fungsi bertambah paling cepat, atau dengan kata lain,
ke arah tempat $D_u f(x)$ bernilai maksimum.

\subsection{Jacobian}

\begin{defn}
Misalkan $U \subset \R^n$ dan
$f \colon U \to \R^n$ suatu pemetaan diferensiabel.  Definisikan
\emph{\myindex{determinan Jacobian}}%
\footnote{Dinamai menurut matematikawan Italia
\href{https://en.wikipedia.org/wiki/Carl_Gustav_Jacob_Jacobi}{Carl Gustav Jacob Jacobi}
(1804--1851).},
atau cukup
\emph{\myindex{Jacobian}}%
\footnote{Matriks dalam \propref{mv:prop:jacobianmatrix} yang merepresentasikan $f'(x)$
disebut
\emph{\myindex{matriks Jacobian}}, atau kadang-kadang, secara membingungkan, juga disebut cukup
\myquote{Jacobian}.},
dari $f$ di $x$ sebagai
\glsadd{not:jacobdet}
\begin{equation*}
J_f(x) \coloneqq \det\bigl( f'(x) \bigr) .
\end{equation*}
Kadang-kadang $J_f$ ditulis sebagai
\begin{equation*}
\frac{\partial(f_1,f_2,\ldots,f_n)}{\partial(x_1,x_2,\ldots,x_n)} .
\end{equation*}
\end{defn}

Notasi terakhir ini mungkin tampak agak membingungkan,
tetapi sangat berguna ketika kita perlu menyatakan
secara tepat variabel dan komponen fungsi yang digunakan,
seperti yang akan kita lakukan, misalnya, dalam teorema fungsi implisit.

Determinan Jacobian $J_f$ merupakan fungsi bernilai real, dan ketika $n=1$, fungsi ini tidak lain adalah
turunan.
Dari aturan rantai dan fakta bahwa $\det(AB) = \det(A)\det(B)$, diperoleh:
\begin{equation*}
J_{f \circ g} (x) = J_f\bigl(g(x)\bigr) J_g(x) .
\end{equation*}

Determinan suatu pemetaan linear memberi tahu kita apa yang terjadi pada
luas atau volume di bawah pemetaan tersebut.
Demikian pula, determinan Jacobian mengukur seberapa besar suatu pemetaan diferensiabel
meregangkan ruang secara lokal dan apakah pemetaan itu membalik orientasi.  Secara khusus, jika determinan
Jacobian
tak nol, kita menduga bahwa pemetaan tersebut invertibel secara lokal (dan
dugaan ini benar, sebagaimana akan kita lihat nanti).

\subsection{Latihan}

\begin{exercise}
Misalkan $\gamma \colon (-1,1) \to \R^n$ dan
$\alpha \colon (-1,1) \to \R^n$ dua kurva diferensiabel
sedemikian sehingga $\gamma(0) = \alpha(0)$ dan $\gamma'(0) = \alpha'(0)$.
Misalkan $F \colon \R^n \to \R$ suatu fungsi diferensiabel.  Tunjukkan bahwa
\begin{equation*}
\frac{d}{dt}\Big|_{t=0}
F\bigl(\gamma(t)\bigr)
=
\frac{d}{dt}\Big|_{t=0}
F\bigl(\alpha(t)\bigr)
.
\end{equation*}
\end{exercise}

\begin{exercise}
Misalkan $f \colon \R^2 \to \R$ diberikan oleh
$f(x,y)
\coloneqq
\sqrt{x^2+y^2}$,
lihat \figureref{fig:distfromorgfunc}.
Tunjukkan bahwa $f$ tidak diferensiabel di titik asal.
\end{exercise}

\begin{myfigureht}
\myincludegraphics{distfromorgfunc}{%
Grafik dalam ruang xyz dari fungsi z sama dengan f dari x, y, yang berbentuk kerucut
menghadap ke atas dari titik asal.  Artinya, grafik ini dapat dibentuk
dengan naik pada sudut 45 derajat sepanjang setiap arah dalam bidang xy menjauhi
titik asal.}
\caption{Grafik $\sqrt{x^2+y^2}$.\label{fig:distfromorgfunc}}
\end{myfigureht}

\begin{samepage}
\begin{exercise}
Dengan hanya menggunakan definisi turunan, tunjukkan bahwa
fungsi-fungsi $f \colon \R^2 \to \R^2$ berikut diferensiabel di titik asal dan
tentukan turunannya.
\begin{enumerate}[a)]
\item
$f(x,y) \coloneqq (1+x+xy,x)$,
\item
$f(x,y) \coloneqq \bigl(y-y^{10},x \bigr)$,
\item
$f(x,y) \coloneqq \bigl( {(x+y+1)}^2 , {(x-y+2)}^2 \bigr)$.
\end{enumerate}
\end{exercise}
\end{samepage}

\begin{exercise}
Misalkan $f \colon \R \to \R$ dan $g \colon \R \to \R$ fungsi-fungsi diferensiabel.
Dengan hanya menggunakan definisi turunan, tunjukkan bahwa
$h \colon \R^2 \to \R^2$ yang didefinisikan oleh $h(x,y)
\coloneqq \bigl(f(x),g(y)\bigr)$ merupakan fungsi diferensiabel, dan tentukan
turunannya di setiap titik $(x,y)$.
\end{exercise}

\begin{exercise} \label{exercise:noncontpartialsexist}
Definisikan fungsi $f \colon \R^2 \to \R$ dengan
(lihat \figureref{fig:xyxsqysqvol2})
\begin{equation*}
f(x,y)
\coloneqq
\begin{cases}
\frac{xy}{x^2+y^2} & \text{jika } (x,y) \neq (0,0), \\
0                  & \text{jika } (x,y) = (0,0).
\end{cases}
\end{equation*}
\begin{enumerate}[a)]
\item
Tunjukkan bahwa turunan parsial
$\frac{\partial f}{\partial x}$ dan
$\frac{\partial f}{\partial y}$ ada di setiap titik (termasuk titik asal).
\item
Tunjukkan bahwa $f$ tidak kontinu di titik asal (dan karenanya tidak
diferensiabel).
\end{enumerate}
\end{exercise}

\begin{myfigureht}
\myincludegraphics{xyxsqysq}{%
Grafik fungsi z sama dengan f dari x, y.
Fungsi tersebut bernilai nol pada sumbu x dan y.
Di kuadran pertama dan ketiga nilainya positif,
sedangkan di kuadran kedua dan keempat nilainya negatif.
Fungsi ini terbatas, tetapi tidak kontinu di titik asal:
jika didekati sepanjang garis y sama dengan x, fungsi tampak
konstan dan positif, sedangkan sepanjang garis y sama dengan minus x,
fungsi tampak konstan dan negatif.  Grafiknya tampak seolah-olah
selembar karet di titik asal ditarik dengan satu tangan dari atas dan
satu tangan dari bawah sehingga titik asal membentuk
lekukan yang memanjang.}
\caption{Grafik $\frac{xy}{x^2+y^2}$.\label{fig:xyxsqysqvol2}}
\end{myfigureht}

\begin{samepage}
\begin{exercise}
Definisikan fungsi $f \colon \R^2 \to \R$ dengan
(lihat \figureref{fig:xsqyxsqysq})
\begin{equation*}
f(x,y)
\coloneqq
\begin{cases}
\frac{x^2y}{x^2+y^2} & \text{jika } (x,y) \neq (0,0), \\
0                    & \text{jika } (x,y) = (0,0).
\end{cases}
\end{equation*}
\begin{enumerate}[a)]
\item
Tunjukkan bahwa turunan parsial
$\frac{\partial f}{\partial x}$ dan
$\frac{\partial f}{\partial y}$ ada di setiap titik.
\item
Tunjukkan bahwa untuk setiap $u \in \R^2$ dengan $\snorm{u}=1$, turunan
arah $D_u f$ ada di setiap titik.
\item
Tunjukkan bahwa $f$ kontinu di titik asal.
\item
Tunjukkan bahwa $f$ tidak diferensiabel di titik asal.
\end{enumerate}
\end{exercise}
\end{samepage}

\begin{myfigureht}
\myincludegraphics{xsqyxsqysq}{%
Grafik fungsi z sama dengan f dari x, y.
Fungsi tersebut bernilai nol pada sumbu x dan y.
Di kuadran pertama dan kedua nilainya positif,
sedangkan di kuadran ketiga dan keempat nilainya negatif.
Meskipun jelas kontinu, gambar memperlihatkan
bahwa bentuknya tetap sama ketika kita terus memperbesar
daerah di sekitar titik asal, sehingga grafik tidak
memiliki bidang singgung di titik asal.}
\caption{Grafik $\frac{x^2y}{x^2+y^2}$.\label{fig:xsqyxsqysq}}
\end{myfigureht}

\begin{samepage}
\begin{exercise}
Misalkan $f \colon \R^n \to \R^n$ injektif, surjektif, diferensiabel di setiap
titik, dan $f^{-1}$ juga diferensiabel di setiap titik.
\begin{enumerate}[a)]
\item
Tunjukkan bahwa $f'(p)$ invertibel di setiap titik $p$ dan hitung
${(f^{-1})}'\bigl(f(p)\bigr)$.  Petunjuk: Tinjau $x = f^{-1}\bigl(f(x)\bigr)$.
\item
Misalkan $g \colon \R^n \to \R^n$ suatu fungsi yang diferensiabel di $q \in \R^n$
dan memenuhi $g(q)=q$.  Misalkan $f(p) = q$ untuk suatu $p \in \R^n$.
Tunjukkan $J_g(q) = J_{f^{-1} \circ g \circ f}(p)$, dengan $J_g$ menyatakan determinan
Jacobian.
\end{enumerate}
\end{exercise}
\end{samepage}

\begin{exercise}
Misalkan $f \colon \R^2 \to \R$ diferensiabel dan memenuhi
$f(x,y) = 0$ jika dan hanya jika $y=0$, serta $\nabla f(0,0) = (0,1)$.
Buktikan bahwa $f(x,y) > 0$ setiap kali $y > 0$, dan
$f(x,y) < 0$ setiap kali $y < 0$.
\end{exercise}

\begin{exnote}
\pagebreak[2]
Seperti pada fungsi satu variabel, $f \colon U \to \R$ memiliki
\emph{\myindex{maksimum relatif}} di $p \in U$ jika terdapat
$\delta >0$ sedemikian sehingga $f(q) \leq f(p)$ untuk setiap $q \in B(p,\delta) \cap U$.
Definisi \emph{\myindex{minimum relatif}} serupa.
\end{exnote}

\begin{exercise} \label{exercise:mv:maximumcritical}
Misalkan $U \subset \R^n$ terbuka dan
$f \colon U \to \R$ diferensiabel.  Misalkan $f$ memiliki maksimum relatif
di $p \in U$.  Tunjukkan bahwa $f'(p) = 0$, yakni pemetaan nol dalam
$L(\R^n,\R)$.  Dengan kata lain, $p$ merupakan
\emph{\myindex{titik kritis}} dari $f$.
\end{exercise}

\begin{exercise}
Misalkan $f \colon \R^2 \to \R$ diferensiabel dan
$f(x,y) = 0$ setiap kali $x^2+y^2 = 1$.
Buktikan bahwa terdapat setidaknya
satu titik $(x_0,y_0)$ sedemikian sehingga
$\frac{\partial f}{\partial x}(x_0,y_0) = \frac{\partial f}{\partial
y}(x_0,y_0) = 0$.
\end{exercise}

\begin{exercise} \label{exercise:peano}
Definisikan $f(x,y) \coloneqq ( x-y^2 ) ( 2 y^2 - x)$.  Grafik $f$ disebut
\emph{\myindex{permukaan Peano}}.%
\footnote{Dinamai menurut matematikawan Italia
\href{https://en.wikipedia.org/wiki/Giuseppe_Peano}{Giuseppe Peano}
(1858--1932).}
\begin{enumerate}[a)]
\item
Tunjukkan bahwa
$(0,0)$ merupakan titik kritis, yakni $f'(0,0) = 0$, pemetaan
linear nol dalam $L(\R^2,\R)$.
\item
Tunjukkan bahwa untuk setiap arah, pembatasan $f$ pada garis yang melalui titik asal
dalam arah tersebut memiliki maksimum relatif di titik asal.  Dengan kata lain,
untuk setiap $(x,y)$ sedemikian sehingga $x^2+y^2=1$, fungsi $g(t) \coloneqq f(tx,ty)$
memiliki maksimum relatif di $t=0$.\\
Petunjuk: Meskipun tidak diperlukan,
\volIref{\sectionref*{vI-sec:taylor} dari jilid I}{\sectionref{sec:taylor}}
membuat
bagian ini lebih mudah.
\item
Tunjukkan bahwa $f$ tidak memiliki maksimum relatif di $(0,0)$.
\end{enumerate}
\end{exercise}

\begin{exercise}
Misalkan $f \colon \R \to \R^n$ diferensiabel dan $\bnorm{f(t)} = 1$ untuk
setiap $t$ (artinya, kita memiliki kurva pada permukaan bola satuan).  Tunjukkan bahwa
$f'(t) \cdot f(t) = 0$ (dengan memperlakukan $f'(t)$ sebagai vektor) untuk setiap $t$.
\end{exercise}

\begin{exercise}
Definisikan $f \colon \R^2 \to \R^2$ dengan $f(x,y) \coloneqq
\bigl(x,y+\varphi(x)\bigr)$ untuk suatu fungsi satu variabel $\varphi$ yang
diferensiabel.  Tunjukkan bahwa $f$ diferensiabel dan tentukan $f'$.
\end{exercise}

\begin{exercise}
Misalkan $U \subset \R^n$ terbuka, $p \in U$, dan
$f \colon U \to \R$,
$g \colon U \to \R$,
$h \colon U \to \R$ merupakan fungsi-fungsi sedemikian sehingga
$f(p) = g(p) = h(p)$, $f$ dan $h$ diferensiabel di $p$,
$f'(p) = h'(p)$, dan
\begin{equation*}
f(x) \leq g(x) \leq h(x) \qquad \text{untuk setiap } x \in U .
\end{equation*}
Tunjukkan bahwa $g$ diferensiabel di $p$ dan
$g'(p) = f'(p) = h'(p)$.
\end{exercise}

\begin{exercise}
Buktikan suatu versi \emph{\myindex{teorema nilai rata-rata}}
untuk fungsi beberapa variabel.  Misalkan $U \subset \R^n$
terbuka, $f \colon U \to \R$ diferensiabel, $p,q \in U$, dan ruas garis
$[p,q] \subset U$.  Buktikan bahwa terdapat $x \in [p,q]$ sedemikian sehingga
$\nabla f (x) \cdot (q-p) = f(q)-f(p)$.
\end{exercise}

\begin{exercise}
Definisikan $f \colon \R^2 \to \R$ dengan $f(x,y) \coloneqq \frac{y(x^2+y^2)}{x}$
ketika $x \neq 0$ dan $f(0,y) \coloneqq 0$.  Tunjukkan bahwa untuk setiap $u \in \R^2$
dengan $\snorm{u}=1$, diperoleh $D_u f(0) = 0$, yakni turunan
arah ada dan bernilai nol dalam setiap arah di titik asal, tetapi
fungsi tersebut tidak kontinu di titik asal (bahkan tidak kontinu setiap kali
$x=0$).
\end{exercise}

%%%%%%%%%%%%%%%%%%%%%%%%%%%%%%%%%%%%%%%%%%%%%%%%%%%%%%%%%%%%%%%%%%%%%%%%%%%%%%

\sectionnewpage
\section{Kekontinuan dan turunan}
\label{sec:svthedercont}

%mbxINTROSUBSECTION

\sectionnotes{1--2 lectures}

\subsection{Memberi batas pada turunan}

Mari kita buktikan suatu \myquote{\myindex{teorema nilai rata-rata}} untuk fungsi bernilai vektor.

\begin{lemma} \label{lemma:mvtmv}
Jika $\varphi \colon [a,b] \to \R^n$ diferensiabel pada $(a,b)$ dan
kontinu pada $[a,b]$, maka terdapat $t_0 \in (a,b)$ sedemikian sehingga
\begin{equation*}
\bnorm{\varphi(b)-\varphi(a)} \leq (b-a) \bnorm{\varphi'(t_0)} .
\end{equation*}
\end{lemma}

\begin{proof}
Dengan menerapkan teorema nilai rata-rata pada fungsi bernilai skalar
$t \mapsto \bigl(\varphi(b)-\varphi(a) \bigr) \cdot \varphi(t)$,
di mana tanda titik menyatakan hasil kali titik, kita memperoleh
$t_0 \in (a,b)$ sedemikian sehingga
\begin{equation*}
\begin{split}
\bnorm{\varphi(b)-\varphi(a)}^2
& =
\bigl( \varphi(b)-\varphi(a) \bigr)
\cdot
\bigl( \varphi(b)-\varphi(a) \bigr)
\\
& =
\bigl(\varphi(b)-\varphi(a) \bigr) \cdot \varphi(b) - 
\bigl(\varphi(b)-\varphi(a) \bigr) \cdot \varphi(a)
\\
& = 
(b-a)
\bigl(\varphi(b)-\varphi(a) \bigr) \cdot \varphi'(t_0) ,
\end{split}
\end{equation*}
di sini kita memperlakukan $\varphi'$ sebagai vektor di $\R^n$ melalui penyalahgunaan
notasi yang telah kita sebutkan dalam bagian sebelumnya.
Jika $\varphi'(t)$ kita pandang sebagai vektor, maka menurut
\exerciseref{exercise:normonedim},
$\bnorm{\varphi'(t)}_{L(\R,\R^n)} = \bnorm{\varphi'(t)}_{\R^n}$.
Dengan kata lain, norma Euklides vektor tersebut sama dengan norma operator
$\varphi'(t)$.

Menurut ketaksamaan Cauchy--Schwarz,
\begin{equation*}
\bnorm{\varphi(b)-\varphi(a)}^2
=
(b-a)\bigl(\varphi(b)-\varphi(a) \bigr) \cdot \varphi'(t_0)
\leq
(b-a)
\bnorm{\varphi(b)-\varphi(a)} \, \bnorm{\varphi'(t_0)} . \qedhere
\end{equation*}
\end{proof}

Ingat bahwa suatu himpunan $U$ disebut konveks
jika setiap kali $p,q \in U$, ruas garis dari
$p$ ke $q$ terletak di dalam $U$.

\begin{prop} \label{mv:prop:convexlip}
Misalkan $U \subset \R^n$ adalah himpunan terbuka yang konveks, $f \colon U \to \R^m$
adalah fungsi diferensiabel, dan $M$ memenuhi
\begin{equation*}
\bnorm{f'(p)} \leq M
\qquad \text{untuk setiap } p \in U.
\end{equation*}
Maka $f$ bersifat Lipschitz dengan konstanta $M$, yaitu
\begin{equation*}
\bnorm{f(p)-f(q)} \leq M \snorm{p-q}
\qquad
\text{untuk setiap } p,q \in U.
\end{equation*}
\end{prop}

\begin{proof}
Tetapkan $p$ dan $q$ di $U$, lalu perhatikan bahwa
$(1-t)p+tq \in U$ untuk setiap $t \in [0,1]$
berdasarkan kekonveksan.
Selanjutnya,
\begin{equation*}
\frac{d}{dt} \Bigl[f\bigl((1-t)p+tq\bigr)\Bigr]
=
f'\bigl((1-t)p+tq\bigr) (q-p) .
\end{equation*}
Menurut \lemmaref{lemma:mvtmv}, terdapat
$t_0 \in (0,1)$ sedemikian sehingga
\begin{equation*}
\begin{split}
\bnorm{f(p)-f(q)} & \leq
\norm{\frac{d}{dt} \Big|_{t=t_0} \Bigl[ f\bigl((1-t)p+tq\bigr) \Bigr] }
\\
& \leq
\norm{f'\bigl((1-t_0)p+t_0q\bigr)} \, \snorm{q-p} \leq
M \snorm{q-p} . \qedhere
\end{split}
\end{equation*}
\end{proof}

\begin{example}
Jika $U$ tidak konveks, proposisi tersebut tidak berlaku. Tinjau
himpunan
\begin{equation*}
U \coloneqq \bigl\{ (x,y) : 0.5 < x^2+y^2 < 2 \bigr\}
\setminus \bigl\{ (x,0) : x < 0 \bigr\} .
\end{equation*}
Untuk $(x,y) \in U$,
misalkan $f(x,y)$ adalah sudut yang dibentuk oleh garis dari titik asal ke $(x,y)$
dengan sumbu $x$ positif.  Kita bahkan memiliki rumus untuk $f$:
\begin{equation*}
f(x,y) = 2 \operatorname{arctan}\left( \frac{y}{x+\sqrt{x^2+y^2}}\right) .
\end{equation*}
Bayangkan tangga spiral dengan ruang kosong di tengahnya.  Lihat
\figureref{mv:fignonlip}.

\begin{myfigureht}
\myincludepdft{nonlip_full}{%
Di sebelah kiri terdapat diagram pada bidang xy berupa daerah berbentuk gelang yang diarsir
(daerah di antara dua lingkaran), berpusat di titik asal, dan memiliki sayatan
sepanjang sumbu x negatif.
Sebuah titik (x,y) diberi label pada daerah gelang dan sebuah garis ditarik dari titik ini
ke titik asal; sudut dari sumbu x positif
ke garis ini diberi label theta sama dengan f dari x, y.
Di sebelah kanan terdapat grafik fungsi ini berupa spiral menaik
yang bermula di bawah bidang xy ketika kita mulai dari sayatan
dan menyusuri daerah gelang berlawanan arah jarum jam, lalu mencapai bidang xy
ketika kita tiba di sumbu x positif, dan berada di atas bidang xy
ketika kita kembali ke sayatan dari sisi lain.}
\caption{Fungsi non-Lipschitz dengan turunan yang terbatas
secara seragam.\label{mv:fignonlip}}
\end{myfigureht}

Fungsi tersebut diferensiabel,
dan turunannya terbatas pada $U$; hal ini tidak sulit dilihat.  Sekarang
perhatikan
apa yang terjadi di dekat sayatan pada anulus di sepanjang sumbu $x$ negatif.
Ketika kita mendekati sayatan ini dari sisi $y$ positif, $f(x,y)$ mendekati $\pi$.
Dari sisi $y$ negatif, $f(x,y)$ mendekati $-\pi$.
Jadi, untuk $\epsilon > 0$ yang kecil, $\babs{f(-1,\epsilon)-f(-1,-\epsilon)}$
mendekati $2\pi$, sedangkan $\bnorm{(-1,\epsilon)-(-1,-\epsilon)} = 2\epsilon$,
yang dapat dibuat sekecil apa pun.  Kesimpulan proposisi tersebut tidak
berlaku untuk $U$ yang tidak konveks ini.
\end{example}

Mari kita selesaikan persamaan diferensial $f' = 0$.

\begin{cor}
Jika $U \subset \R^n$ terbuka dan terhubung, $f \colon U \to \R^m$
diferensiabel,
dan $f'(x) = 0$ untuk setiap $x \in U$, maka $f$ konstan.
\end{cor}

\begin{proof}
Untuk setiap $x \in U$, terdapat bola $B(x,\delta) \subset U$.  Bola
$B(x,\delta)$ bersifat konveks.  Karena
$\bnorm{f'(y)} \leq 0$ untuk setiap $y \in B(x,\delta)$, menurut proposisi tersebut,
$\bnorm{f(x)-f(y)} \leq 0 \snorm{x-y} = 0$.  Jadi, $f(x) = f(y)$ untuk setiap $y \in
B(x,\delta)$.
Oleh karena itu, $f^{-1}(c)$ terbuka untuk setiap $c \in \R^m$.

Andaikan $c_0 \in \R^m$ sedemikian sehingga
$f^{-1}(c_0) \neq \emptyset$. 
Karena $f$ juga kontinu,
kedua himpunan
\begin{equation*}
U' = f^{-1}(c_0), \qquad U'' = f^{-1}\bigl(\R^m\setminus\{c_0\}\bigr)
\end{equation*}
terbuka dan saling lepas, serta $U = U' \cup U''$.  Karena $U'$ tidak kosong
dan $U$ terhubung, maka
$U'' = \emptyset$.  Jadi, $f(x) = c_0$ untuk setiap $x \in U$.
\end{proof}

\subsection{Fungsi yang diferensiabel secara kontinu}

\begin{defn}
Misalkan $U \subset \R^n$ terbuka.
Kita mengatakan bahwa $f \colon U \to \R^m$
\emph{\myindex{diferensiabel secara kontinu}},
atau berkelas $C^1(U)$\glsadd{not:C1},
jika $f$ diferensiabel dan $f' \colon U \to L(\R^n,\R^m)$
kontinu.
\end{defn}

\begin{prop} \label{mv:prop:contdiffpartials}
Misalkan $U \subset \R^n$ terbuka dan
$f \colon U \to \R^m$.  Fungsi
$f$ diferensiabel secara kontinu jika dan hanya jika 
turunan parsial $\frac{\partial f_k}{\partial x_j}$
ada dan kontinu untuk setiap $k$ dan $j$.
\end{prop}

Tanpa kekontinuan turunan parsial, pernyataan tersebut tidak berlaku.  Keberadaan
turunan parsial saja tidak berarti bahwa $f$ diferensiabel;
bahkan, $f$ mungkin tidak kontinu.  Lihat latihan
untuk bagian sebelumnya dan juga untuk bagian ini.

\begin{proof}
Kita telah membuktikan bahwa jika $f$ diferensiabel, maka
turunan parsialnya ada.  Turunan parsial
merupakan entri-entri matriks yang merepresentasikan $f'(x)$.  Jika
$f' \colon U \to L(\R^n,\R^m)$ kontinu, maka entri-entri tersebut
kontinu, sehingga turunan parsialnya kontinu.

Untuk membuktikan arah sebaliknya,
andaikan turunan parsialnya ada dan kontinu.
Tetapkan $p \in U$.  Jika kita menunjukkan bahwa $f'(p)$ ada, pembuktian selesai, sebab
entri-entri matriks yang merepresentasikan $f'(p)$ adalah turunan
parsial, dan jika entri-entri tersebut merupakan fungsi kontinu,
fungsi bernilai matriks $f'$ juga kontinu.

Kita menggunakan induksi terhadap dimensi.  Pertama,
kesimpulan tersebut benar ketika $n=1$
(latihan; perhatikan bahwa $f$ bernilai vektor).
Dalam hal ini, $f'(p)$ pada dasarnya merupakan turunan yang dibahas dalam
\volIref{\chapterref*{vI-der:chapter}}{\chapterref{der:chapter}}.
Andaikan kesimpulan tersebut benar untuk $\R^{n-1}$.
Artinya,
jika kita membatasi fungsi pada $n-1$ variabel pertama, fungsi itu diferensiabel.
Ketika mengambil turunan parsial terhadap $x_1$ hingga $x_{n-1}$,
tidak ada bedanya apakah kita meninjau $f$ atau pembatasan $f$ pada himpunan tempat
$x_n$ bernilai tetap.
Dalam uraian berikut, dengan sedikit penyalahgunaan notasi,
kita memandang $\R^{n-1}$ sebagai subhimpunan $\R^n$, yaitu himpunan dalam $\R^n$ tempat $x_n = 0$.
Dengan kata lain, kita mengidentifikasi vektor $(x_1,x_2,\ldots,x_{n-1})$ dan
$(x_1,x_2,\ldots,x_{n-1},0)$.

Dengan $p \in U$ yang telah ditetapkan, misalkan
\begin{equation*}
A \coloneqq 
\begin{bmatrix}
\frac{\partial f_1}{\partial x_1}(p)
& \ldots &
\frac{\partial f_1}{\partial x_n}(p)
\\
\vdots & \ddots & \vdots
\\
\frac{\partial f_m}{\partial x_1}(p)
& \ldots &
\frac{\partial f_m}{\partial x_n}(p)
\end{bmatrix} ,
%mbxSTARTIGNORE
\qquad
%mbxENDIGNORE
%mbxlatex \quad
A' \coloneqq 
\begin{bmatrix}
\frac{\partial f_1}{\partial x_1}(p)
& \ldots &
\frac{\partial f_1}{\partial x_{n-1}}(p)
\\
\vdots & \ddots & \vdots
\\
\frac{\partial f_m}{\partial x_1}(p)
& \ldots &
\frac{\partial f_m}{\partial x_{n-1}}(p)
\end{bmatrix} ,
%mbxSTARTIGNORE
\qquad
%mbxENDIGNORE
%mbxlatex \quad
v \coloneqq 
\begin{bmatrix}
\frac{\partial f_1}{\partial x_n}(p)
\\
\vdots
\\
\frac{\partial f_m}{\partial x_n}(p)
\end{bmatrix} .
\end{equation*}
Ambil sembarang $\epsilon > 0$.  Karena $U$ terbuka, hipotesis induksi memberikan
suatu $\delta > 0$; dengan memperkecilnya jika perlu, kita dapat memastikan $B(p,\delta) \subset U$ dan
untuk setiap $h' \in \R^{n-1}$ dengan $\snorm{h'} < \delta$, berlaku
\begin{equation*}
\bnorm{f(p+h') - f(p) - A' h'} \leq \epsilon \snorm{h'} .
\end{equation*}
Dengan memperkecil $\delta$ jika perlu, kekontinuan turunan parsial memastikan
bahwa
\begin{equation*}
\abs{\frac{\partial f_k}{\partial x_n}(p+h)
      - \frac{\partial f_k}{\partial x_n}(p)} < \epsilon
\end{equation*}
untuk setiap $k$ dan setiap $h \in \R^n$ dengan $\snorm{h} < \delta$.

Andaikan $h = h' + t e_n$ adalah vektor dalam $\R^n$, dengan $h' \in \R^{n-1}$,
$t \in \R$, sedemikian sehingga
$\snorm{h} < \delta$.  Maka $\snorm{h'} \leq \snorm{h} < \delta$.
Perhatikan bahwa $Ah = A' h' + tv$.
\begin{equation*}
\begin{split}
%mbxlatex &
\bnorm{f(p+h) - f(p) - Ah}
%mbxSTARTIGNORE
&
%mbxENDIGNORE
%mbxlatex \\
%mbxlatex & \qquad \qquad
= \bnorm{f(p+h' + t e_n) - f(p+h') - tv + f(p+h') - f(p) - A' h'}
\\
&
%mbxlatex \qquad \qquad
\leq \bnorm{f(p+h' + t e_n) - f(p+h') -tv} + \bnorm{f(p+h') - f(p) -
A' h'}
\\
&
%mbxlatex \qquad \qquad
\leq \bnorm{f(p+h' + t e_n) - f(p+h') -tv} + \epsilon \snorm{h'} .
\end{split}
\end{equation*}
Karena semua turunan parsial ada, untuk setiap $k$ teorema nilai rata-rata memberikan
suatu $\theta_k \in [0,t]$ (atau $[t,0]$ jika $t < 0$), sedemikian sehingga
\begin{equation*}
f_k(p+h' + t e_n) - f_k(p+h') =
t \frac{\partial f_k}{\partial x_n}(p+h'+\theta_k e_n).
\end{equation*}
Kita mempunyai $\snorm{h'+\theta_k e_n} \leq \snorm{h} < \delta$,
sehingga penaksiran dapat dituntaskan sebagai berikut:
\begin{equation*}
\begin{split}
\bnorm{f(p+h) - f(p) - Ah}
& \leq \bnorm{f(p+h' + t e_n) - f(p+h') -tv} + \epsilon \snorm{h'}
\\
& \leq \sqrt{\sum_{k=1}^m {\left(t\frac{\partial f_k}{\partial x_n}(p+h'+\theta_k e_n) -
t \frac{\partial f_k}{\partial x_n}(p)\right)}^2} + \epsilon \snorm{h'}
\\
& \leq \sqrt{m}\, \epsilon \sabs{t} + \epsilon \snorm{h'}
\\
& \leq (\sqrt{m}+1)\epsilon \snorm{h} . \qedhere
\end{split}
\end{equation*}
\end{proof}

Penerapan yang umum adalah membuktikan bahwa suatu fungsi
diferensiabel.  Sebagai contoh, kita dapat menunjukkan bahwa semua polinom
diferensiabel, bahkan diferensiabel secara kontinu,
dengan menghitung turunan parsialnya.

\begin{cor}
Suatu polinom $p \colon \R^n \to \R$ dalam beberapa variabel,
\begin{equation*}
p(x_1,x_2,\ldots,x_n)
=
\sum_{0 \leq j_1+j_2+\cdots+j_n \leq d}
c_{j_1,j_2,\ldots,j_n}
\,
x_1^{j_1}
x_2^{j_2}
\cdots
x_n^{j_n}
\end{equation*}
bersifat diferensiabel secara kontinu.
\end{cor}

\begin{proof}
Tinjau turunan parsial $p$ terhadap variabel $x_n$.
Tuliskan $p$ sebagai
\begin{equation*}
p(x) = \sum_{j=0}^d p_j(x_1,\ldots,x_{n-1}) \, x_n^j ,
\end{equation*}
di mana $p_j$ adalah polinom dalam satu variabel lebih sedikit.
Maka
\begin{equation*}
\frac{\partial p}{\partial x_n}(x)
= \sum_{j=1}^d p_j(x_1,\ldots,x_{n-1}) \, j x_n^{j-1} ,
\end{equation*}
yang kembali merupakan polinom.
Jadi, turunan parsial polinom ada dan kembali berupa polinom.
Berdasarkan kekontinuan operasi aljabar, polinom merupakan fungsi kontinu.
Oleh karena itu, $p$ diferensiabel secara kontinu.
\end{proof}

\subsection{Latihan}

\begin{exercise}
Definisikan $f \colon \R^2 \to \R$ sebagai
\begin{equation*}
f(x,y) \coloneqq
\begin{cases}
(x^2+y^2)\sin\bigl({(x^2+y^2)}^{-1}\bigr) & \text{jika } (x,y) \neq (0,0), \\
0                                         & \text{jika } (x,y) = (0,0).
\end{cases}
\end{equation*}
Tunjukkan bahwa $f$ diferensiabel di titik asal, tetapi tidak 
diferensiabel secara kontinu.
\\
Catatan: Anda boleh menggunakan pengetahuan tentang sinus dan kosinus dari kalkulus.
\end{exercise}

\begin{exercise}
Misalkan $f \colon \R^2 \to \R$ adalah fungsi dalam
\exerciseref{exercise:noncontpartialsexist}, yaitu
\begin{equation*}
f(x,y)
\coloneqq
\begin{cases}
\frac{xy}{x^2+y^2} & \text{jika } (x,y) \neq (0,0), \\
0                  & \text{jika } (x,y) = (0,0).
\end{cases}
\end{equation*}
Hitung turunan parsial 
$\frac{\partial f}{\partial x}$ dan
$\frac{\partial f}{\partial y}$ pada setiap titik, lalu tunjukkan bahwa keduanya bukan
fungsi kontinu.
\end{exercise}

\begin{exercise}
Misalkan $B(0,1) \subset \R^2$ adalah bola satuan, yaitu himpunan yang diberikan oleh
$x^2 + y^2 < 1$.
Andaikan $f \colon B(0,1) \to \R$ adalah fungsi diferensiabel
sedemikian sehingga $\babs{f(0,0)} \leq 1$,
dan 
$\babs{\frac{\partial f}{\partial x}} \leq 1$ dan
$\babs{\frac{\partial f}{\partial y}} \leq 1$ di setiap
titik dalam $B(0,1)$.
\begin{enumerate}[a)]
\item
Tentukan $M \in \R$ sedemikian sehingga $\bnorm{f'(x,y)} \leq M$
untuk setiap $(x,y) \in
B(0,1)$.
\item
Tentukan $B \in \R$ sedemikian sehingga
$\babs{f(x,y)} \leq B$
untuk setiap $(x,y) \in
B(0,1)$.
\end{enumerate}
\end{exercise}

\begin{exercise}
Definisikan $\varphi \colon [0,2\pi] \to \R^2$ dengan $\varphi(t) =
\bigl(\sin(t),\cos(t)\bigr)$.  Hitung $\varphi'(t)$ untuk setiap $t$.  Hitung
$\bnorm{\varphi'(t)}$ untuk setiap $t$.  Perhatikan bahwa $\varphi'(t)$ tidak pernah sama dengan vektor nol,
namun $\varphi(0) = \varphi(2\pi)$.  Oleh karena itu, teorema Rolle tidak berlaku
untuk fungsi bernilai vektor.
\end{exercise}

\begin{exercise}
Misalkan $f \colon \R^2 \to \R$ adalah fungsi sedemikian sehingga
$\frac{\partial f}{\partial x}$ dan
$\frac{\partial f}{\partial y}$ ada pada setiap titik dan terdapat $M
\in \R$
sedemikian sehingga 
$\babs{\frac{\partial f}{\partial x}} \leq M$ dan
$\babs{\frac{\partial f}{\partial y}} \leq M$ pada setiap titik.  Tunjukkan bahwa $f$
kontinu.
\end{exercise}

\begin{samepage}
\begin{exercise}
Misalkan $f \colon \R^2 \to \R$ adalah fungsi dan
$M \in \R$, sedemikian sehingga
untuk setiap $(x,y) \in \R^2$ dengan $x^2+y^2=1$, fungsi $g(t) \coloneqq f(xt,yt)$
diferensiabel
dan $\babs{g'(t)} \leq M$ untuk setiap $t$.
\begin{enumerate}[a)]
\item
Tunjukkan bahwa $f$ kontinu di $(0,0)$.
\item
Berikan contoh fungsi $f$ semacam itu yang diskontinu di setiap titik lain dalam
$\R^2$.\\
Petunjuk: Ingat kembali cara kita mengonstruksi fungsi yang tidak kontinu di setiap titik pada $[0,1]$.
\end{enumerate}
\end{exercise}
\end{samepage}

\begin{exercise}
Andaikan $r \colon \R^n \setminus X \to \R$ adalah fungsi rasional, yaitu
$p \colon \R^n \to \R$ dan
$q \colon \R^n \to \R$ adalah polinom,
$q$ tidak identik nol,
$X = q^{-1}(0)$, dan
$r = \frac{p}{q}$.
Tunjukkan bahwa $r$ diferensiabel secara kontinu.
\end{exercise}

\begin{exercise}
Andaikan $f \colon \R^n \to \R$ dan $h \colon \R^n \to \R$ adalah dua 
fungsi diferensiabel sedemikian sehingga $f'(x) = h'(x)$ untuk setiap $x \in \R^n$.
Buktikan bahwa
jika $f(0) = h(0)$, maka $f(x) = h(x)$ untuk setiap $x \in \R^n$.
\end{exercise}

\begin{exercise}
Buktikan kasus dasar
pada \propref{mv:prop:contdiffpartials}.  Artinya, buktikan bahwa
jika $n=1$ dan 
\myquote{turunan parsial ada dan kontinu}, maka fungsi tersebut
diferensiabel secara kontinu.  Perhatikan bahwa $f$ bernilai vektor.
\end{exercise}

\begin{exercise}
Andaikan $U \subset \R^n$ terbuka, $f \colon U \to \R^m$
diferensiabel, terdapat $M$ sedemikian sehingga $\bnorm{f'(p)} \leq M$
untuk setiap $p \in U$, dan $K \subset U$ adalah himpunan kompak.  Buktikan bahwa
terdapat $M'$ (dengan $M' \geq M$) sedemikian sehingga
untuk setiap $p,q \in K$ berlaku $\bnorm{f(p)-f(q)} \leq M' \snorm{p-q}$.
Bandingkan dengan \propref{mv:prop:convexlip}.
\end{exercise}


%%%%%%%%%%%%%%%%%%%%%%%%%%%%%%%%%%%%%%%%%%%%%%%%%%%%%%%%%%%%%%%%%%%%%%%%%%%%%%

\sectionnewpage
\section{Teorema fungsi invers dan fungsi implisit}
\label{sec:svinvfuncthm}

%mbxINTROSUBSECTION

\sectionnotes{2--3 lectures}

Secara intuitif, jika suatu fungsi diferensiabel secara kontinu, maka secara lokal fungsi itu
\myquote{berperilaku seperti} turunannya (yang merupakan fungsi linear).
Gagasan teorema fungsi invers adalah bahwa jika suatu fungsi
diferensiabel secara kontinu dan turunannya invertibel, maka fungsi tersebut
invertibel (secara lokal).


\begin{thm}[Teorema fungsi invers]\index{teorema fungsi invers}
\label{thm:inverse}
Misalkan $U \subset \R^n$ himpunan terbuka dan
$f \colon U \to \R^n$ fungsi yang diferensiabel secara kontinu.
Andaikan $p \in U$ dan $f'(p)$ invertibel
(yaitu, $J_f(p) \neq 0$).
Maka terdapat himpunan terbuka $V, W \subset \R^n$ sedemikian sehingga
$p \in V \subset U$, $f(V) = W$, dan $f|_V$ injektif.  
Dengan demikian, terdapat fungsi $g \colon W \to V$ sedemikian sehingga
$g(y) \coloneqq (f|_V)^{-1}(y)$.
Selain itu, $g$ diferensiabel secara kontinu
dan 
\begin{equation*}
g'(y) = {\bigl(f'(x)\bigr)}^{-1}, \qquad \text{untuk setiap } x \in V, y = f(x).
\end{equation*}
\end{thm}
Lihat \figureref{fig:inversefuncRn}.

\begin{myfigureht}
\myincludepdft{inversefuncRn}{%
Diagram dua himpunan dan pemetaan di antara keduanya.
Di sebelah kiri terdapat himpunan U yang diarsir dengan batas putus-putus dan subhimpunan V yang lebih kecil,
diarsir tipis dengan batas putus-putus.  Titik p dan x ditandai di dalam V.
Di sebelah kanan terdapat himpunan berarsir tipis berlabel W sama dengan f dari V, dengan dua
titik: satu ditandai f dari p dan satu lagi ditandai y.  Dua panah berlabel f mengarah dari kiri
ke kanan; satu dari titik p ke titik f dari p
dan satu lagi dari titik x ke titik y.  Dua panah berlabel g
mengarah sebaliknya di antara titik-titik yang sama.}
\caption{Skema teorema fungsi invers dalam $\R^n$.\label{fig:inversefuncRn}}
\end{myfigureht}

Untuk membuktikan teorema tersebut, kita menggunakan prinsip pemetaan
kontraksi dari
\volIref{\chapterref*{vI-ms:chapter}}{\chapterref{ms:chapter}},
yang telah kita gunakan
untuk membuktikan teorema Picard.
Ingat bahwa pemetaan $f \colon X \to Y$ di antara ruang metrik
$(X,d_X)$ dan $(Y,d_Y)$ merupakan kontraksi 
jika terdapat $0 \leq k < 1$ sedemikian sehingga
\begin{equation*}
d_Y\bigl(f(p),f(q)\bigr) \leq k \, d_X(p,q)
\qquad \text{untuk setiap } p,q \in X.
\end{equation*}
Prinsip pemetaan kontraksi menyatakan bahwa jika $f \colon X \to X$
merupakan kontraksi dan $X$ ruang metrik lengkap,
maka terdapat tepat satu titik tetap, yaitu
terdapat tepat satu $x \in X$ sedemikian sehingga $f(x) = x$.

\begin{proof}
Tuliskan $A = f'(p)$.  Karena $f'$ kontinu dan domainnya terbuka, kita dapat memilih bola terbuka
$V \subset U$ yang berpusat di $p$ sedemikian sehingga
\begin{equation*}
\bnorm{A-f'(x)} < \frac{1}{2\snorm{A^{-1}}}
\qquad \text{untuk setiap } x \in V.
\end{equation*}
Akibatnya, turunan $f'(x)$ invertibel untuk setiap $x \in V$
menurut \propref{prop:finitedimpropinv}.

Untuk $y \in \R^n$, definisikan $\varphi_y \colon V \to \R^n$ dengan
\begin{equation*}
\varphi_y (x) \coloneqq x + A^{-1}\bigl(y-f(x)\bigr) .
\end{equation*}
Karena $A^{-1}$ injektif,
$\varphi_y(x) = x$ ($x$ adalah titik tetap) jika dan hanya jika
$y-f(x) = 0$, atau dengan kata lain $f(x)=y$.  Dengan menggunakan aturan rantai, kita memperoleh
\begin{equation*}
\varphi_y'(x) = I - A^{-1} f'(x) = A^{-1} \bigl( A-f'(x) \bigr) .
\end{equation*}
Jadi, untuk $x \in V$, kita mempunyai
\begin{equation*}
\bnorm{\varphi_y'(x)} \leq \snorm{A^{-1}} \, \bnorm{A-f'(x)} < \nicefrac{1}{2} .
\end{equation*}
Karena $V$ sebuah bola, himpunan tersebut konveks.  Oleh karena itu,
\begin{equation*}
\bnorm{\varphi_y(x_1)-\varphi_y(x_2)} \leq \frac{1}{2} \snorm{x_1-x_2} 
\qquad
\text{untuk setiap } x_1,x_2 \in V.
\end{equation*}
Dengan kata lain, $\varphi_y$ adalah kontraksi yang didefinisikan pada $V$, meskipun sejauh ini
kita belum mengetahui jangkauan $\varphi_y$.  Kita belum dapat
menerapkan teorema titik
tetap, tetapi kita dapat mengatakan bahwa $\varphi_y$ 
memiliki paling banyak satu titik tetap dalam $V$:
Jika $\varphi_y(x_1) = x_1$ dan
$\varphi_y(x_2) = x_2$, maka
$\snorm{x_1-x_2} = \bnorm{\varphi_y(x_1)-\varphi_y(x_2)} \leq
\frac{1}{2} \snorm{x_1-x_2}$, sehingga $x_1 = x_2$.
Artinya, terdapat paling banyak satu $x \in V$
sedemikian sehingga $f(x) = y$, sehingga $f|_V$ injektif.

Misalkan $W \coloneqq f(V)$ dan misalkan $g \colon W \to V$ adalah invers dari $f|_V$.
Kita perlu menunjukkan bahwa $W$ terbuka.  Ambil $y_0 \in W$.
Terdapat tepat satu $x_0 \in V$ sedemikian sehingga $f(x_0) = y_0$.
Ambil $r > 0$ yang cukup kecil sehingga bola tertutup $C(x_0,r) \subset V$
($r > 0$ semacam itu ada karena $V$ terbuka).

Andaikan $y$ memenuhi
\begin{equation*}
\snorm{y-y_0} <
\frac{r}{2\snorm{A^{-1}}} .
\end{equation*}
Jika kita menunjukkan bahwa $y \in W$, maka kita telah menunjukkan bahwa $W$ terbuka.
Jika $x_1 \in
C(x_0,r)$, maka
\begin{equation*}
\begin{split}
\bnorm{\varphi_y(x_1)-x_0}
& \leq
\bnorm{\varphi_y(x_1)-\varphi_y(x_0)} +
\bnorm{\varphi_y(x_0)-x_0} \\
& \leq
\frac{1}{2}\snorm{x_1-x_0} +
\bnorm{A^{-1}(y-y_0)} \\
& \leq
\frac{1}{2}r +
\snorm{A^{-1}} \, \snorm{y-y_0} \\
& <
\frac{1}{2}r +
\snorm{A^{-1}}
\frac{r}{2\snorm{A^{-1}}} = r .
\end{split}
\end{equation*}
Jadi, $\varphi_y$ memetakan $C(x_0,r)$ ke dalam $B(x_0,r) \subset C(x_0,r)$.  Pemetaan ini
merupakan kontraksi pada $C(x_0,r)$, dan $C(x_0,r)$ lengkap (karena setiap subhimpunan tertutup dari
$\R^n$ lengkap).
Terapkan prinsip pemetaan kontraksi untuk memperoleh titik tetap $x$,
yaitu\ $\varphi_y(x) = x$.  Artinya, $f(x) = y$, dan $y \in
f\bigl(C(x_0,r)\bigr) \subset f(V) = W$.  Oleh karena itu, $W$ terbuka.

Selanjutnya kita perlu menunjukkan bahwa $g$ diferensiabel secara kontinu dan menghitung
turunannya.  Pertama, mari kita tunjukkan bahwa fungsi tersebut diferensiabel.
Misalkan $y \in W$ dan $k \in \R^n$, $k \neq 0$, sedemikian sehingga $y+k \in W$.
Karena $f|_V$ merupakan pemetaan injektif dan surjektif dari $V$ ke $W$,
terdapat tepat satu $x \in V$ dan tepat satu
$h \in \R^n$, $h \neq 0$ serta $x+h \in V$, sedemikian sehingga
$f(x) = y$ dan $f(x+h) = y+k$.
Dengan kata lain, $g(y) = x$ dan $g(y+k) = x+h$.  Lihat
\figureref{fig:inversefuncRn2}.
\begin{myfigureht}
\myincludepdft{inversefuncRn2}{%
Diagram dua himpunan dan pemetaan di antara keduanya.
Di sebelah kiri terdapat himpunan V yang diarsir dengan batas putus-putus dan
titik x tambah h serta x yang ditandai.
Di sebelah kanan terdapat himpunan W yang diarsir dengan titik y tambah k
serta y yang ditandai.
Dua panah berlabel f mengarah dari kiri
ke kanan; satu dari titik x tambah h ke titik y tambah k
dan satu lagi dari titik x ke titik y.  Dua panah berlabel g
mengarah sebaliknya di antara titik-titik yang sama.}
\caption{Membuktikan bahwa $g$ diferensiabel.\label{fig:inversefuncRn2}}
\end{myfigureht}

Kita masih dapat
memperoleh informasi dari fakta bahwa $\varphi_y$ merupakan kontraksi.
\begin{equation*}
\varphi_y(x+h)-\varphi_y(x) = h + A^{-1} \bigl( f(x)-f(x+h) \bigr) = h - A^{-1} k .
\end{equation*}
Maka
\begin{equation*}
\snorm{h-A^{-1}k} = \bnorm{\varphi_y(x+h)-\varphi_y(x)} \leq
\frac{1}{2}\snorm{x+h-x} = \frac{\snorm{h}}{2}.
\end{equation*}
Menurut ketaksamaan segitiga terbalik, $\snorm{h} - \snorm{A^{-1}k} \leq
\frac{1}{2}\snorm{h}$.
Jadi,
\begin{equation*}
\snorm{h} \leq 2 \snorm{A^{-1}k} \leq 2 \snorm{A^{-1}} \, \snorm{k}.
\end{equation*}
Secara khusus, ketika $k$ menuju 0, $h$ juga menuju 0.

Karena $x \in V$, kita mengetahui bahwa $f'(x)$ invertibel.
Misalkan $B \coloneqq \bigl(f'(x)\bigr)^{-1}$, yang kita duga merupakan turunan
$g$ di $y$.  Maka
\begin{equation*}
\begin{split}
\frac{\bnorm{g(y+k)-g(y)-Bk}}{\snorm{k}}
& =
\frac{\snorm{h-Bk}}{\snorm{k}}
\\
& =
\frac{\bnorm{h-B\bigl(f(x+h)-f(x)\bigr)}}{\snorm{k}}
\\
& =
\frac{\bnorm{B\bigl(f(x+h)-f(x)-f'(x)h\bigr)}}{\snorm{k}}
\\
& \leq
\snorm{B}
\frac{\snorm{h}}{\snorm{k}}\,
\frac{\bnorm{f(x+h)-f(x)-f'(x)h}}{\snorm{h}}
\\
& \leq
2\snorm{B} \, \snorm{A^{-1}}
\frac{\bnorm{f(x+h)-f(x)-f'(x)h}}{\snorm{h}} .
\end{split}
\end{equation*}
Ketika $k$ menuju 0, $h$ juga menuju 0.  Jadi, ruas kanan menuju 0 karena $f$
diferensiabel, sehingga
ruas kiri juga menuju 0.  Dengan demikian,
$B$ tepat merupakan $g'(y)$ yang kita cari.

Kita telah menunjukkan bahwa $g$ diferensiabel; mari kita tunjukkan bahwa fungsi tersebut berkelas $C^1(W)$.
Fungsi $g \colon W \to V$ kontinu (karena diferensiabel),
$f'$ adalah fungsi kontinu dari $V$
ke $L(\R^n)$, dan $X \mapsto X^{-1}$ adalah fungsi kontinu pada
himpunan operator invertibel.
Karena
$g'(y) = {\bigl( f'\bigl(g(y)\bigr)\bigr)}^{-1}$ merupakan komposisi
ketiga fungsi
kontinu tersebut, maka komposisi tersebut kontinu.
\end{proof}

\begin{cor} \label{cor:IVTopenmap}
Andaikan $U \subset \R^n$ terbuka dan $f \colon U \to \R^n$ adalah pemetaan yang
diferensiabel secara kontinu sedemikian sehingga $f'(x)$ invertibel untuk setiap $x \in U$.  Maka
untuk setiap himpunan terbuka $V \subset U$, himpunan $f(V)$ terbuka ($f$ disebut
\emph{\myindex{pemetaan terbuka}}).
\end{cor}

\begin{proof}
Tanpa mengurangi keumuman, andaikan $U=V$.
Untuk setiap $y \in f(V)$, pilih $x \in f^{-1}(y)$ (mungkin terdapat lebih
dari satu titik semacam itu); lalu, berdasarkan teorema fungsi invers, terdapat
lingkungan $x$ dalam $V$ yang citranya merupakan suatu lingkungan dari $y$.  Dengan demikian,
$f(V)$ terbuka.
\end{proof}

\begin{example}
Teorema tersebut beserta korolarinya tidak berlaku jika $f'(x)$ tidak invertibel untuk
suatu~$x$.
Pemetaan $f(x,y) \coloneqq (x,xy)$ memetakan $\R^2$ pada himpunan
$\R^2 \setminus \bigl\{ (0,y) : y \neq 0 \bigr\}$, yang tidak terbuka maupun tertutup.
Faktanya, $f^{-1}(0,0) = \bigl\{ (0,y) : y \in \R \bigr\}$.  Perilaku buruk ini
hanya terjadi pada sumbu $y$; pada setiap titik lain, fungsi tersebut invertibel
secara lokal.  Jika kita menghindari sumbu $y$, $f$ bahkan injektif.
\end{example}

\begin{example}
Fakta bahwa $f'(x)$ invertibel di setiap titik tidak
berarti bahwa $f$
injektif.  Fungsi tersebut injektif \myquote{secara lokal}, tetapi belum tentu
\myquote{secara global.}
Tinjau $f \colon \R^2 \setminus \bigl\{ (0,0) \bigr\} \to
\R^2 \setminus \bigl\{ (0,0) \bigr\}$ yang didefinisikan
oleh $f(x,y) \coloneqq (x^2-y^2,2xy)$.
Pembaca dipersilakan memverifikasi pernyataan-pernyataan berikut.
Pemetaan $f$ diferensiabel dan turunannya invertibel.
Di sisi lain, $f$ bersifat 2-ke-1 secara global: Untuk setiap
$(a,b)$ yang bukan titik asal, terdapat tepat dua
solusi bagi $x^2-y^2=a$ dan $2xy=b$ ($f$ juga surjektif).  Perhatikan bahwa
setelah Anda menunjukkan adanya paling sedikit satu solusi,
mengganti $x$ dan $y$ dengan $-x$ dan $-y$ menghasilkan solusi lain.
\end{example}

Keinvertibelan turunan bukan syarat yang perlu,
melainkan hanya syarat cukup, untuk memiliki invers kontinu dan menjadi pemetaan
terbuka.  Fungsi $f(x) \coloneqq x^3$ merupakan pemetaan terbuka dari $\R$
ke $\R$ dan injektif secara global dengan invers kontinu, meskipun
inversnya tidak diferensiabel di $x=0$.

\medskip

Sebagai catatan tambahan, terdapat masalah terkenal yang terkait dan hingga kini belum sepenuhnya terpecahkan,
yang disebut \emph{\myindex{konjektur Jacobian}}.  Jika $F \colon \R^n \to
\R^n$ adalah pemetaan polinomial (setiap komponennya merupakan polinom) dan $J_F$
(determinan Jacobian) adalah konstanta
tak nol, apakah $F$ memiliki invers polinom?
Teorema fungsi invers memberikan invers $C^1$ lokal,
tetapi pertanyaannya adalah apakah kita selalu dapat
menemukan invers polinom global.
Jawabannya afirmatif untuk $n=1$, dan baru-baru ini telah ditemukan contoh tandingan
untuk $n \geq 3$.

\subsection{Teorema fungsi implisit}

Teorema fungsi invers merupakan kasus khusus teorema fungsi
implisit, yang akan kita buktikan selanjutnya.  Agak ironis, kita 
membuktikan teorema fungsi implisit dengan menggunakan teorema fungsi invers.
Salah satu cara menyatakan teorema fungsi invers adalah bahwa
persamaan $x-f(y) = 0$ dapat diselesaikan terhadap $y$ sebagai fungsi dari $x$ jika turunannya
terhadap $y$, yaitu $-f'(y)$, invertibel; hal ini ekuivalen dengan $f'(y)$ invertibel.
Dengan demikian, terdapat (secara lokal) suatu
fungsi $g$ sedemikian sehingga $x-f\bigl(g(x)\bigr) = 0$.

Secara umum, persamaan $f(x,y) = 0$ tidak selalu
dapat diselesaikan terhadap $y$ sebagai fungsi dari $x$.
Sebagai contoh, pada umumnya tidak ada solusi
jika $f(x,y)$ sebenarnya tidak bergantung pada $y$.
Untuk contoh yang lebih menarik, perhatikan bahwa $x^2+y^2-1 = 0$ mendefinisikan lingkaran satuan, dan
kita dapat secara lokal menyelesaikan persamaan tersebut terhadap $y$ sebagai fungsi dari $x$ apabila 1) kita berada di dekat
suatu titik pada lingkaran satuan dan 2)~kita tidak berada di titik
tempat lingkaran memiliki garis singgung vertikal, yaitu tempat
$\frac{\partial f}{\partial y} = 0$.

Kita tetapkan beberapa notasi.  Misalkan $(x,y) \in
\R^{n+m}$ menyatakan blok koordinat $(x_1,\ldots,x_n,y_1,\ldots,y_m)$.  
Dengan demikian, kita dapat menuliskan suatu
pemetaan linear $A \in L(\R^{n+m},\R^m)$ sebagai
$A = [ A_x ~ A_y ]$ sehingga $A(x,y) = A_x x + A_y y$,
dengan $A_x \in L(\R^n,\R^m)$ dan
$A_y \in L(\R^m)$.
Pertama, kita berikan versi linear teorema tersebut.

\begin{prop}
Misalkan $A = [A_x~A_y] \in L(\R^{n+m},\R^m)$ dan andaikan 
$A_y$ invertibel.  Jika $B = - {(A_y)}^{-1} A_x$, maka
\begin{equation*}
0 = A ( x, Bx) = A_x x + A_y Bx .
\end{equation*}
Selain itu, untuk setiap $x \in \R^n$, terdapat tepat satu $y \in \R^m$
yang memenuhi $A(x,y) = 0$, yaitu $y=Bx$.
\end{prop}

Pembuktiannya langsung: Kita menyelesaikan $A_x x +A_y y = 0$ dan memperoleh $y = Bx$.
Pendekatan lain adalah \myquote{melengkapi basis,} yaitu menambahkan
baris-baris pada matriks hingga kita memperoleh matriks invertibel:
Operator di $L(\R^{n+m})$ yang diberikan oleh
$(x,y) \mapsto (x,A_x x + A_y y)$
bersifat invertibel, dan pemetaan $B$
dapat diperoleh langsung dari inversnya.
Mari kita tunjukkan bahwa hal yang sama dapat dilakukan untuk fungsi berkelas $C^1$.

\begin{thm}[Teorema fungsi implisit]\index{teorema fungsi implisit}
\label{thm:implicit}
Misalkan $U \subset \R^{n+m}$ himpunan terbuka dan $f \colon U \to \R^m$
pemetaan berkelas $C^1(U)$.  Misalkan $(p,q) \in U$ titik sedemikian sehingga
$f(p,q) = 0$ dan
\begin{equation*}
\frac{\partial(f_1,\ldots,f_m)}{\partial(y_1,\ldots,y_m)} (p,q)  \neq 0 .
\end{equation*}
Maka terdapat
himpunan terbuka $W \subset \R^n$ dengan $p \in W$,
himpunan terbuka $W' \subset \R^m$ dengan $q \in W'$,
sedemikian sehingga $W \times W' \subset U$,
serta
pemetaan $g \colon W \to W'$ berkelas $C^1(W)$, dengan $g(p) = q$, dan
untuk setiap $x \in W$, titik $g(x)$ merupakan satu-satunya titik di $W'$
sedemikian sehingga 
\begin{equation*}
f\bigl(x,g(x)\bigr) = 0 .
\end{equation*}
Selain itu, jika $A = [ A_x ~ A_y ] = f'(p,q)$, maka
\begin{equation*}
g'(p) = -{(A_y)}^{-1}A_x .
\end{equation*}
\end{thm}

Syarat
$\frac{\partial(f_1,\ldots,f_m)}{\partial(y_1,\ldots,y_m)} (p,q) =
\det(A_y)  \neq 0$
semata-mata berarti bahwa $A_y$ invertibel.  Jika $n=m=1$, syarat tersebut 
adalah $\frac{\partial f}{\partial y}(p,q) \neq 0$, dan $W$ serta $W'$ merupakan 
interval terbuka.  Lihat \figureref{fig:implicitfunc}.
\begin{myfigureht}
\myincludepdft{implicitfunc}{%
Diagram sebuah lingkaran pada bidang xy yang berlabel f dari x koma y sama dengan 0.
Sebuah titik berlabel p koma q ditandai pada lingkaran di kuadran pertama.
Titik itu berada di dalam persegi panjang kecil berarsir yang berlabel W kali W aksen.
Proyeksi persegi panjang berarsir pada sumbu x dan y masing-masing
berlabel W dan W aksen; keduanya terletak di antara titik asal
dan titik potong lingkaran dengan sumbu yang bersesuaian.}
\caption{Teorema fungsi implisit untuk $f(x,y) = x^2+y^2-1$ pada $U=\R^2$ dan
$(p,q)$ di kuadran pertama.\label{fig:implicitfunc}}
\end{myfigureht}

\begin{proof}
Definisikan $F \colon U \to \R^{n+m}$ dengan $F(x,y) \coloneqq \bigl(x,f(x,y)\bigr)$.
Jelas bahwa $F$ berkelas $C^1$, dan kita ingin menunjukkan bahwa turunannya
di $(p,q)$ invertibel.
Mari kita hitung turunannya.  Hasil bagi
\begin{equation*}
\frac{\bnorm{f(p+h,q+k) - f(p,q) - A_x h - A_y k}}{\bnorm{(h,k)}}
\end{equation*}
menuju nol ketika $\bnorm{(h,k)} = \sqrt{\snorm{h}^2+\snorm{k}^2}$ menuju nol.
Karena itu, hasil bagi berikut juga menuju nol:
\begin{multline*}
\frac{\bnorm{F(p+h,q+k)-F(p,q) - (h,A_x h+A_y k)}}{\bnorm{(h,k)}}
\\
\begin{aligned}
& =
\frac{\bnorm{\bigl(h,f(p+h,q+k)-f(p,q)\bigr) - (h,A_x h+A_y
k)}}{\bnorm{(h,k)}}
\\
& =
\frac{\bnorm{f(p+h,q+k) - f(p,q) - A_x h - A_y k}}{\bnorm{(h,k)}} .
\end{aligned}
\end{multline*}
Jadi, turunan $F$ di $(p,q)$ memetakan $(h,k)$ ke $(h,A_x h+A_y k)$.
Dalam bentuk matriks blok, turunan tersebut adalah
$\left[\begin{smallmatrix}I & 0\\A_x & A_y\end{smallmatrix}\right]$.  Jika 
$(h,A_x h+A_y k) = (0,0)$, maka $h=0$, sehingga $A_y k = 0$.  Karena $A_y$
injektif, $k=0$.  Dengan demikian, $F'(p,q)$ injektif, dan karena itu invertibel.
Kita terapkan teorema fungsi invers.

Dengan kata lain, terdapat himpunan terbuka $V \subset \R^{n+m}$ dengan
$F(p,q) = (p,0) \in V$,
serta pemetaan
$G \colon V \to \R^{n+m}$ berkelas $C^1$, sedemikian sehingga $F\bigl(G(x,s)\bigr) = (x,s)$ untuk
setiap $(x,s) \in V$, $G$ injektif, dan $G(V)$ merupakan himpunan terbuka yang termuat dalam $U$; pada
$G(V)$, pemetaan $F$ dan $G$ saling invers. % (dengan $x \in \R^n$ dan $s \in \R^m$).
Tuliskan $G = (G_1,G_2)$ (berturut-turut $n$ komponen pertama dan $m$ komponen berikutnya dari $G$).
Maka
\begin{equation*}
F\bigl(G_1(x,s),G_2(x,s)\bigr) = \Bigl(G_1(x,s),f\bigl(G_1(x,s),G_2(x,s) \bigr)\Bigr)
= (x,s) .
\end{equation*}
Jadi, $x = G_1(x,s)$ dan $f\bigl(G_1(x,s),G_2(x,s)\bigr) = f\bigl(x,G_2(x,s)\bigr) = s$.
Untuk setiap $x$ sedemikian sehingga $(x,0) \in V$, dengan menyubstitusikan $s=0$ kita memperoleh
\begin{equation*}
f\bigl(x,G_2(x,0)\bigr) = 0 .
\end{equation*}
Karena himpunan $G(V)$ terbuka dan $(p,q) \in G(V)$,
terdapat himpunan-himpunan terbuka
$\widetilde{W}$ dan $W'$ sedemikian sehingga $\widetilde{W} \times W' \subset G(V)$, dengan $p
\in \widetilde{W}$ dan
$q \in W'$.  Karena $V$ terbuka dan $(p,0) \in V$, setelah memperkecil
$\widetilde{W}$ jika perlu, kita juga dapat memastikan bahwa $(x,0) \in V$ untuk setiap
$x \in \widetilde{W}$.
Ambil $W \coloneqq \bigl\{ x \in \widetilde{W} : G_2(x,0) \in W' \bigr\}$.
Fungsi yang memetakan $x$ ke $G_2(x,0)$ berkelas $C^1$ pada $\widetilde{W}$; karena
$W'$ terbuka, $W$ terbuka.  Selain itu, $G(p,0)=(p,q)$; jadi $p \in W$.
Karena $W \subset \widetilde{W}$, kita juga mempunyai $W \times W' \subset G(V) \subset U$.
Definisikan
$g \colon W \to W'$ dengan $g(x) \coloneqq G_2(x,0)$.  Fungsi ini berkelas $C^1$,
$g(p)=q$, dan persamaan di atas memberikan $f\bigl(x,g(x)\bigr)=0$ untuk setiap $x \in W$.
Jika $y \in W'$ juga memenuhi $f(x,y)=0$, maka
$F(x,y)=(x,0)=F\bigl(x,g(x)\bigr)$.  Kedua titik yang menjadi argumen $F$ tersebut berada di $G(V)$,
dan $F$ injektif pada $G(V)$; oleh karena itu, $y=g(x)$.

Selanjutnya, diferensiasikan pemetaan
\begin{equation*}
x\mapsto f\bigl(x,g(x)\bigr)
\end{equation*}
di $p$.
Pemetaan tersebut adalah pemetaan nol, sehingga turunannya nol.
Dengan menggunakan aturan rantai,
\begin{equation*}
0 = A\bigl(h,g'(p)h\bigr) = A_xh + A_yg'(p)h
\end{equation*}
untuk setiap $h \in \R^{n}$,
dan kita memperoleh turunan $g$ yang diinginkan.
\end{proof}

Dengan kata lain, dalam konteks teorema tersebut, kita mempunyai
$m$ persamaan dengan $n+m$ peubah tak diketahui:
\begin{align*}
& f_1 (x_1,\ldots,x_n,y_1,\ldots,y_m) = 0 , \\
& f_2 (x_1,\ldots,x_n,y_1,\ldots,y_m) = 0 , \\
& \qquad \qquad \qquad  \vdots \\
& f_m (x_1,\ldots,x_n,y_1,\ldots,y_m) = 0 .
\end{align*}
Teorema tersebut menjamin adanya solusi lokal di sekitar $(p,q)$ jika $f(p,q)=0$,
$f=(f_1,f_2,\ldots,f_m)$ merupakan pemetaan berkelas $C^1$
(komponen-komponennya berkelas $C^1$: turunan parsial terhadap semua variabel ada
dan kontinu) dan matriks
\begin{equation*}
\begin{bmatrix}
\frac{\partial f_1}{\partial y_1}
&
\frac{\partial f_1}{\partial y_2}
& \ldots &
\frac{\partial f_1}{\partial y_m}
\\[6pt]
\frac{\partial f_2}{\partial y_1}
&
\frac{\partial f_2}{\partial y_2}
& \ldots &
\frac{\partial f_2}{\partial y_m}
\\
\vdots & \vdots & \ddots & \vdots
\\
\frac{\partial f_m}{\partial y_1}
&
\frac{\partial f_m}{\partial y_2}
& \ldots &
\frac{\partial f_m}{\partial y_m}
\end{bmatrix}
\end{equation*}
invertibel di $(p,q)$.

\begin{example}
Tinjau himpunan yang diberikan oleh $x^2+y^2-{(z+1)}^3 = -1$ dan $e^x+e^y+e^z = 3$
di sekitar titik $(0,0,0)$.
Himpunan ini merupakan himpunan nol dari pemetaan
\begin{equation*}
f(x,y,z) = \bigl(x^2+y^2-{(z+1)}^3+1,e^x+e^y+e^z-3\bigr) ,
\end{equation*}
yang turunannya adalah
\begin{equation*}
f' =
\begin{bmatrix}
2x & 2y & -3{(z+1)}^2 \\
e^x & e^y & e^z
\end{bmatrix} .
\end{equation*}
Matriks
\begin{equation*}
\begin{bmatrix}
2(0) & -3{(0+1)}^2 \\
e^0 & e^0
\end{bmatrix}
=
\begin{bmatrix}
0 & -3 \\
1 & 1
\end{bmatrix}
\end{equation*}
bersifat invertibel.  Oleh karena itu, di sekitar $(0,0,0)$ kita dapat menyatakan $y$ dan $z$
masing-masing sebagai fungsi berkelas $C^1$ dari $x$ sedemikian sehingga, untuk $x$ di sekitar $0$,
\begin{equation*}
x^2+y(x)^2-{\bigl(z(x)+1\bigr)}^3 = -1,
\qquad
e^x+e^{y(x)}+e^{z(x)} = 3 .
\end{equation*}
Dengan kata lain, di sekitar titik asal, himpunan solusi merupakan
kurva mulus di $\R^3$ yang melalui titik asal.
Teorema tersebut tidak memberi tahu kita cara menemukan $y(x)$ dan $z(x)$ secara eksplisit;
teorema itu hanya menjamin bahwa kedua fungsi tersebut ada.
\end{example}

Pengamatan yang menarik dan terkadang berguna dari pembuktian tersebut adalah bahwa kita sebenarnya menyelesaikan persamaan
$f\bigl(x,G_2(x,s)\bigr) = s$ untuk setiap $(x,s)$ dalam suatu lingkungan dari $(p,0)$, bukan hanya
untuk $s=0$.

\begin{remark}
Terdapat versi teorema tersebut untuk tingkat keterdiferensialan sebarang:
Jika $f$ memiliki turunan kontinu hingga orde $k$ (lihat bagian berikutnya), maka solusinya juga memiliki
turunan kontinu hingga orde $k$.
\end{remark}


\subsection{Latihan}

\begin{exercise}
Misalkan $C \coloneqq \bigl\{ (x,y) \in \R^2 : x^2+y^2 = 1 \bigr\}$.
\begin{enumerate}[a)]
\item
Nyatakan $y$ sebagai fungsi dari $x$ di sekitar $(0,1)$ (yaitu, carilah fungsi $g$
dari teorema fungsi implisit untuk suatu lingkungan dari titik $(p,q) = (0,1)$).
\item
Nyatakan $y$ sebagai fungsi dari $x$ di sekitar $(0,-1)$.
\item
Nyatakan $x$ sebagai fungsi dari $y$ di sekitar $(-1,0)$.
\end{enumerate}
\end{exercise}

\begin{exercise}
Definisikan $f \colon \R^2 \to \R^2$ dengan $f(x,y) \coloneqq
\bigl(x,y+h(x)\bigr)$ untuk suatu fungsi $h \colon \R \to \R$ yang diferensiabel
secara kontinu.
\begin{enumerate}[a)]
\item
Tunjukkan bahwa $f$ injektif dan surjektif.
\item
Hitung $f'$.  (Pastikan Anda menjelaskan mengapa $f'$ ada.)
\item
Tunjukkan bahwa $f'$ invertibel di setiap titik dan hitung
${\bigl(f'(x,y)\bigr)}^{-1}$.
\end{enumerate}
\end{exercise}

\begin{exercise}
Definisikan $f \colon \R^2 \to \R^2 \setminus \bigl\{ (0,0) \bigr\}$ dengan
$f(x,y) \coloneqq \bigl(e^x\cos(y),e^x\sin(y)\bigr)$.
\begin{enumerate}[a)]
\item
Tunjukkan bahwa $f$ surjektif.
\item
Tunjukkan bahwa $f'$ invertibel di setiap titik.
\item
Tunjukkan bahwa $f$ tidak injektif; bahkan, untuk setiap $(a,b) \in \R^2
\setminus \bigl\{ (0,0) \bigr\}$,
terdapat tak berhingga banyak titik berbeda $(x,y) \in \R^2$ sedemikian sehingga 
$f(x,y) = (a,b)$.
\end{enumerate}
Oleh karena itu, turunan yang invertibel di setiap titik tidak berarti bahwa
$f$ invertibel secara global.\\
Catatan: Silakan gunakan pengetahuan Anda tentang sinus dan kosinus dari kalkulus.
\end{exercise}

\begin{exercise}
Carilah pemetaan $f \colon \R^n \to \R^n$ yang injektif, surjektif,
dan diferensiabel secara kontinu, tetapi $f'(0) = 0$.  Petunjuk: Generalisasikan $f(x) = x^3$ dari satu
dimensi ke $n$ dimensi.
\end{exercise}

\begin{exercise}
Tinjau $z^2 + xz + y =0$ di $\R^3$.  Carilah ekspresi $D(x,y)$ sedemikian sehingga,
jika $D(x_0,y_0) \neq 0$ dan $z^2+x_0z+y_0 = 0$ untuk suatu $z \in \R$,
maka pada suatu lingkungan dari $(x_0,y_0)$ terdapat
tepat dua fungsi berbeda $r_1(x,y)$
dan $r_2(x,y)$ yang diferensiabel secara kontinu, sedemikian sehingga kedua nilai
$z=r_1(x,y)$ dan $z=r_2(x,y)$ merupakan solusi dari $z^2 + xz + y =0$.
Apakah Anda mengenali bentuk $D$ dari aljabar?
\end{exercise}


\begin{exercise}
Andaikan $f \colon (a,b) \to \R^2$ diferensiabel secara kontinu dan
komponen pertama (komponen $x$) dari $f'(t)$ tidak sama dengan 0
untuk setiap $t \in (a,b)$.
Buktikan bahwa terdapat interval terbuka $I \subset \R$ dan
fungsi $g \colon I \to \R$ yang diferensiabel secara kontinu
sedemikian sehingga 
$(x,y) \in f\bigl((a,b)\bigr)$ jika dan hanya jika $x \in I$ dan $y=g(x)$.
Dengan kata lain, himpunan
$f\bigl((a,b)\bigr)$ merupakan grafik dari $g$.
\end{exercise}

\begin{samepage}
\begin{exercise}
Definisikan $f \colon \R^2 \to \R^2$ dengan
\begin{equation*}
f(x,y) \coloneqq
\begin{cases}
\bigl(x^2 \sin (\nicefrac{1}{x}) + \nicefrac{x}{2} , y \bigr) &
 \text{jika } x \neq 0, \\
(0,y) &
 \text{jika } x=0.
\end{cases}
\end{equation*}
\begin{enumerate}[a)]
\item
Tunjukkan bahwa $f$ diferensiabel di setiap titik.
\item
Tunjukkan bahwa $f'(0,0)$ invertibel.
\item
Tunjukkan bahwa $f$ tidak injektif pada lingkungan mana pun dari titik asal (fungsi ini
tidak invertibel secara lokal; dengan kata lain, kesimpulan teorema fungsi invers tidak berlaku).
\item
Tunjukkan bahwa $f$ tidak diferensiabel secara kontinu.
\end{enumerate}
Catatan: Silakan gunakan pengetahuan Anda tentang sinus dan kosinus dari kalkulus.
\end{exercise}
\end{samepage}

\begin{exercise}[Koordinat polar] \label{mv:exercise:polarcoordinates}
\index{koordinat polar}
Definisikan pemetaan $F(r,\theta) \coloneqq \bigl(r \cos(\theta), r \sin(\theta) \bigr)$.
\begin{enumerate}[a)]
\item
Tunjukkan bahwa $F$ diferensiabel secara kontinu (untuk setiap $(r,\theta) \in
\R^2$).
\item
Hitung $F'(0,\theta)$ untuk setiap $\theta$.
\item
Tunjukkan bahwa jika $r \neq 0$, maka $F'(r,\theta)$ invertibel; oleh karena itu,
$F$ memiliki invers lokal di sekitar setiap titik dengan $r \neq 0$.
\item
Tunjukkan bahwa $F \colon \R^2 \to \R^2$ surjektif dan, untuk setiap titik $(x,y) \in
\R^2$, himpunan $F^{-1}(x,y)$ tak berhingga.
\item
Tunjukkan bahwa $F \colon \R^2 \to \R^2$ bukan pemetaan terbuka.
Perhatikan bahwa $F|_{(0,\infty) \times \R}$ merupakan pemetaan terbuka
menurut \corref{cor:IVTopenmap}.
Petunjuk: Ke manakah
persegi panjang terbuka kecil seperti $(-\epsilon,\epsilon) \times
(-\epsilon,\epsilon)$ dipetakan?
\item
Tunjukkan bahwa $F|_{(0,\infty) \times [0,2\pi)}$ merupakan bijeksi ke
$\R^2 \setminus \bigl\{ (0,0) \bigr\}$.
\end{enumerate}
Catatan: Silakan gunakan pengetahuan Anda tentang sinus dan kosinus dari kalkulus.
\end{exercise}

\begin{exercise}
Misalkan $H \coloneqq \bigl\{ (x,y) \in \R^2 : y > 0 \bigr\}$, dan untuk $(x,y) \in H$
definisikan
\begin{equation*}
F(x,y) \coloneqq \left(
\frac{x^2+y^2-1}{x^2+2y+y^2+1}
,~
\frac{-2x}{x^2+2y+y^2+1}
\right) .
\end{equation*}
Buktikan bahwa $F$ merupakan pemetaan bijektif dari $H$ ke $B(0,1)$, bahwa $F$
diferensiabel secara kontinu pada $H$, dan bahwa inversnya juga diferensiabel
secara kontinu.
\end{exercise}

\begin{exercise}
Andaikan $U \subset \R^2$ terbuka dan $f \colon U \to \R$ adalah
fungsi berkelas $C^1$ sedemikian
sehingga $\nabla f(x,y) \neq 0$ untuk setiap $(x,y) \in U$.  Tunjukkan bahwa setiap
himpunan level tak kosong merupakan kurva reguler berkelas $C^1$ secara lokal.  Dengan kata lain,
untuk setiap $(x,y) \in U$, terdapat $\delta > 0$, lingkungan terbuka $V \subset U$ dari $(x,y)$,
dan fungsi $\gamma \colon (-\delta,\delta) \to V$ berkelas $C^1$ dengan
$\gamma(0)=(x,y)$ serta $\gamma'(t) \neq 0$ untuk setiap $t \in (-\delta,\delta)$, sedemikian sehingga
$\gamma\bigl((-\delta,\delta)\bigr)
= \bigl\{(u,v) \in V : f(u,v)=f(x,y)\bigr\}$.
\end{exercise}

\begin{exercise}
Andaikan $U \subset \R^2$ terbuka dan $f \colon U \to \R$ adalah
fungsi berkelas $C^1$ sedemikian
sehingga $\nabla f(x,y) \neq 0$ untuk setiap $(x,y) \in U$.
Tunjukkan bahwa untuk setiap $(x,y) \in U$ terdapat lingkungan terbuka $V \subset U$ dari $(x,y)$,
himpunan terbuka $W \subset \R^2$, serta fungsi bijektif berkelas $C^1$
$g \colon W \to V$ yang inversnya berkelas $C^1$, sedemikian sehingga
himpunan-himpunan level dari $f \circ g$ merupakan garis-garis horizontal di $W$; yaitu, untuk setiap
$c \in (f \circ g)(W)$ terdapat $t_0=t_0(c)$ sedemikian sehingga
$\bigl\{(s,t) \in W : (f \circ g)(s,t)=c\bigr\}
=\bigl\{ (s,t_0) \in \R^2 : s \in \R, (s,t_0) \in W \bigr\}$.
Dengan kata lain, kurva-kurva level dapat \myquote{diluruskan} secara lokal.
\end{exercise}

%%%%%%%%%%%%%%%%%%%%%%%%%%%%%%%%%%%%%%%%%%%%%%%%%%%%%%%%%%%%%%%%%%%%%%%%%%%%%%

\sectionnewpage
\section{Turunan berorde lebih tinggi}
\label{sec:mvhighordders}

\sectionnotes{kurang dari satu sesi kuliah, opsional, lihat juga bagian opsional
\volIref{\sectionref*{vI-sec:taylor} dalam jilid I}{\sectionref{sec:taylor}}}

Misalkan $U \subset \R^n$ himpunan terbuka dan $f \colon U \to \R$ suatu fungsi.
Tuliskan koordinat sebagai $x = (x_1,x_2,\ldots,x_n) \in \R^n$.
Jika $\frac{\partial f}{\partial x_{\ell}}$ ada di setiap titik dalam $U$,
maka turunan tersebut juga merupakan fungsi $\frac{\partial f}{\partial x_{\ell}}
\colon U \to \R$.  Oleh karena itu, kita dapat membicarakan turunan
parsialnya.  Kita menyatakan 
turunan parsial $\frac{\partial f}{\partial x_{\ell}}$ terhadap
$x_m$ dengan\glsadd{not:multivarpartder}
\begin{equation*}
\frac{\partial^2 f}{\partial x_m \partial x_{\ell}}
\coloneqq
\frac{\partial \Bigl( \frac{\partial f}{\partial x_{\ell}} \Bigr)}{\partial x_m} .
\end{equation*}
Jika $m={\ell}$, agar lebih ringkas kita menuliskan 
$\frac{\partial^2 f}{\partial x_{\ell}^2}$.

Kita mendefinisikan turunan berorde lebih tinggi secara induktif.
Untuk suatu bilangan bulat $k \geq 2$, andaikan ${\ell}_1,{\ell}_2,\ldots,{\ell}_k$ adalah bilangan bulat antara $1$ dan $n$, dan
andaikan 
\begin{equation*}
\frac{\partial^{k-1} f}{\partial x_{{\ell}_{k-1}} \partial x_{{\ell}_{k-2}}
\cdots \partial x_{{\ell}_1}}
\end{equation*}
ada dan dapat diturunkan terhadap peubah $x_{{\ell}_{k}}$.
Maka turunan parsial terhadap peubah tersebut dinyatakan dengan
\begin{equation*}
\frac{\partial^{k} f}{\partial x_{{\ell}_{k}} \partial x_{{\ell}_{k-1}}
\cdots \partial x_{{\ell}_1}}
\coloneqq 
\frac{\partial \Bigl( \frac{\partial^{k-1} f}{\partial x_{{\ell}_{k-1}} \partial
x_{{\ell}_{k-2}} \cdots \partial x_{{\ell}_1}} \Bigr)}{\partial x_{{\ell}_{k}}} .
\end{equation*}
Turunan semacam itu disebut
\emph{\myindex{turunan parsial orde $k$}}.

Terkadang notasi $f_{x_{\ell} x_m}$\glsadd{not:subpartder} digunakan untuk
$\frac{\partial^2 f}{\partial x_m \partial x_{\ell}}$.  Dibandingkan dengan notasi pecahan,
notasi subskrip ini menuliskan urutan variabel diferensiasi secara terbalik, dan hal itu dapat menjadi penting.

\begin{defn}
Untuk suatu bilangan bulat positif $k$, andaikan $U \subset \R^n$ himpunan terbuka dan
$f \colon U \to \R$ suatu fungsi.  Kita mengatakan bahwa $f$ adalah
\emph{fungsi yang diferensiabel secara kontinu hingga orde $k$}%
\index{fungsi yang diferensiabel secara kontinu hingga orde k@diferensiabel secara kontinu hingga orde $k$}\index{diferensiabel secara kontinu},
atau fungsi berkelas $C^k$\glsadd{not:Ck}, jika semua turunan parsial untuk setiap orde
sampai dengan orde $k$ ada dan kontinu.
\end{defn}

Jadi, fungsi yang diferensiabel secara kontinu, atau berkelas $C^1$, adalah fungsi yang memiliki semua
turunan parsial orde pertamanya
ada dan kontinu; berdasarkan \propref{mv:prop:contdiffpartials}, hal ini sesuai dengan
definisi kita sebelumnya.  Sebenarnya, kita
dapat mensyaratkan hanya bahwa semua turunan parsial tepat berorde $k$ ada dan
kontinu.
Keberadaan turunan parsial berorde lebih rendah jelas
diperlukan bahkan untuk mendefinisikan turunan parsial orde $k$.
Hal ini mengikuti secara rekursif: setiap turunan orde $j<k$ memiliki turunan parsial orde pertama
berupa turunan orde $j+1$, yang telah diketahui kontinu.

Jika turunan-turunan parsial kontinu, kita dapat menukar urutan diferensiasi parsial.

\begin{prop} \label{mv:prop:swapders}
Andaikan $U \subset \R^n$ terbuka dan $f \colon U \to \R$ fungsi berkelas $C^2$,
serta $\ell$ dan $m$ bilangan bulat antara $1$ dan $n$.  Maka
\begin{equation*}
\frac{\partial^2 f}{\partial x_m \partial x_{\ell}}
=
\frac{\partial^2 f}{\partial x_{\ell} \partial x_m} .
\end{equation*}
\end{prop}

\begin{proof}
Jika $\ell=m$, kesimpulannya langsung.  Jadi, andaikan $\ell \neq m$.
Tetapkan $p \in U$, dan misalkan $e_{\ell}$ serta $e_m$ adalah vektor basis standar.
Pilih dua bilangan positif $s$ dan $t$ yang cukup kecil sehingga
$p+s_0e_{\ell} +t_0e_m \in U$ apabila
$0 \leq s_0 \leq s$ dan $0 \leq t_0 \leq t$.
Setiap $s$ dan $t$ yang cukup kecil memenuhi syarat ini karena $U$ terbuka, sehingga
memuat bola terbuka kecil (atau kotak terbuka, jika Anda mau) di sekitar $p$.

Terapkan teorema nilai rata-rata pada fungsi
\begin{equation*}
\tau \mapsto f(p+se_{\ell} + \tau e_m)-f(p + \tau e_m) ,
\end{equation*}
pada interval $[0,t]$
untuk memperoleh $t_0 \in (0,t)$
sedemikian sehingga
\begin{equation*}
%mbxlatex \begin{aligned}
%mbxlatex &
\frac{f(p+se_{\ell} + te_m)- f(p+t e_m) - f(p+s e_{\ell})+f(p)}{t}
%mbxlatex \\
%mbxlatex & \qquad \qquad \qquad \qquad
=
\frac{\partial f}{\partial x_m}(p + s e_{\ell} + t_0 e_m)
-
\frac{\partial f}{\partial x_m}(p + t_0 e_m) .
%mbxlatex \end{aligned}
\end{equation*}
Dengan cara serupa, terdapat bilangan $s_0 \in (0,s)$ sedemikian sehingga
\begin{equation*}
\frac{\frac{\partial f}{\partial x_m}(p + s e_{\ell} + t_0 e_m)
-
\frac{\partial f}{\partial x_m}(p + t_0 e_m)}{s}
=
\frac{\partial^2 f}{\partial x_{\ell} \partial x_m}(p + s_0 e_{\ell} + t_0 e_m) .
\end{equation*}
Dengan kata lain,
\begin{equation*}
%mbxlatex \begin{aligned}
g(s,t)
%mbxlatex &
\coloneqq
\frac{f(p+se_{\ell} + te_m)- f(p+t e_m) - f(p+s e_{\ell})+f(p)}{st}
%mbxlatex \\
%mbxlatex &
=
\frac{\partial^2 f}{\partial x_{\ell} \partial x_m}(p + s_0 e_{\ell} + t_0 e_m) .
%mbxlatex \end{aligned}
\end{equation*}

\begin{myfigureht}
\myincludepdft{der2orderflip}{%
Diagram pada sebuah bidang
dengan titik p yang ditandai.
Panah horizontal ke kanan berlabel e subskrip el,
dan panah vertikal ke atas berlabel e subskrip m.
Titik p ditambah s kali e subskrip el ditandai
pada panah horizontal.
Pada panah vertikal terdapat titik
p ditambah t kali e subskrip m dan, di bawahnya,
titik 
p ditambah t subskrip nol kali e subskrip m.
Di atas p ditambah s kali e subskrip el dan di sebelah kanan p ditambah t kali e subskrip m
terdapat titik p ditambah s kali e subskrip el ditambah t kali e subskrip m.
Di atas p ditambah s kali e subskrip el dan di sebelah kanan p ditambah t subskrip nol kali e subskrip m
terdapat titik p ditambah s kali e subskrip el ditambah t subskrip nol kali e subskrip m.
Terakhir, di antara titik
p ditambah t subskrip nol kali e subskrip m dan
p ditambah s kali e subskrip el ditambah t subskrip nol kali e subskrip m
terdapat titik
p ditambah s subskrip nol kali e subskrip el ditambah t subskrip nol kali e subskrip m.}
\caption{Menggunakan teorema nilai rata-rata untuk menaksir
turunan parsial orde kedua dengan
suatu hasil bagi selisih.\label{fig:der2orderflip}}
\end{myfigureht}

Lihat \figureref{fig:der2orderflip}.
Nilai $s_0$ dan $t_0$ bergantung pada $s$ dan $t$,
tetapi $0 < s_0 < s$ dan
$0 < t_0 < t$.
Ambil himpunan $(0,\epsilon) \times
(0,\epsilon)$ sebagai domain fungsi $g$, untuk suatu $\epsilon > 0$ yang kecil.
Ketika $(s,t) \in (0,\epsilon) \times (0,\epsilon)$ menuju $(0,0)$,
titik
$(s_0,t_0)$ juga menuju $(0,0)$.
Berdasarkan kekontinuan turunan parsial orde kedua,
\begin{equation*}
\lim_{(s,t) \to (0,0)} g(s,t) = 
\frac{\partial^2 f}{\partial x_{\ell} \partial x_m}(p) .
\end{equation*}

Sekarang pertukarkan peran $s$ dan $t$ (serta $\ell$ dan $m$).
Terapkan teorema nilai rata-rata pada
fungsi $\sigma \mapsto f(p+\sigma e_{\ell} + te_m)-f(p + \sigma
e_{\ell})$ untuk
memperoleh $s_1 \in (0,s)$ sedemikian sehingga
\begin{equation*}
%mbxlatex \begin{aligned}
%mbxlatex &
\frac{f(p+ se_{\ell} + te_m )- f(p+s e_{\ell}) - f(p+t e_m)+f(p)}{s}
%mbxlatex \\
%mbxlatex & \qquad \qquad \qquad \qquad
=
\frac{\partial f}{\partial x_{\ell}}(p + s_1 e_{\ell} + t e_m)
-
\frac{\partial f}{\partial x_{\ell}}(p + s_1 e_{\ell}) .
%mbxlatex \end{aligned}
\end{equation*}
Dengan menerapkan teorema nilai rata-rata sekali lagi, terdapat $t_1 \in (0,t)$ sedemikian sehingga
\begin{equation*}
\frac{\frac{\partial f}{\partial x_{\ell}}(p + s_1 e_{\ell} + t e_m)
-
\frac{\partial f}{\partial x_{\ell}}(p + s_1 e_{\ell})}{t}
=
\frac{\partial^2 f}{\partial x_m \partial x_{\ell}}(p + s_1 e_{\ell} + t_1 e_m) .
\end{equation*}
Jadi, $g(s,t) = \frac{\partial^2 f}{\partial x_m \partial
x_{\ell}}(p + s_1 e_{\ell} + t_1 e_m)$ untuk fungsi $g$ yang sama seperti di atas.
Seperti sebelumnya,
\begin{equation*}
\lim_{(s,t) \to (0,0)} g(s,t) = 
\frac{\partial^2 f}{\partial x_m \partial x_{\ell}}(p) .
\end{equation*}
Oleh karena itu, kedua turunan parsial tersebut sama.
\end{proof}

Kesimpulan proposisi tersebut tidak selalu berlaku tanpa asumsi bahwa kedua turunan parsial campuran
yang bersangkutan kontinu.  Lihat \exerciseref{exercise:mixedpartialsunequal}.
Perhatikan juga bahwa kita sebenarnya tidak memerlukan fungsi berkelas $C^2$---kita hanya memerlukan
kedua turunan parsial orde kedua yang terlibat bersifat kontinu.

\subsection{Latihan}

\begin{exercise}
Andaikan $f \colon U \to \R$ fungsi berkelas $C^2$ untuk suatu himpunan terbuka $U \subset
\R^n$ dan $p \in U$, serta misalkan $\ell,m \in \{1,\ldots,n\}$.
Gunakan pembuktian \propref{mv:prop:swapders} untuk mencari suatu ekspresi
yang hanya memuat nilai-nilai $f$ (padanan hasil bagi selisih
untuk turunan pertama), yang limitnya adalah
$\frac{\partial^2 f}{ \partial x_{\ell} \partial x_m}(p)$.
\end{exercise}

\begin{exercise} \label{exercise:mixedpartialsunequal}
Definisikan
\begin{equation*}
f(x,y) \coloneqq
\begin{cases}
\frac{xy(x^2-y^2)}{x^2+y^2} & \text{jika } (x,y) \neq (0,0),\\
0                           & \text{jika } (x,y) = (0,0).
\end{cases}
\end{equation*}
Tunjukkan bahwa
\begin{enumerate}[a)]
\item
Semua turunan parsial orde pertama ada dan kontinu.
\item
Turunan parsial
$\frac{\partial^2 f}{\partial x \partial y}$ dan
$\frac{\partial^2 f}{\partial y \partial x}$ ada, tetapi tidak kontinu
di $(0,0)$, dan
$\frac{\partial^2 f}{\partial x \partial y}(0,0) \neq 
\frac{\partial^2 f}{\partial y \partial x}(0,0)$.
\end{enumerate}
\end{exercise}

\begin{exercise}
Untuk suatu bilangan bulat positif $k$, misalkan $f \colon U \to \R$ fungsi berkelas $C^k$ untuk suatu himpunan terbuka $U \subset \R^n$
dan $p \in U$.
Andaikan ${\ell}_1,{\ell}_2,\ldots,{\ell}_k$ bilangan bulat antara $1$
dan $n$, serta $\sigma=(\sigma_1,\sigma_2,\ldots,\sigma_k)$ suatu
permutasi dari $(1,2,\ldots,k)$.  Buktikan bahwa
\begin{equation*}
\frac{\partial^{k} f}{\partial x_{{\ell}_{k}} \partial x_{{\ell}_{k-1}}
\cdots \partial x_{{\ell}_1}} (p)
=
\frac{\partial^{k} f}{\partial x_{{\ell}_{\sigma_k}} \partial
x_{{\ell}_{\sigma_{k-1}}}
\cdots \partial x_{{\ell}_{\sigma_1}}} (p) .
\end{equation*}
\end{exercise}

\begin{exercise}
Untuk suatu bilangan bulat positif $k$, andaikan $\varphi \colon \R^2 \to \R$ fungsi berkelas $C^k$
sedemikian sehingga
$\varphi(0,\theta) = \varphi(0,\psi)$ untuk setiap $\theta,\psi \in \R$,
serta
$\varphi(r,\theta) = \varphi(r,\theta+2\pi)$
dan
$\varphi(r,\theta) = \varphi(-r,\theta+\pi)$
untuk setiap $r,\theta \in \R$.
Misalkan $F(r,\theta) \coloneqq \bigl(r \cos(\theta), r \sin(\theta) \bigr)$ dari 
\exerciseref{mv:exercise:polarcoordinates}.  Tunjukkan bahwa fungsi
$g \colon \R^2 \to \R$ yang, untuk setiap $(x,y) \in \R^2$, didefinisikan dengan memilih sembarang
$(r,\theta) \in F^{-1}\bigl(\{(x,y)\}\bigr)$ dan menetapkan
$g(x,y) \coloneqq \varphi(r,\theta)$ terdefinisi dengan baik (perhatikan bahwa
cabang invers lokal untuk $F$ hanya tersedia di sekitar titik-titik tak nol), dan bahwa
pembatasannya pada $\R^2 \setminus \bigl\{ (0,0) \bigr\}$ merupakan fungsi berkelas $C^k$.
\\
Catatan: Anda boleh menggunakan pengetahuan Anda tentang sinus dan kosinus dari kalkulus.
\end{exercise}

\begin{exercise}
Andaikan $f \colon \R^2 \to \R$ fungsi berkelas $C^2$.
Untuk setiap $(x,y) \in \R^2$, hitung
\begin{equation*}
\lim_{t \to 0}
\frac{f(x+t,y)+f(x-t,y)+f(x,y+t)+f(x,y-t) - 4f(x,y)}{t^2}
\end{equation*}
Nyatakan limit tersebut dalam bentuk turunan-turunan parsial $f$.
\end{exercise}

\begin{exercise}
Andaikan $f \colon \R^2 \to \R$ fungsi sedemikian sehingga
semua turunan parsial orde pertama dan kedua ada.  Selain itu,
andaikan semua turunan parsial orde kedua merupakan fungsi terbatas.
Buktikan bahwa $f$ diferensiabel secara kontinu.
\end{exercise}

\begin{samepage}
\begin{exercise}
Ikuti strategi di bawah ini untuk
membuktikan versi sederhana berikut dari uji turunan kedua bagi
fungsi yang didefinisikan pada $\R^2$ (dengan menggunakan $(x,y)$ sebagai koordinat):  \emph{Andaikan $f \colon \R^2
\to \R$ fungsi berkelas $C^2$ dengan
titik kritis di titik asal, $f'(0,0) = 0$.  Jika
\begin{equation*}
\frac{\partial^2 f}{\partial x^2} (0,0) > 0 \qquad \text{dan} \qquad
\frac{\partial^2 f}{\partial x^2} (0,0) \frac{\partial^2 f}{\partial y^2}
(0,0) -
{\left(\frac{\partial^2 f}{\partial x \partial y} (0,0) \right)}^2 > 0 ,
\end{equation*}
maka $f$ mencapai minimum lokal ketat pada $(0,0)$.}
Gunakan teknik berikut:  Pertama, tanpa mengurangi keumuman, andaikan
bahwa $f(0,0) = 0$.  Kemudian buktikan:
\begin{enumerate}[a)]
\item
Terdapat $A \in L(\R^2)$ yang invertibel sedemikian sehingga $g = f \circ A$
memenuhi $\frac{\partial^2 g}{\partial x \partial y} (0,0) = 0$, dan
$\frac{\partial^2 g}{\partial x^2} (0,0) =
\frac{\partial^2 g}{\partial y^2} (0,0) = 2$.
\item
Untuk setiap $\epsilon > 0$, terdapat $\delta > 0$
sedemikian sehingga 
$\babs{g(x,y) - (x^2 + y^2)} \leq \epsilon (x^2+y^2)$ untuk setiap
$(x,y) \in B\bigl((0,0),\delta\bigr)$.\\
Petunjuk: Anda dapat menggunakan teorema Taylor satu peubah.
\item
Ini berarti bahwa $g$, dan karena itu $f$, mencapai minimum lokal
ketat pada $(0,0)$.
\end{enumerate}
Catatan: Anda harus menghindari godaan untuk sekadar
menerapkan uji turunan kedua satu variabel sepanjang garis-garis yang melalui
titik asal; lihat \exerciseref{exercise:peano}.
\end{exercise}
\end{samepage}
